% Public paper source; file identities are recorded in MANIFEST.json.
\documentclass[11pt,reqno]{amsart}
\usepackage{lmodern}
\usepackage{amsmath,amssymb}
\usepackage{fvextra}
\IfFileExists{lltjp-fancyvrb.sty}{\usepackage{lltjp-fancyvrb}}{}
\usepackage{luatexja-fontspec}
\setmainfont{Latin Modern Roman}[SmallCapsFont=lmromancaps10-regular.otf]
\setsansfont{Latin Modern Sans}
\setmonofont{Latin Modern Mono}[Scale=0.9]
\setmainjfont{HaranoAjiMincho-Regular.otf}[BoldFont=HaranoAjiMincho-Bold.otf,ItalicFont=HaranoAjiMincho-Regular.otf,BoldItalicFont=HaranoAjiMincho-Bold.otf,SmallCapsFont=HaranoAjiMincho-Regular.otf]
\setsansjfont{HaranoAjiGothic-Medium.otf}[BoldFont=HaranoAjiGothic-Bold.otf,ItalicFont=HaranoAjiGothic-Medium.otf,SmallCapsFont=HaranoAjiGothic-Medium.otf]
\setmonojfont{HaranoAjiGothic-Medium.otf}
\renewcommand{\abstractname}{要旨}
\renewcommand{\proofname}{証明}
\renewcommand{\refname}{参考文献}
\renewcommand{\appendixname}{付録}
\renewcommand{\tablename}{表}
\renewcommand{\figurename}{図}
\makeatletter
\def\@setthanks{\def\thanks##1{\par##1}\thankses}
\makeatother
\theoremstyle{definition}
\usepackage[letterpaper,textwidth=6in,textheight=8.7in,centering]{geometry}
\usepackage{booktabs,array,longtable}
\usepackage{needspace}
\usepackage{xurl}
\usepackage[hidelinks,unicode]{hyperref}
\newtheorem{theorem}{定理}
\newtheorem{corollary}[theorem]{系}
\newtheorem{lemma}[theorem]{補題}
\newtheorem{proposition}[theorem]{命題}
\theoremstyle{definition}
\newtheorem*{definition}{定義}
\theoremstyle{remark}
\newtheorem*{remark}{注意}
\providecommand{\tightlist}{\setlength{\itemsep}{0pt}\setlength{\parskip}{0pt}}
\setlength{\emergencystretch}{2em}
\raggedbottom
\hypersetup{pdftitle={有理数を底とする床関数列の合成数項に対する有限グラフ証明書},
  pdfauthor={Yuri Odagiri}}
\pdfvariable suppressoptionalinfo 767
\listfiles
\title[有理数を底とする床関数列の有限グラフ証明書]{有理数を底とする床関数列の合成数項に対する有限グラフ証明書}
\author{Yuri Odagiri}
\address{Yuri Odagiri: 個人研究者}
\email{ixixizko@gmail.com}
\thanks{本稿は英語論文の全文日本語訳である。英語版を正本とする。}
\date{}
% BEGIN PUBLIC ADDITION figure-preamble
\newcommand{\DoNotLoadEpstopdf}{}
\def\FigureJapanese{1}
% Shared paper illustrations. Load in the preamble.
\usepackage{tikz}
\usetikzlibrary{arrows.meta,positioning,calc,patterns,shapes.geometric}
\providecommand{\figtext}[2]{\ifdefined\FigureJapanese#2\else#1\fi}
\definecolor{figblue}{RGB}{30,70,110}
\definecolor{figgray}{RGB}{105,105,105}
\tikzset{
  paperfig/.style={font=\fontsize{9}{11}\selectfont,>=Stealth,line width=.65pt},
  fbox/.style={draw,rounded corners=2pt,align=center,inner sep=5pt,fill=white},
  fnode/.style={draw,circle,minimum size=6mm,inner sep=1pt,fill=white},
  fedge/.style={->},
  femph/.style={draw=figblue,line width=1.05pt},
  flabel/.style={fill=white,inner sep=1.5pt,align=center},
  ftitle/.style={font=\fontsize{9.5}{12}\selectfont\bfseries,align=center}
}

% END PUBLIC ADDITION figure-preamble

% BEGIN PUBLIC ADDITION japanese-footer-spacing
\setlength{\footskip}{24pt}
% END PUBLIC ADDITION japanese-footer-spacing

% BEGIN PUBLIC ADDITION doi
\thanks{DOI: \url{https://doi.org/10.5281/zenodo.22887589}.}
\hypersetup{pdfsubject={DOI: 10.5281/zenodo.22887589}}
% END PUBLIC ADDITION doi

\begin{document}
% Markdown AST block 0: Finite graph certificates for composite terms in rational floor sequences
% Title appears in front matter.

% Markdown AST block 1: Abstract
\begin{abstract}

% Markdown AST block 2: Para
任意の実数 \(\xi>0\) に対し、\(\lfloor\xi(7/5)^n\rfloor\) が \(2,3,5,11,13\) の少なくとも一つで割り切れることが無限に多くの正整数 \(n\) について成り立ち、また、\(\lfloor\xi(5/2)^n\rfloor\) が \(2,3,7,11,13,17,19,23,29,31\) の少なくとも一つで割り切れることが無限に多くの正整数 \(n\) について成り立つことを、計算機援用により証明する。したがって、どちらの数列にも合成数の項が無限に存在する。証明では、指定された素数を避ける軌道を有限ラベル付きグラフのウォークとして表す。ここでウォーク（walk）とは、辺の向きに従って頂点を順にたどる列であり、同じ頂点や辺を何度通ってもよい。頂点は剰余と小数部のセルを記録し、辺には符号付きの繰り上がりを付す。出力が巡回位相のみに依存する成分は、初等的な非周期性補題によって除外される。語ラベル付きの構成では、辺に有限の繰り上がり語を付し、決定的な道を厳密に縮約するとともに、語の各中間時刻で新しい素数の条件を課す。比 \(7/5\) については、待ち時間の明示的な上界も証明する。すなわち、\(P=\lfloor\xi\rfloor>13\) ならば、添字が \(428+12L(P)\) 以下、\(L(P)=\min\{h\ge0:5^h\ge P+2\}\) のある項が、五つの素数のいずれかで割り切れる。さらに、この証明書による手法の適用範囲を明らかにする。縮約、素数冪、および \(2ab\) を割る素数を用いても、有限に認証可能な底のクラスは広がらない。\(a\ge2\operatorname{rad}M\) のときには、この手法は停止しない。また、補助素数の積 \(R\) は、各素数が \(2ab\) を割らないとき、\(a\le2R\,J(N_0)\) を満たす必要がある。ここで \(J\) はJacobsthal関数であり、初等的な評価から \(R\ge a^{1-o(1)}\) が従う。採用した有限計算はすべて、科学的中核を共有しない二つの実装で実行し、比較対象としたすべての有限データが要素ごとに一致することを確認した。Leanによる証明は、二つの整除性定理と、それらから従う合成数性の系を扱う。\(7/5\) の1ステップ証明はカーネルで検査され、語ラベル付きの証明は、カーネルで検査された健全性定理と \nolinkurl{native_decide} による有限評価を組み合わせる。後者では、コンパイラとネイティブ評価に対する追加の信頼を要する。

% Markdown AST block 3: 1. Introduction
\end{abstract}
\maketitle
\enlargethispage{7pt}

\Needspace{6\baselineskip}
\section{序論}\label{sec:1}

% Markdown AST block 4: Para
有理数の冪の整数部分に関する算術には、単純な非整数の底に対してさえ困難がある。数列が合成数の項を無限に含むことを示す有用な方法は、固定された有限個の素数を選び、無限に多くの項がそのいずれかを約数にもつことを示すことである。項がそれらの素数をすべて上回れば、整除性から合成数性が従う。Dubickas \hyperref[ref-dub06]{Dub06, §2} にならい、任意の \(\xi>0\) について \(\lfloor\xi b^n\rfloor\) がこの性質をもつような有限素数集合を、底 \(b\) に関して\emph{回避不能}（unavoidable）であるという。整数の底では、\(b\in\{2,3,4,6\}\) に対してこのような集合が存在する。一方、\(b=5\) の場合（\hyperref[ref-dn]{DN}。\hyperref[ref-dub06]{Dub06, §2} での紹介による）と、\(b-1\) が平方因子をもつ任意の \(b\ge3\) の場合（\hyperref[ref-dub06]{Dub06, Theorem 4}）には存在せず、それ以外の整数の底でも存在しないと予想されている（\hyperref[ref-dub06]{Dub06, Conjecture 3}）。この用語でいえば、次に述べる結果は有理数の底 \(3/2\)、\(4/3\)、\(5/4\) に対する回避不能集合を与え、定理1は \(7/5\) と \(5/2\) に対する回避不能集合を与える。

% Markdown AST block 5: Para
1967年にFormanとShapiro \hyperref[ref-fs]{FS} は、\(3/2\) と \(4/3\) の冪の整数部分が合成数を無限に含むことを証明した。歴史的経緯については \hyperref[ref-dn]{DN, pp.635--636} を参照されたい。同文献で指摘されているとおり、彼らの証明は、任意の固定された正の係数 \(\xi\) を掛けた場合にも適用できる。DubickasとNovikasは、さらに底 \(5/4\) を扱い、三つの底のそれぞれに対して明示的な整除性の結果を述べた。すなわち、任意の \(\xi>0\) に対し、
\[
\lfloor\xi(3/2)^n\rfloor,\qquad
\lfloor\xi(4/3)^n\rfloor,\qquad
\lfloor\xi(5/4)^n\rfloor
\]
の各数列の無限に多くの項が、それぞれ次の素数集合に属する約数をもつ。
\[
\{2,5,7,11\},\qquad \{2,3,5\},\qquad \{2,3,7,11,13\};
\]
\hyperref[ref-dn]{DN, Theorem 1, p.636} および \hyperref[ref-n]{N, Theorem 1.1, p.24} を参照されたい。

% Markdown AST block 6: Para
DubickasとNovikasは、素数 \(\{2,3,5,11\}\) に関する対応する主張を、\emph{最近接整数}
\(\lfloor\xi(7/5)^n+1/2\rfloor\)
に対しても証明した。これは \hyperref[ref-dn]{DN, Theorem 4, p.637} の第3の主張であり、\hyperref[ref-n]{N, Theorem 1.5(iii), p.26} にも述べられている。比 \(5/2\) については、任意の \(\xi>0\) に対し、\emph{シフトした}数列 \(\lfloor\xi(5/2)^n\rfloor-1\) に、\(2\)、\(3\) または \(5\) で割り切れる項が無限に存在することを証明した。\hyperref[ref-dn]{DN, Theorem 3, p.637} および \hyperref[ref-n]{N, Theorem 1.4, p.25} を参照されたい。その証明中では、\(\lfloor\xi(5/2)^n\rfloor-1+30k\) にも、任意の固定整数 \(k\) に対して同じ議論が適用できると指摘されている（\hyperref[ref-dn]{DN, p.646}、\hyperref[ref-n]{N, p.26}）。ここでは丸め方とシフトの取り方が重要である。本論文の結果は、シフトを加えない整数部分を扱う。\(5/2\) ではシフトによって許容される繰り上がりの文字集合が変わり、本論文の証明書では異なる素数集合を用いる。

% Markdown AST block 7: Theorem 1 (Main Theorem).
\begin{samepage}
\begin{theorem}[主定理]\leavevmode\label{thm:1}
任意の \(\xi\in\mathbb R\) で \(\xi>0\) を満たすものと、任意の \(N\in\mathbb N\) に対して、次が成り立つ。

% Markdown AST block 8: OrderedList
\begin{enumerate}
\def\labelenumi{(\roman{enumi})}
\item
  次を満たす整数 \(n\ge\max(N,1)\) が存在する。
  \[
  \gcd\!\left(\left\lfloor\xi(7/5)^n\right\rfloor,4290\right)>1,
  \qquad 4290=2\cdot3\cdot5\cdot11\cdot13;
  \tag{1}\label{eq:1}
  \]
\item
  次を満たす整数 \(n\ge\max(N,1)\) が存在する。
  \[
  \begin{aligned}
  \gcd\!&\left(\left\lfloor\xi(5/2)^n\right\rfloor,40112098026\right)>1,\\
  40112098026&=2\cdot3\cdot7\cdot11\cdot13\cdot17\cdot19\cdot23\cdot29\cdot31.
  \end{aligned}
  \tag{2}\label{eq:2}
  \]
\end{enumerate}
\label{thm-end:1}
\end{theorem}
\end{samepage}

% Markdown AST block 9: Para
(i)と(ii)の添字はそれぞれ別々に量化されており、共通の添字の存在は主張していない。(2)の法は素数 \(5\) を含まない。

% Markdown AST block 10: Para
ここから以下では \(\mathbb N=\{0,1,2,\ldots\}\) とし、整数 \(q\) に対して \(\gcd(q,M)=\gcd(|q|,M)\) を用いる。任意の \(N\) に対する定式化は、条件を満たす正の添字全体が無限集合をなすことを正確に表している。正整数が\emph{合成数}であるとは、\(uv\) と書け、整数の因子が \(u,v>1\) を満たすことをいう。

% Markdown AST block 11: Corollary 2.
\begin{samepage}
\begin{corollary}\leavevmode\label{thm:2}
任意の \(\xi>0\) に対して、次の各数列は合成数の項を無限に含む。
\[
\bigl(\lfloor\xi(7/5)^n\rfloor\bigr)_{n\ge1}\qquad \text{and}\qquad \bigl(\lfloor\xi(5/2)^n\rfloor\bigr)_{n\ge1}
\]
\label{thm-end:2}
\end{corollary}
\end{samepage}

% Markdown AST block 12: Para
比 \(7/5\) については、証明書から待ち時間の明示的な上界も得られる。次のようにおく。
\[
L(P)=\min\{h\ge0:5^h\ge P+2\}.
\]

% Markdown AST block 13: Theorem 3 (Quantitative Theorem).
\begin{samepage}
\begin{theorem}[定量的定理]\leavevmode\label{thm:3}
\(\xi>0\) が \(P=\lfloor\xi\rfloor>13\) を満たすとする。このとき、次を満たす整数 \(n\) が存在し、その項は合成数である。
\[
\begin{aligned}
0&\le n\le428+12L(P)\\
\text{and}\qquad
\gcd\!&\left(\left\lfloor\xi(7/5)^n\right\rfloor,4290\right)>1,
\end{aligned}
\tag{3}\label{eq:3}
\]
したがって、任意の \(N\in\mathbb N\) で \(Q_N=\lfloor\xi(7/5)^N\rfloor>13\) を満たすものに対し、ある \(\lfloor\xi(7/5)^n\rfloor\) は合成数となり、その添字は
\[
N\le n\le N+428+12L(Q_N)
\]
を満たす。
\label{thm-end:3}
\end{theorem}
\end{samepage}

% Markdown AST block 14: Para
関連する小数部分の研究には、ある \(\xi>0\) が
\[
0\le\{\xi(3/2)^n\}<1/2\qquad\text{for all }n\ge0.
\]
を満たすかどうかというMahlerの問題 \hyperref[ref-m68]{M68, p.313} がある。Flatto、Lagarias、Pollington \hyperref[ref-flp]{FLP, Theorem 1.4, p.129} は、互いに素な整数 \(p>q\ge2\) と任意の \(\xi>0\) に対して、
\[
\limsup_{n\to\infty}\{\xi(p/q)^n\}
-\liminf_{n\to\infty}\{\xi(p/q)^n\}\ge1/p.
\]
を証明した。Dubickas \hyperref[ref-dub08]{Dub08} は、これらの小数部分がすべて指定された区間または区間の和集合に入るかどうかについて、さらに結果を得ている。このような閉じ込めの問題は、本論文で小数部のセルを用いる背景となる。\hyperref[ref-dn]{DN, p.636} で説明されているとおり、値域の評価だけでは、合成数性を導くのに十分な、偶数の項が無限に存在するという結論は得られない。

% Markdown AST block 15: Para
他の実数の底について、AlkauskasとDubickas \hyperref[ref-ad04]{AD04} は、正の冪の整数部分がすべて合成数となるPisot数や、整数部分の数列が任意の整数等差数列と無限回交わる超越数を構成した。これらは特定の底を構成する結果である。本論文では底を指定された有理数に固定し、係数 \(\xi\) をすべての正の実数にわたって動かす。

% Markdown AST block 16: Para
\textbf{手法。} 証明では、既存の非周期性による障害と、明示的な有限証明書を組み合わせる。繰り上がり列 \(c_n=bQ_{n+1}-aQ_n\) の非周期性は、すでに \hyperref[ref-dn]{DN, Lemma 2, p.638} および \hyperref[ref-n]{N, Lemma 1.7, pp.28--29} に含まれており、操作語への制約を用いて周期性を強制する方法にも \hyperref[ref-dn]{DN, pp.638--639 and 647} に先例がある。本論文では、この障害について整数だけを用いた証明を示し、グラフの構成、全辺を対象とする位相判定、および尾部保存の議論を明示する。1ステップの構成（第3節）では、各辺に一つの繰り上がりを記録し、小数部のセルを細分することで(1)を証明する。語ラベル付きの構成（第4節）では、辺に有限の繰り上がり語を付し、出力を変えずに決定的な道を縮約し、各語のすべての中間時刻において剰余を検査しながら、一度に一つずつ新しい素数でグラフを持ち上げる。これによって(2)を証明し、(1)の第2の証明も得る。完全な1ステップグラフ上の有限ランク証明書と、非周期性補題の有限版を組み合わせると、定理3が得られる（第6節）。

% Markdown AST block 17: Para
Akiyama、Frougny、Sakarovitch \hyperref[ref-afs08]{AFS08} は、有理数を底とする表現と、その算術を処理する有限トランスデューサを展開した。彼らの非周期性の結果では、分母の増大する冪による整除性を用いる（\hyperref[ref-afs08]{AFS08, Lemma 8, p.65 and Proposition 26, p.75}）。彼らの整数表現の言語は正則ではない（\hyperref[ref-afs08]{AFS08, Corollary 7, p.65}）。これに対し、本論文の有限グラフは、素数を避ける軌道が満たす必要のある剰余と小数部のセルの制約を記録する。そのような軌道は必ずウォークを与えるが、ウォークが軌道を表すとは限らない。この区別により、表現言語全体の有限オートマトンを必要とせずに、有限証明書を得ることができる。

% Markdown AST block 18: Para
Stephanのプレプリント \hyperref[ref-s26]{S26}（2026年8月15日付）は、\hyperref[ref-dub06]{Dub06, Conjecture 2} の \(b=7\) の場合、すなわち任意の \(\xi>0\) に対して数列 \(\lfloor\xi7^n\rfloor\) が合成数の項を無限に多く含むことを証明し、同じ証明書により、底 \(7\) の無限右切り詰め素数（infinite right truncatable prime）が存在しないことも示している。その手法は本論文の手法と密接に関連する。そのグラフの頂点は、\(M=\prod_{p\le31}p\) を法とする剰余のうち \(M\) と互いに素なものであり、末尾に付け加える各桁に一つの辺が対応する。証明書が主張するのは、両側無限の道の上にない状態を取り除いた後、どの強連結成分も相異なる二つの閉路を含まないことである。このとき、すべての無限の道は最終的に周期的となり、最終的に周期的な桁の語は、Dubickasの定理 \hyperref[ref-dub09]{Dub09, Theorem 4}（定理46の後で述べる）によって合成数の項を無限に多く与える。証明書は、一度に一つずつ素数を加えていく法の梯子を射影とランク関数を用いて登ることで得られ、Leanで検証されている。その有限計算は \nolinkurl{native_decide} で評価される（\hyperref[ref-s26]{S26, §§4--6 and Appendix A}）。

% Markdown AST block 18 (continued): Para
本論文は \hyperref[ref-s26]{S26} といくつかの要素を共有している。仮想的な素数を避ける軌道がすべて無限ウォークを与えるような有限の剰余グラフ、認証済みのグラフを一度に一つずつ新しい素数で持ち上げること（第4.4節）、グラフアルゴリズムの形式的証明の代わりに検査器が確認するランク関数（第6.2節および第9.4節）、そして語ラベル付きの有限評価を同じく \nolinkurl{native_decide} で受け入れるLeanによる検証（第9.4節）である。一方、二つの構成は三つの点で異なる。第一に、\hyperref[ref-s26]{S26} のグラフは剰余だけからなる。非整数の底では、すべての繰り上がり語が軌道によって実現されるわけではない。本論文の状態は小数部のセルも記録し、グラフは外側近似である。すなわち、すべての軌道はウォークを与えるが、ウォークが軌道を表すとは限らない。第二に、どちらの証明書も、すべての無限ウォークのラベル列が周期的になる成分（\hyperref[ref-s26]{S26} では単純閉路、本論文では補題7の位相判定条件に合格する成分）を特定するが、それらを正反対に扱う。底 \(7\) では、そのような成分に実際の素数を避ける軌道が乗る。\hyperref[ref-s26]{S26} のグラフの核は法 \(M\) でもなお 371702 個の状態をもち、\hyperref[ref-dj22]{DJ22, Theorem 1.2, p.167} からは、任意の有限素数集合に対して適切な正の係数を選ぶとすべての項がその素数集合を避けることが従うので、回避不能集合は存在しない。\hyperref[ref-s26]{S26} はこれらの成分を残し、その周期的な語にDubickasの定理を適用する。非整数の底では、補題4により、どの軌道もそのような成分にとどまり続けることはないので、その成分は削除される。本論文の証明書はグラフを空グラフにまで縮小し、それによって定理1の回避不能集合を与える。\hyperref[ref-dj22]{DJ22} の結果は \(7/5\) や \(5/2\) には適用されない。第三に、すべての中間時刻で素数条件を課す語ラベル付き縮約、定理3の待ち時間の上界、および第7節と第8節の正規形定理と障害定理は、\hyperref[ref-s26]{S26} には含まれていない。\hyperref[ref-s26]{S26} の計算は手法上の比較対象であり、本論文の証明が依存するものではない。本論文の具体的な貢献は、シフトを加えない \(7/5\) および \(5/2\) の数列に対する有限証明書とその厳密な検証、すべての中間時刻で素数条件を課す語ラベル付き縮約、待ち時間の上界、および以下で述べる構造的結果である。

% Markdown AST block 19: Para
\textbf{手法の適用範囲。} 第7節と第8節では、この証明書による手法で何が可能で、何が不可能かを調べる。操作を厳密に固定し（初期完全グラフ、位相除去、厳密な縮約、セル細分、法の拡大、繰り上がりの制限）、正規形定理を証明する。すなわち、あるパラメータの選択に対して有限証明書が存在することと、あるセル数 \(K\) とある有限素数集合 \(S\) が存在し、その素数はいずれも \(2ab\) を割らず、これらのパラメータに対して、一つの明示的なグラフのすべての巡回的強連結成分が位相判定に合格することは同値である。縮約によって認証可能なクラスは広がらず、法に含まれる素数冪は取り除くことができ、\(2ab\) を割る素数は繰り上がりの文字集合への制限に帰着する。この正規形では結論に現れる素数集合が変わり、\(2ab\) を割るすべての素数を含むものとなる。続く二つの障害定理により、\(a\ge2\operatorname{rad}M\) のとき、およびある2文字の繰り上がり語が法 \(M\) で単元剰余を許すときに、この手法が停止しないことを示す。第2の障害から必要条件 \(a\le2R\,J(N_0)\) が得られる。ここで、\(R\) は \(M\) を割り \(2ab\) を割らない素数の積、\(N_0\) は \(ab\) を割る奇素数の積、\(J\) はJacobsthal関数である。初等的な評価 \(J(N)\le3^{\omega(N)}\) から \(R\ge a^{1-o(1)}\) が得られる。これらの障害定理は、固定した手続きが停止しないという主張である。いかなる算術的命題に対する反例でもなく、ある底について合成数の項が無限に存在しないと主張するものでもない。

% Markdown AST block 20: Para
\textbf{検証。} 採用した有限計算はすべて、数学的仕様から科学的中核（辺の列挙、成分分解、認証、縮約、持ち上げ）をそれぞれ別々に実装した二つのプログラムで実行し、完全なグラフデータを要素ごとに比較した。第9節で、それらの作成経緯と比較を説明する。Leanによる証明は、定理1の両方の主張と系2を確立する。\(7/5\) の1ステップ証明では、カーネルで検査された有限証明書を用いる。語ラベル付きの証明では、カーネルで検査された健全性定理と \nolinkurl{native_decide} による有限評価を用いる。第9.4節で、形式化の正確な範囲と、後者の証明で追加されるコンパイラおよびネイティブ評価への信頼を述べる。

% Markdown AST block 21: Para
第2節では非周期性補題を証明する。第3節では1ステップ外側グラフと有限認証を、第4節では語ラベル付き証明書を示す。第5節で(1)と(2)の有限計算を述べる。第6節では定理3を証明する。第7節と第8節には構造的結果と障害に関する結果をまとめる。第9節では検証と形式化の正確な範囲を説明する。第10節では付録に証明を掲げる追加の結果を述べ、第11節で未解決問題を挙げる。

% BEGIN PUBLIC ADDITION figure:certificate-flow
二つの有限認証手順は、以下に示す尾部保存の議論を共有する。

\begin{figure}[htbp]
\centering
\begin{tikzpicture}[paperfig]
\node[fbox,text width=11.5cm] (orbit) at (0,0)
 {\figtext{Assume a true orbit eventually avoids the chosen primes}{真の軌道が選んだ素数を最終的にすべて回避すると仮定}\\
 $Q_n=\lfloor\xi(a/b)^n\rfloor,\qquad c_n=bQ_{n+1}-aQ_n$};
\node[fbox,text width=10cm] (graph) at (0,-1.45)
 {\figtext{Finite outer graph: every such orbit gives a walk}{有限外側グラフ：そのような軌道はすべてウォークを与える}\\
 \figtext{congruence + cell compatibility + carry alphabet}{合同条件＋セル整合条件＋キャリー集合}};
\draw[fedge] (orbit)--(graph);
\node[fbox,text width=10cm] (remove) at (0,-3)
 {\figtext{Decompose into SCCs; test all internal letters}{SCC分解し、内部辺の全文字について位相判定}\\
 \figtext{Remove certified cyclic SCCs and acyclic parts}{認証された巡回SCCと非巡回部分を除去}};
\draw[fedge] (graph)--(remove);
\node[fbox,text width=5cm] (ref) at (-3.2,-5.2)
 {\textbf{\figtext{One-step method}{一段階法}}\\[4pt]
 $K\longmapsto fK$\\\figtext{Refine retained cells; rebuild edges}{残存セルを細分し、辺を再構成}};
\node[fbox,text width=5cm] (word) at (3.2,-5.2)
 {\textbf{\figtext{Word method}{語ラベル法}}\\[4pt]
 $K\ \text{\figtext{fixed}{は固定}}$\\\figtext{Contract paths; lift by the next prime}{経路を縮約し、次の素数で持ち上げ}};
\draw[fedge] (remove.south) -- ++(0,-.32) -| (ref.north);
\draw[fedge] (remove.south) -- ++(0,-.32) -| (word.north);
\draw[fedge,dashed] (ref.west) -- ++(-.5,0) |- (remove.west);
\draw[fedge,dashed] (word.east) -- ++(.5,0) |- (remove.east);
\node[fbox,text width=11.5cm] (empty) at (0,-7.2)
 {\figtext{If the retained graph becomes empty: contradiction}{残存グラフが空になれば矛盾}\\
 \figtext{A true prime-avoiding tail survives every prescribed operation}{真の素数回避軌道の尾部は各操作を生き残る}};
\draw[fedge] (ref.south)-- ++(0,-.35) -| ([xshift=-3.2cm]empty.north);
\draw[fedge] (word.south)-- ++(0,-.35) -| ([xshift=3.2cm]empty.north);
\end{tikzpicture}

\caption{一段階法と語ラベル法の論理的な流れ。軌道からグラフへの矢印は外側近似への包含を表す。一段階法ではセル細分を、語ラベル法ではセル数を固定した縮約と素数による持ち上げを用いる。破線矢印は反復を表す。}
\label{fig:certificate-flow}
\end{figure}
% END PUBLIC ADDITION figure:certificate-flow

% Markdown AST block 22: 2. Carries and nonperiodicity
\Needspace{6\baselineskip}
\section{繰り上がりと非周期性}\label{sec:2}

% Markdown AST block 23: Para
一般的な議論を通じて、
\[
a,b\in\mathbb N,\qquad a>b\ge2,\qquad \gcd(a,b)=1,\qquad \xi>0,
\]
とし、\(r=a/b\) とおく。除法は \(\mathbb R\) におけるものとする。\(n\ge0\) に対して、
\[
x_n=\xi r^n,\qquad Q_n=\lfloor x_n\rfloor\in\mathbb Z,
\qquad \theta_n=x_n-Q_n\in[0,1),
\tag{4}\label{eq:4}
\]
および符号付き整数の繰り上がり
\[
c_n=bQ_{n+1}-aQ_n=a\theta_n-b\theta_{n+1}.
\tag{5}\label{eq:5}
\]
を定義する。したがって、
\[
1-b\le c_n\le a-1,\qquad bQ_{n+1}=aQ_n+c_n.
\tag{6}\label{eq:6}
\]
が成り立つ。実数における厳密な境界 \(-b<c_n<a\) から、(6)の整数境界が従う。別の添字の付け方では \(e_{n+1}=c_n\) と書け、\(e_0\) の値は必要ない。\(x_{n+1}=rx_n>x_n\) であるから、数列 \((Q_n)\) は単調非減少であり、\(Q_n\ge Q_0\) がすべての \(n\) に対して成り立つ。

% Markdown AST block 24: Para
また、\(Q_n\to+\infty\) を、次の明示的な意味で用いる。
\[
\forall T\in\mathbb Z\ \exists N\in\mathbb N\ \forall n\ge N,
\qquad Q_n>T.
\tag{7}\label{eq:7}
\]
実際、\(r=1+\delta\) と書き、\(\delta>0\) とすると、Bernoulliの不等式から \(r^n\ge1+n\delta\) が得られ、
\(Q_n>\xi(1+n\delta)-1\)
となる。

% Markdown AST block 25: Lemma 4 (nonperiodicity).
\begin{samepage}
\begin{lemma}[非周期性]\leavevmode\label{thm:4}
繰り上がり列 \((c_n)\) は最終的周期的ではない。すなわち、\(N\ge0\) と \(L\ge1\) であって、
\[
c_{n+L}=c_n
\]
がすべての \(n\ge N\) に対して成り立つようなものは存在しない。
\label{thm-end:4}
\end{lemma}
\end{samepage}

% Markdown AST block 26: Proof.
\begin{proof}\leavevmode\label{proof:4}
まず、整数列 \((D_k)_{k\ge0}\) が \(bD_{k+1}=aD_k\) を満たすとする。帰納法により、
\[
b^kD_k=a^kD_0.
\]
を得る。\(a^k\) と \(b^k\) は互いに素なので、\(b^k\mid D_0\) がすべての \(k\) に対して成り立つ。\(b\ge2\) より、その符号に関する仮定なしに \(D_0=0\) が従う。

% Markdown AST block 27: Para
次に、\(c\) が周期 \(L\) を \(N\) 以降でもつと仮定する。\(m\ge N\) を固定し、
\[
D_k=Q_{m+k+L}-Q_{m+k}.
\]
とおく。(6)の二つの漸化式の差を取り、周期性を用いると、\(bD_{k+1}=aD_k\) が得られる。先ほどの考察から、\(Q_{m+L}=Q_m\) が従う。これはすべての \(m\ge N\) について成り立つので、特に
\(Q_{N+kL}=Q_N\)
がすべての \(k\) に対して成り立ち、(7)に矛盾する。\qedhere
\end{proof}

% Markdown AST block 28: Para
この証明では、整数の漸化式と増大性だけを用いている。仮定 \(b\ge2\) と \(\gcd(a,b)=1\) は不可欠である。比が整数の場合には、適切な係数に対して繰り上がりが恒等的にゼロとなりうる。先に引用した既存の非周期性補題も、最終的周期性を排除する。尾部を新しい正の係数 \(\xi r^N\) に吸収できるためである。

% Markdown AST block 29: Para
(6)を反復した形も記しておく。\(C_0=0\) および
\[
C_{i+1}=aC_i+b^ic_{u+i},
\]
とすると、帰納法によって
\[
b^iQ_{u+i}=a^iQ_u+C_i\qquad(i\ge0),
\tag{8}\label{eq:8}
\]
が成り立つ。これは、
\[
b^{i+1}Q_{u+i+1}=b^i(aQ_{u+i}+c_{u+i})=a(a^iQ_u+C_i)+b^ic_{u+i}.
\]
による。

% Markdown AST block 30: 3. One-step outer graphs and finite certification
\Needspace{6\baselineskip}
\section{1ステップ外側グラフと有限認証}\label{sec:3}

% Markdown AST block 31: 3.1 The graphs
\subsection{グラフ}\label{sec:3.1}

% Markdown AST block 32: Para
整数 \(M\ge1\) と \(K\ge1\) を固定する。\(M\) と互いに素な剰余と、小数部のセルを記録する状態は、
\[
V(M,K)=\{(q,j)\in\mathbb Z^2:0\le q<M,\ \gcd(q,M)=1,\ 0\le j<K\}.
\tag{9}\label{eq:9}
\]
である。\(M=1\) のとき、剰余座標は \(q=0\) である。\(W\subseteq V(M,K)\) に対して、有向ラベル付きグラフ \(G_K(W)\) を \(W\) 上に定義する。辺 \((q,j)\xrightarrow{e}(z,k)\) が存在することと、両端点が \(W\) に属して次の符号付き整数条件が成り立つことは同値である。
\[
\begin{aligned}
1-b&\le e\le a-1, &&\text{(E1)}\\
bz&\equiv aq+e\pmod M, &&\text{(E2)}\\
aj-b(k+1)&<eK<a(j+1)-bk. &&\text{(E3)}
\end{aligned}
\tag{10}\label{eq:10}
\]
端点が同じでもラベルが異なる辺は区別する。\emph{完全な}1ステップグラフを \(G(M,K)=G_K(V(M,K))\) と書く。

% Markdown AST block 33: Para
実際の軌道の状態を
\[
\sigma_K(n)=\bigl(Q_n\bmod M,\ \lfloor K\theta_n\rfloor\bigr),
\tag{11}\label{eq:11}
\]
で定義する。剰余は \(\{0,\ldots,M-1\}\) 内の代表元で表す。

% Markdown AST block 34: Proposition 5 (orbit inclusion).
\begin{samepage}
\begin{proposition}[軌道の包含]\leavevmode\label{thm:5}
\(\sigma_K(n),\sigma_K(n+1)\in W\) ならば、
\[
\sigma_K(n)\xrightarrow{c_n}\sigma_K(n+1)
\]
は \(G_K(W)\) の辺である。
\label{thm-end:5}
\end{proposition}
\end{samepage}

% Markdown AST block 35: Proof.
\begin{proof}\leavevmode\label{proof:5}
条件(E1)は(6)そのものであり、(5)を法 \(M\) で還元すると(E2)が得られる。\(j=\lfloor K\theta_n\rfloor\) および \(k=\lfloor K\theta_{n+1}\rfloor\) と書く。このとき、
\[
j\le K\theta_n<j+1,\qquad k\le K\theta_{n+1}<k+1,
\]
であるから、
\[
aj-b(k+1)<aK\theta_n-bK\theta_{n+1}=Kc_n<a(j+1)-bk.
\]
となる。これは(E3)である。小数部分がゼロとなる場合でも、両方の不等式は厳密である。\qedhere
\end{proof}

% Markdown AST block 36: Para
この証明では、実数の恒等式(5)と半開セルだけを用いており、\(Q_n\) の整数性は使っていない。この点は第7節で用いる。グラフは外側近似である。\(M\) の素因数をすべて最終的に避ける各軌道はグラフ内のウォークを与えるが、逆にグラフのウォークが実現可能であるという主張は用いない。実数 \(u,v\) が \(j\le u<j+1\)、\(k\le v<k+1\)、\(au-bv=eK\) を満たすという条件について、条件(E3)はそのような組が存在するための必要十分条件であり、隣接する二つのセルの整合性のみを表す。

% Markdown AST block 37: Para
これらの条件により、整数演算による完全な列挙が可能になる。\(j,e\) を固定すると、(E3)と \(0\le k<K\) は、
\[
\max\!\left(0,\left\lfloor\frac{aj-eK}{b}\right\rfloor\right)
\le k\le
\min\!\left(K-1,\left\lfloor\frac{a(j+1)-eK-1}{b}\right\rfloor\right).
\tag{12}\label{eq:12}
\]
と同値である。例えば、\(A=aj-eK\) ならば、\(A<b(k+1)\) は \(\lfloor A/b\rfloor\le k\) と同値である。他方の境界は、厳密な整数不等式 \(bk<B\) を \(bk\le B-1\) に置き換えると得られる。(12)に現れる床関数は、分子が負の場合も含め、すべて数学的な床関数である。ゼロ方向への切り捨てを用いると、列挙は誤ったものとなる。整数区間が空ならば、辺は一つも生じない。

% Markdown AST block 38: Para
\(b\) の法 \(M\) における可逆性は仮定しない。(E2)のすべての解の求め方を具体的に示すため、\(g=\gcd(b,M)\) および \(t=aq+e\) とおく。\(g\mid t\) でなければ解は存在しない。この整除性が成り立つならば、約分した合同式
\[
(b/g)z_0\equiv t/g\pmod{M/g},\qquad 0\le z_0<M/g.
\]
を解く。\(M/g=1\) では \(z_0=0\) を用いる。それ以外の場合、係数 \(b/g\) は法 \(M/g\) で可逆である。代表元全体は
\[
z=z_0+\ell(M/g),\qquad 0\le\ell<g.
\tag{13}\label{eq:13}
\]
である。その中から、\(\gcd(z,M)=1\) と \((z,k)\in W\) を満たすものをちょうど残す。同値な方法として、1ステップの実装のように、すべての単元剰余 \(z\) を列挙し、\(bz\bmod M\) の値でまとめてもよい。\(M=4290\) の場合には \(\gcd(5,4290)=5\) なので、\(5\) による法 \(4290\) での除算はできない。

% BEGIN PUBLIC ADDITION figure:cell-intersection
整合する一組のセルを見ると、(E3)の両方の不等号が厳密である理由が分かる。

\begin{figure}[htbp]
\centering
\begin{tikzpicture}[paperfig]
\begin{scope}[x=2.3cm,y=2.3cm]
\draw[fedge] (.65,.65)--(2.35,.65);
\node at (1.65,.27) {$u=K\theta_n$};
\draw[fedge] (.65,.65)--(.65,2.35) node[above] {$v=K\theta_{n+1}$};
\fill[black!5] (1,1) rectangle (2,2);
\draw (1,2)--(1,1)--(2,1);
\draw[dashed] (1,2)--(2,2)--(2,1);
\draw[femph] (1.2,1)--(1.6,2);
\fill[figblue] (1.2,1) circle(1.4pt);
\draw[figblue,fill=white,line width=.8pt] (1.6,2) circle(1.7pt);
\node[flabel,rotate=63.435] at (1.4,1.52) {$5u-2v=4$};
\foreach \x in {1,2}{\draw (\x,.62)--(\x,.68);\node[below] at(\x,.62){$\x$};}
\foreach \y in {1,2}{\draw (.62,\y)--(.68,\y);\node[left] at(.62,\y){$\y$};}
\end{scope}
\node[align=left,anchor=west,text width=6.5cm] at (6,4.3)
 {$a=5,\ b=2,\ K=4$\\[4pt]
 $j=k=1,\quad e=1$\\[8pt]
 $[j,j+1)\times[k,k+1)$\\[4pt]
 \figtext{Solid: included lower boundaries}{実線：含まれる下側境界}\\
 \figtext{Dashed: excluded upper boundaries}{破線：含まれない上側境界}};
\node[fbox,align=center,text width=6.3cm] at (9.2,1.8)
 {$aj-b(k+1)<eK<a(j+1)-bk$\\[5pt]
 $1<4<8$\\[4pt]
 \figtext{Both endpoint bounds are strict}{両端の不等号は厳密}};
\end{tikzpicture}

\caption{\(a=5,b=2,K=4,j=k=1,e=1\)におけるセル整合条件。直線は半開正方形と、\(u=(4+2v)/5\), \(1\le v<2\)の部分で交わる。下端は含まれ、上端は含まれない。セル対の判定は局所的であり、軌道全体の存在を主張するものではない。}
\label{fig:cell-intersection}
\end{figure}
% END PUBLIC ADDITION figure:cell-intersection

% Markdown AST block 39: 3.2 Components with periodic output
\Needspace{6\baselineskip}
\subsection{周期的出力をもつ成分}\label{sec:3.2}

% Markdown AST block 40: Para
強連結成分（SCC）とは、有向道による相互到達可能性に関する同値類であり、長さゼロの道も許す。内部辺を含むSCCを\emph{巡回的}と呼ぶ。したがって、1頂点からなるSCCが巡回的であることと、自己ループをもつことは同値である。

% Markdown AST block 41: Lemma 6.
\begin{samepage}
\begin{lemma}\leavevmode\label{thm:6}
有限有向グラフの任意の無限ウォークは、最終的に一つの巡回的SCC内にとどまる。
\label{thm-end:6}
\end{lemma}
\end{samepage}

% Markdown AST block 42: Proof.
\begin{proof}\leavevmode\label{proof:6}
ウォークは、あるSCCを出て別のSCCに移った後には、元のSCCに戻ることができない。戻れたならば、二つの成分が相互に到達可能となるからである。したがって、成分を移る回数は有限回に限られる。最後の成分は無限の尾部を含むので、内部辺をもつ。\qedhere
\end{proof}

% Markdown AST block 43: Para
\(C\) を \(G_K(W)\) の巡回的SCCとする。根 \(\rho\in C\) を選び、\(h(v)\) を、\(\rho\) から \(v\) へ至り \(C\) 内にとどまる道の最短の長さとする。次のように定義する。
\[
d=\gcd\{\,|h(u)+1-h(v)|:
u\xrightarrow{e}v\text{ is an internal edge of }C\,\}.
\tag{14}\label{eq:14}
\]
\(C\) 内には空でない閉ウォークが存在する。差 \(h(u)+1-h(v)\) を長さ \(\ell>0\) の閉ウォークに沿って足し合わせると、\(\ell\) を得る。したがって、これらの差がすべてゼロとなることはなく、\(d>0\) である。すべての内部辺は、
\[
h(v)\equiv h(u)+1\pmod d.
\tag{15}\label{eq:15}
\]
を満たす。同じ和を取る議論から、どの閉ウォークの長さも \(d\) で割り切れる。したがって、空でない任意の閉ウォークに沿う最初の \(d\) 個の出発位相は、法 \(d\) のすべての剰余を実現する。特に、各位相には少なくとも一つの出る内部辺がある。

% Markdown AST block 44: Para
各位相 \(s\in\{0,\ldots,d-1\}\) に対して、始点 \(u\) が \(h(u)\equiv s\pmod d\) を満たす\emph{すべての}内部辺のラベルを集める。それぞれの集合が単集合であるとき、\(C\) を認証し、その集合を \(\{\lambda_s\}\) と書く。

% Markdown AST block 45: Lemma 7 (phase criterion).
\begin{samepage}
\begin{lemma}[位相判定条件]\leavevmode\label{thm:7}
認証された成分内にとどまる任意の無限ウォークのラベル列は周期的であり、\(d\) を周期にもつ。
\label{thm-end:7}
\end{lemma}
\end{samepage}

% Markdown AST block 46: Proof.
\begin{proof}\leavevmode\label{proof:7}
ウォーク \(v_0\xrightarrow{\ell_0}v_1\xrightarrow{\ell_1}\cdots\) に対して、(15)から
\[
h(v_i)\equiv h(v_0)+i\pmod d.
\]
が得られる。したがって、
\[
\ell_i=\lambda_{(h(v_0)+i)\bmod d},
\qquad \ell_{i+d}=\ell_i.\qedhere
\]
\end{proof}

% Markdown AST block 47: Para
この判定は、多重辺も含めてすべての内部辺を対象とし、木の辺や選択した一つの閉路だけを対象とするものではない。\(C\) から出る辺は、ウォークが \(C\) 内にとどまる段階では無関係である。後の補題12では、任意の整数値関数 \(h\) に対して健全性を証明する。完全な判定を得るには、\(h\) を上記のように共通の根からの道の長さとして選ぶ。最短の道である必要はない。定理21と系22により、判定結果は根やこれらの道の選択に依存しないことが示される。

% Markdown AST block 48: Para
(14)の整数 \(d\) と、位相ラベル語の原始周期 \(p\) は区別する。結果を報告する際には、\((\lambda_0,\ldots,\lambda_{d-1})\) を最短の反復語で置き換え、さらにその巡回シフトのうち辞書式順序で最小のものに置き換える。この語の長さは \(p\mid d\) であり、\(p<d\) となる場合もある。この正規化は比較と表示のためのものであり、補題7は周期の最小性も特定の開始位相も必要としない。

% Markdown AST block 49: 3.3 Tail preservation and finite certification
\Needspace{6\baselineskip}
\subsection{尾部保存と有限認証}\label{sec:3.3}

% Markdown AST block 50: Para
\(R_K(W)\) を、\(G_K(W)\) の巡回的SCCのうち位相判定条件を満たさないものすべての和集合とする。それ以外の頂点は除去する。

% Markdown AST block 51: Lemma 8 (retained tail).
\begin{samepage}
\begin{lemma}[尾部の保持]\leavevmode\label{thm:8}
\(\sigma_K(n)\in W\) が十分大きなすべての \(n\) に対して成り立つならば、\(\sigma_K(n)\in R_K(W)\) が十分大きなすべての \(n\) に対して成り立つ。
\label{thm-end:8}
\end{lemma}
\end{samepage}

% Markdown AST block 52: Proof.
\begin{proof}\leavevmode\label{proof:8}
命題5により、与えられた尾部はラベル \(c_n\) をもつ無限ウォークとなる。補題6により、このウォークは最終的に一つの巡回的SCC内にとどまる。その成分が認証されていたならば、補題7により \(c_n\) が最終的周期的となり、補題4に矛盾する。したがって、その成分は保持されなければならない。\qedhere
\end{proof}

% Markdown AST block 53: Para
整数 \(f\ge2\) に対して、細分を次で定義する。
\[
\operatorname{Refine}_f(U)
=\{(q,fj+i):(q,j)\in U,\ 0\le i<f\}.
\tag{16}\label{eq:16}
\]

% Markdown AST block 54: Lemma 9 (refinement).
\begin{samepage}
\begin{lemma}[細分]\leavevmode\label{thm:9}
\(\sigma_K(n)\in U\) ならば、
\[
\sigma_{fK}(n)\in\operatorname{Refine}_f(U).
\]
\label{thm-end:9}
\end{lemma}
\end{samepage}

% Markdown AST block 55: Proof.
\begin{proof}\leavevmode\label{proof:9}
\(j=\lfloor K\theta_n\rfloor\) から、
\[
fj\le fK\theta_n<f(j+1).
\]
を得る。したがって、
\[
\lfloor fK\theta_n\rfloor=fj+i
\]
が、一意な整数 \(0\le i<f\) に対して成り立つ。剰余座標は変わらない。\qedhere
\end{proof}

% Markdown AST block 56: Para
\(K_0\ge1\) と \(W_0=V(M,K_0)\) から出発し、
\[
K_t=K_0f^t,\qquad U_t=R_{K_t}(W_t),\qquad
W_{t+1}=\operatorname{Refine}_f(U_t).
\tag{17}\label{eq:17}
\]
を計算する。新しい各段階では、(10)から辺を再計算する。

% Markdown AST block 57: Theorem 10 (finite certification).
\begin{samepage}
\begin{theorem}[有限認証]\leavevmode\label{thm:10}
\(U_s=\varnothing\) がある有限の \(s\) に対して成り立つならば、任意の \(\xi>0\) と任意の \(N\in\mathbb N\) に対して、ある \(n\ge\max(N,1)\) が存在し、
\[
\gcd\!\left(\left\lfloor\xi(a/b)^n\right\rfloor,M\right)>1.
\tag{18}\label{eq:18}
\]
を満たす。
\label{thm-end:10}
\end{theorem}
\end{samepage}

% Markdown AST block 58: Proof.
\begin{proof}\leavevmode\label{proof:10}
逆に、ある \(\xi>0\) に対して、整数 \(Q_n\) が \(M\) と互いに素であることが、十分大きなすべての \(n\) について成り立つと仮定する。(9)と(11)により、対応する状態は \(W_0\) に属する。補題8により、さらに先の尾部は \(U_0\) に入り、補題9により、より細かい尺度で見た同じ尾部は \(W_1\) に入る。この二つの含意を有限個の段階にわたって反復すると、ある尾部が \(U_s\) に入ることになるが、この集合は空である。この矛盾により(18)が証明される。\qedhere
\end{proof}

% Markdown AST block 59: Para
尾部の開始時刻は各段階で遅くなりうるし、\(\xi\) に依存してよい。一様な開始時刻も、すべての有限接頭部の保存も必要ない。特に、細分は真の軌道の状態について証明されており、粗いグラフのすべてのウォークが細かい尺度への持ち上げをもつとは仮定していない。グラフ自体は \(\xi\) を入力として受け取らない。すべての正の実数係数への移行は、これらの補題によって保証されるのであり、初期値の標本抽出や長い数値軌道の計算によるものではない。

% Markdown AST block 60: Para
定理10は、成功した有限計算の存在を仮定する。任意に \(a,b,M,K_0,f\) を選んだときの停止を主張するものではない。

% Markdown AST block 61: Para
法も段階ごとに拡大してよい。\(U\subseteq V(M,K)\) と整数 \(h,f\ge1\) に対し、\(M'=hM\)、\(K'=fK\) とおき、
\[\begin{aligned}
&U'=\{(q+\ell M,\ fj+i):(q,j)\in U,\ 0\le\ell<h,\ 0\le i<f,\\\
&\gcd(q+\ell M,hM)=1\}\subseteq V(M',K').
\end{aligned}\]
と定める。

% Markdown AST block 62: Lemma 11 (joint enlargement).
\begin{samepage}
\begin{lemma}[同時拡大]\leavevmode\label{thm:11}
\(\sigma_K(n)\in U\) が法 \(M\) の状態として成り立ち、\(\gcd(Q_n,hM)=1\) であるならば、
\[
\sigma_{K'}(n)\in U'
\]
が法 \(M'\) の状態として成り立つ。
\label{thm-end:11}
\end{lemma}
\end{samepage}

% Markdown AST block 63: Proof.
\begin{proof}\leavevmode\label{proof:11}
\(q=Q_n\bmod M\) および \(q'=Q_n\bmod hM\) とする。このとき \(q'\equiv q\pmod M\) と \(0\le q'<hM\) が成り立つので、\(q'=q+\ell M\) と書け、\(0\le\ell<h\) であり、
\[
\gcd(q',hM)=\gcd(Q_n,hM)=1.
\]
が成り立つ。セル座標は補題9で扱える。\qedhere
\end{proof}

% Markdown AST block 64: Para
途中のすべての法は最終的な法を割るので、最終的な法と互いに素な尾部は途中の各法とも互いに素である。したがって、補題9の代わりに補題11を用いれば、定理10の帰納法がそのまま通用する。

% BEGIN PUBLIC ADDITION figure:cell-refinement
次の例のように、セル細分は各真の軌道点を含むセルを保存する。

\begin{figure}[htbp]
\centering
\begin{tikzpicture}[paperfig]
\node[anchor=east] at(-.2,1.45){$K=4,\ j=1$};
\draw[line width=1pt] (0,1.45)--(9,1.45);
\foreach \x/\l in {0/{1/4},9/{1/2}} {\node[above=3pt] at(\x,1.45){$\l$};}
\fill (0,1.45) circle(2pt);\draw[fill=white] (9,1.45) circle(2pt);
\node[anchor=east] at(-.2,0){$fK=12$};
\foreach \x/\j in {0/3,3/4,6/5}{
 \draw[line width=1pt] (\x,0)--++(3,0);
 \fill (\x,0) circle(2pt);\draw[fill=white] (\x+3,0) circle(2pt);
 \node[below=3pt] at(\x+1.5,0){$\j$};
}
\node[above=3pt] at(3,0){$1/3$};\node[above=3pt] at(6,0){$5/12$};
\draw[fedge,dashed] (4.5,1.2)--(4.5,.3) node[midway,left] {$f=3$};
\draw[femph,dotted] (5.4,1.45)--(5.4,0);
\fill[figblue] (5.4,1.45) circle(2pt);\fill[figblue] (5.4,0) circle(2pt);
\node[above=3pt,align=center] at(5.4,1.5){$\theta=2/5$};
\node[fbox,text width=11.5cm] at (3.8,-1.25)
 {$\lfloor4\theta\rfloor=1,\qquad\lfloor12\theta\rfloor=4=3\cdot1+1$\\[5pt]
 $\sigma_K(n)\in U\quad\Longrightarrow\quad
 \sigma_{fK}(n)\in\operatorname{Refine}_f(U)$};
\node[align=center,text width=12cm] at(3.8,-2.5)
 {\figtext{The same true orbit point lies in exactly one finer cell.}{同じ真の軌道点は、細分後のちょうど一つのセルに入る。}\\
 \figtext{Edges are rebuilt; an arbitrary coarse walk need not lift.}{辺は再構成する。任意の粗いウォークが持ち上がるとは限らない。}};
\end{tikzpicture}

\caption{\(K=4\)のセル\([1/4,1/2)\)は、\(fK=12\)ではセル\(3,4,5\)に分かれる。点\(\theta=2/5\)はセル\(4\)に入る。剰余座標は変わらない。これは真の軌道状態についての包含であり、任意のグラフ上のウォークが持ち上がるという主張ではない。}
\label{fig:cell-refinement}
\end{figure}
% END PUBLIC ADDITION figure:cell-refinement

% Markdown AST block 65: 4. Word-labelled certificates
\Needspace{6\baselineskip}
\section{語ラベル付き証明書}\label{sec:4}

% Markdown AST block 66: 4.1 Word-labelled graphs and coverings
\subsection{語ラベル付きグラフと被覆}\label{sec:4.1}

% Markdown AST block 67: Para
\(\mathcal C\subseteq[1-b,a-1]\cap\mathbb Z\) を繰り上がりの有限集合とし、\emph{文字集合}と呼ぶ。\emph{語ラベル付きグラフ} \(G=(V,E)\) は、有限頂点集合 \(V\) と有限辺集合 \(E\) からなり、各辺 \(e=(u,w,v)\) について \(u,v\in V\) であり、\(w\in\mathcal C^{+}\) は空でない有限語である。\(|w|\ge1\) を辺の長さとする。辺は三つ組として同一視する。すなわち、同じ三つ組は一つの辺であり、端点が同じでも語が異なる辺は区別する。無限ウォークは
\[
v_0\xrightarrow{w_0}v_1\xrightarrow{w_1}\cdots
\]
であり、各辺は \((v_i,w_i,v_{i+1})\in E\) を満たす。その\emph{出力}は連結
\[
w_0w_1w_2\cdots\in\mathcal C^{\mathbb N},
\]
であり、長さゼロの辺が存在しないので無限語となる。第3節の1ステップグラフは、すべての語の長さが1である特別な場合である。

% Markdown AST block 68: Para
\((c_n)\) を真の軌道の繰り上がり列とし、\(n_0\ge0\) とする。\(G\) が\emph{\(n_0\) からの尾部を被覆する}といい、\(G\models n_0\) と書くのは、無限ウォーク \((v_i,w_i)_{i\ge0}\) が \(G\) に存在し、時刻列 \(t_0=n_0<t_1<t_2<\cdots\) とともに
\[
t_{i+1}=t_i+|w_i|,\qquad
w_i=(c_{t_i},c_{t_i+1},\ldots,c_{t_{i+1}-1})\qquad(i\ge0);
\]
を満たすときである。すなわち、そのウォークの出力が \((c_n)_{n\ge n_0}\) に等しいことをいう。以下の各操作 \(\Phi\)（除去、縮約、持ち上げ）について、\emph{尾部保存}を証明する。つまり、\(G\models n_0\) と、その操作に固有の算術的仮定が成り立つならば、\(\Phi(G)\models n_1\) がある \(n_1\ge n_0\) に対して成り立つことを示す。開始時刻は各操作で遅くなってよく、\(\xi\) に依存してよい。有限接頭部は保存されない。

% Markdown AST block 69: 4.2 The weighted phase criterion and removal
\Needspace{6\baselineskip}
\subsection{重み付き位相判定条件と除去}\label{sec:4.2}

% Markdown AST block 70: Para
\(C\) を \(G\) の巡回的SCC、すなわち内部辺 \((u,w,v)\) で \(u,v\in C\) を満たすものを少なくとも一つ含むSCCとする。\(h:C\to\mathbb Z\) を\emph{任意の}整数値関数とし、内部辺 \(e=(u,w,v)\) に対して、
\[
\begin{aligned}
\delta(e)&=h(u)+|w|-h(v),\\
d&=d_C=\gcd\{|\delta(e)|:e\text{ an internal edge of }C\}.
\end{aligned}
\]
とおく。

% Markdown AST block 71: Lemma 12 (weighted phase criterion).
\begin{samepage}
\begin{lemma}[重み付き位相判定条件]\leavevmode\label{thm:12}
(a) \(d\ge1\) であり、すべての内部辺は
\[
h(v)\equiv h(u)+|w|\pmod d;
\]
を満たす。\(C\) 内の空でない任意の閉ウォークの総語長は、\(d\) の正の倍数である。

% Markdown AST block 72: OrderedList
\begin{enumerate}
\def\labelenumi{(\alph{enumi})}
\setcounter{enumi}{1}
\tightlist
\item
  写像 \(\lambda:\mathbb Z/d\to\mathcal C\) が存在し、すべての内部辺 \(e=(u,w,v)\) とすべての文字位置 \(0\le i<|w|\) に対して、
  \[
  w_i=\lambda\bigl((h(u)+i)\bmod d\bigr).
  \tag{19}\label{eq:19}
  \]
  が成り立つとする。このとき、任意の無限ウォーク \((v_i,w_i)\) が \(C\) 内にあるならば、その出力 \(y=w_0w_1\cdots\) は、
  \[
  y_t=\lambda((h(v_0)+t)\bmod d)
  \]
  をすべての \(t\ge0\) に対して満たす。特に、\(y_{t+d}=y_t\) がすべての \(t\) に対して成り立つ。
\end{enumerate}
\label{thm-end:12}
\end{lemma}
\end{samepage}

% Markdown AST block 73: Para
\(C\) が（\(h\) に関して）\emph{認証される}とは、(19)がある \(\lambda\) に対して成り立つことをいう。

% Markdown AST block 74: Proof.
\begin{proof}\leavevmode\label{proof:12}
(a) 内部辺 \(e_1=(u,w,v)\) を一つ取る。強連結性により、道 \(e_2\cdots e_m\) を \(v\) から \(u\) へ \(C\) の内部で取れる（\(u=v\) の場合は空の道でよい）。閉ウォーク \(e_1\cdots e_m\) に沿って足し合わせると、\(h\) の項が望遠鏡和として相殺され、
\[
\sum_{k\le m}\delta(e_k)=\sum_{k\le m}|w_k|\ge1.
\]
となる。したがって、すべての \(\delta(e)\) がゼロとなることはなく、\(d\ge1\) である。合同式は最大公約数の定義から従う。また、空でない閉ウォークの総長は、その \(\delta\) の和に等しく、\(d\) の正の倍数である。

% Markdown AST block 75: OrderedList
\begin{enumerate}
\def\labelenumi{(\alph{enumi})}
\setcounter{enumi}{1}
\tightlist
\item
  \(\tau_i=\sum_{k<i}|w_k|\) とおく。(a)を用いてウォークに沿って帰納法を行うと、
  \[
  h(v_i)\equiv h(v_0)+\tau_i\pmod d.
  \]
  を得る。位置 \(t=\tau_i+k\) にある出力文字は、\(0\le k<|w_i|\) のとき、
  \[
  (w_i)_k=\lambda((h(v_i)+k)\bmod d)=\lambda((h(v_0)+t)\bmod d).\qedhere
\]
  である。
\end{enumerate}
\end{proof}

% Markdown AST block 76: Para
補題12は、任意の整数値関数 \(h\) に対する十分条件である。以下のアルゴリズムでは、根 \(\rho\in C\) を選び、各 \(v\in C\) に対して道 \(\pi_v\) を \(\rho\) から \(v\) へ \(C\) 内で取り、
\[
h(v)=|\operatorname{out}(\pi_v)|.
\]
とおく。したがって、\(h(v)\) は辺の本数ではなく総語長である。\(\rho\) への道は空でもよく、それ以外の道も最短である必要はない。実装では、グラフ探索で最初に見つかった根からの道を用いる。

% Markdown AST block 77: Para
任意の \(h\) を明示的に指定する場合を除き、以下で\emph{認証される}または\emph{認証されない}というときは、これらの根からの道の長さを用いる完全な判定(19)を指す。この規約のもとで、判定は内部のすべての無限出力の最終的周期性について完全であり、その結果は根と道の選択に依存しない（定理21および系22）。資源上限によって判定が未完に終わった場合は、判定未了である。

% Markdown AST block 78: Para
補題12(a)により、各位相 \(s\in\mathbb Z/d\) には、ある内部辺の文字が少なくとも一つ置かれる。総長 \(\ell\ge d\) の閉ウォークが、位相 \(h(v_0),\ldots,h(v_0)+\ell-1\) を占めるからである。すべての語の長さが1の場合、(19)は第3.2節の位相判定条件となる。

% Markdown AST block 79: Para
\textbf{除去。} \(\mathcal U\) を、\(G\) の認証されない巡回的SCCの族とする。次のように定義する。
\[
R(G)=\Bigl(\bigcup_{C\in\mathcal U}C,\
\{(u,w,v)\in E:\exists C\in\mathcal U,\ u,v\in C\}\Bigr):
\]
認証されない成分の内部辺だけを保持する。成分間の辺、非巡回的成分、および認証された成分は除去する。

% Markdown AST block 80: Proposition 13 (tail preservation under removal).
\begin{samepage}
\begin{proposition}[除去による尾部保存]\leavevmode\label{thm:13}
\(G\models n_0\) ならば、\(R(G)\models n_1\) がある \(n_1\ge n_0\) に対して成り立つ。
\label{thm-end:13}
\end{proposition}
\end{samepage}

% Markdown AST block 81: Proof.
\begin{proof}\leavevmode\label{proof:13}
被覆を与えるウォークは有限グラフ内の無限ウォークなので、補題6により、\(i_1\) と巡回的SCC \(C\) が存在し、\(v_i\in C\) がすべての \(i\ge i_1\) に対して成り立つ。辺 \((v_i,w_i,v_{i+1})\)、\(i\ge i_1\) は \(C\) の内部にある。もし \(C\) が認証されていたならば、補題12(b)により、
\[
c_{n+d}=c_n
\]
がすべての \(n\ge t_{i_1}\) に対して、\(d\ge1\) として成り立ち、補題4に矛盾する。したがって、\(C\in\mathcal U\) であり、\((v_i,w_i)_{i\ge i_1}\) は \(R(G)\) のウォークである。\(n_1=t_{i_1}\) とすればよい。\qedhere
\end{proof}

% Markdown AST block 82: Para
内部辺だけを保持することは健全である。補題6により、尾部のすべての辺は一つの成分の内部にあるからである。認証されない成分の内部辺を除去することは健全ではない。より多くの辺を保持しても健全性は保たれるが、報告する件数は変わる。

% BEGIN PUBLIC ADDITION figure:phase-test
位相判定は、巡回する位置に割り当てた文字の整合性を調べる操作として読める。

\begin{figure}[htbp]
\centering
\begin{tikzpicture}[paperfig]
\foreach \x/\s in {0/L,7/R}{
 \begin{scope}[xshift=\x cm]
 \node[fnode] (\s u) at (0,0) {$u$};
 \node[fnode] (\s v) at (3.3,0) {$v$};
 \node[below=3pt of \s u] {$h(u)=0$};
 \node[below=3pt of \s v] {$h(v)=1$};
 \draw[fedge] (\s u) to[bend left=30] node[flabel,above] {$2$} (\s v);
 \draw[fedge] (\s v) to[bend left=30] node[flabel,below] {$4$} (\s u);
 \node at (1.65,-1.35) {$\delta(u\to v)=0,\quad\delta(v\to u)=2$};
 \node at (1.65,-1.85) {$d=2$};
 \end{scope}
}
\node[ftitle] at(1.65,1.3){\figtext{Certified cycle}{認証される閉路}};
\node[ftitle] at(8.65,1.3){\figtext{One extra edge}{辺を一本追加}};
\draw[fedge,dashed] (Ru) to node[flabel,above] {$4$} (Rv);
\node[fbox,text width=5.2cm] at (1.65,-2.85)
 {$\lambda(0)=2,\quad\lambda(1)=4$\\[3pt]
 \figtext{Output: $242424\ldots$ (up to shift)}{出力：$242424\ldots$（開始位相を除く）}};
\node[fbox,text width=5.2cm] at (8.65,-2.85)
 {$\lambda(0)=2\ \text{\figtext{and}{かつ}}\ \lambda(0)=4$\\[3pt]
 \figtext{Conflict at phase $0$; not certified}{位相 $0$ で不整合：認証されない}};
\node[align=center,text width=12cm] at (5.15,-4)
 {\figtext{In both panels, $h$ is the total word length of a chosen path from the root $u$.}{両図とも、$h$ は根 $u$ から選んだ経路の総語長である。}};
\end{tikzpicture}

\caption{一文字のラベルを持つ二つの抽象的な語グラフ。両図とも、\(h(u)=0,h(v)=1\)は根\(u\)からの経路の総語長であり、\(d=2\)となる。左の閉路は認証される。右では破線の辺を追加すると位相ゼロに異なる二文字が置かれ、完全な判定は失敗する。これらは判定法の例であり、\(5/2\)の特定のキャリー集合の例ではない。}
\label{fig:phase-test}
\end{figure}
% END PUBLIC ADDITION figure:phase-test

% Markdown AST block 83: 4.3 Exact contraction of deterministic paths
\Needspace{6\baselineskip}
\subsection{決定的な道の厳密な縮約}\label{sec:4.3}

% Markdown AST block 84: Para
\(G'=R(G)\) とする。まず、同一の三つ組を併合する。これはウォークの集合を変えない。出次数 \(\deg^+(u)\) は、併合後の \(G'\) で数える。

% Markdown AST block 85: Lemma 14.
\begin{samepage}
\begin{lemma}\leavevmode\label{thm:14}
(a) \(G'\) の各頂点は \(\deg^+\ge1\) を満たす。(b) 成分 \(C\) が \(G'\) の成分であり、その各頂点が \(\deg^+=1\) を満たすならば、\(C\) は単純有向閉路である（自己ループをもつ1頂点の場合を含む）。また、\(C\) は、閉路に沿った根からの道の長さに関して、補題12の意味で認証される。
\label{thm-end:14}
\end{lemma}
\end{samepage}

% Markdown AST block 86: Proof.
\begin{proof}\leavevmode\label{proof:14}
(a) \(|C|\ge2\) の場合には、強連結性により、各頂点はそこから出る内部辺をもつ。\(|C|=1\) の場合には、巡回性により自己ループをもつ。

% Markdown AST block 87: Para
\textbf{(b)} 後続頂点が一意であるとき、内部辺は関数グラフをなす。強連結な関数グラフは一つの閉路
\[
v_1\xrightarrow{w_1}\cdots\xrightarrow{w_m}v_1.
\]
である。\(L=\sum|w_k|\)、\(h(v_1)=0\) とおき、
\[
h(v_{k+1})=h(v_k)+|w_k|.
\]
とする。このとき \(\delta(e_k)=0\) が \(k<m\) に対して成り立ち、\(\delta(e_m)=L\) および \(d=L\) である。各位相 \(0\le s<L\) には、ちょうど一つの文字、すなわち \(s\) 番目に当たる \(w_1\cdots w_m\) の文字が置かれるので、(19)が成り立つ。\qedhere
\end{proof}

% Markdown AST block 88: Para
したがって、除去において根からの道の長さによる完全な判定を用いる場合、出次数1の頂点だけからなる成分が縮約段階に到達することはない。それでも実装では、安全策として、そのような成分に一つの\emph{アンカー}頂点を残す（資源上限によって位相判定を省略した場合に用いる）。第5節の計算では、この規則が働くことは一度もなかった。

% Markdown AST block 89: Definition (macro vertices and contraction).
\begin{definition}[マクロ頂点と縮約]
\(B\) を、\(G'\) の頂点で \(\deg^+\ge2\) を満たすすべてのものと、そのような頂点を含まない各成分のアンカー一つ（頂点順序で最小の頂点）からなる集合とする。\(u\in B\) と、各辺 \(e=(u,w,v)\) で \(u\) から出るもの に対して、次を行う。\(v\notin B\) である間、一意な辺 \((v,w',v')\) で \(v\) から出るもの をたどり、\(w\leftarrow ww'\)、\(v\leftarrow v'\) と置き換える。初めて \(v\in B\) となったとき、\emph{マクロ辺} \((u,w,v)\) を記録する。縮約 \(P(G')\) の頂点集合は \(B\) であり、辺集合はマクロ辺の集合である（同一の三つ組は併合する）。
\end{definition}

% Markdown AST block 90: Lemma 15.
\begin{samepage}
\begin{lemma}\leavevmode\label{thm:15}
定義中の反復は停止する。また、\(G'\) の任意の無限ウォークは \(B\) を無限回訪れる。
\label{thm-end:15}
\end{lemma}
\end{samepage}

% Markdown AST block 91: Proof.
\begin{proof}\leavevmode\label{proof:15}
反復が停止しないとすると、訪れる頂点はすべて出次数1であり、頂点数は有限なので、反復は出次数1の頂点からなる閉路 \(Z\) に入る。\(Z\) の各頂点から出る一意な辺は \(Z\) 内にとどまるため、\(Z\) から、それを含む成分 \(C\) 内の補集合へ向かう辺は存在しない。強連結性から \(C=Z\) が従う。これは出次数1の頂点だけからなる成分であり、そのアンカーは \(B\cap Z\) に属するので、反復はそこで停止する。これは矛盾である。ある時点以降に \(B\) を避ける無限ウォークにも、同じ議論を適用できる。\qedhere
\end{proof}

% Markdown AST block 92: Proposition 16 (exactness of contraction).
\begin{samepage}
\begin{proposition}[縮約の厳密性]\leavevmode\label{thm:16}
\(G'\) の任意の無限ウォークで、\(B\) の頂点から始まるものは、\(P(G')\) の無限ウォークを定める。後者の頂点列は、前者が \(B\) を訪れる頂点の列であり、出力は同じ無限語であり、その \(k\) 番目のマクロ辺の長さは、対応する区間の総長に等しい。逆に、\(P(G')\) の任意の無限ウォークは、\(G'\) の少なくとも一つの無限ウォークからこのようにして得られる。したがって、\(G'\models n_1\) ならば、
\[
P(G')\models n_2,
\]
が成り立つ。ここで \(n_2\) は、\(n_1\) 以降に被覆を与えるウォークが \(B\) を最初に訪れる時刻である。
\label{thm-end:16}
\end{proposition}
\end{samepage}

% Markdown AST block 93: Proof.
\begin{proof}\leavevmode\label{proof:16}
頂点 \(u\in B\) でウォークは辺 \(e=(u,w,v)\) を選ぶ。\(v\notin B\) である間は一意な出る辺をたどることが強制され、定義中の反復と同じ頂点と語を通り、有限ステップの後に \(B\) の次の頂点へ到達する（補題15）。

% Markdown AST block 94: Para
したがって、\(B\) への連続する訪問の間の各区間は、ちょうど一つのマクロ辺に対応する。逆に、各マクロ辺は \(G'\) の少なくとも一つの道に展開できる。同一の三つ組を併合した後では、始点、語、終点が同じ二つの区間は同じマクロ辺を与えるため、この対応は全射であるが、単射とは限らない。

% Markdown AST block 95: Para
この対応は区間ごとに成り立ち、出力はいずれも同じ語の連結である。尾部保存に用いるのは順方向だけである。時刻 \(n_2\) の存在は補題15から従い、被覆に関する主張が得られる。\qedhere
\end{proof}

% Markdown AST block 96: Para
ここでは、縮約を除去と持ち上げの間だけで用いる。セル数 \(K\) は語を用いる手法の初期段階で固定し、セル座標が関与するのは初期辺を通じてのみである。縮約後のセルの細分は用いず、そのような細分が可能であるとも主張しない。

% BEGIN PUBLIC ADDITION figure:word-contraction
縮約は出力語の区切りを変えるが、すべての文字とその順序を保存する。

\begin{figure}[htbp]
\centering
\begin{tikzpicture}[paperfig]
\node[ftitle,anchor=west] at(-.4,1.7){\figtext{Before contraction}{縮約前}};
\node[fnode] (a) at(0,0){$u$};\node[fnode] (x) at(4,0){$x$};\node[fnode] (b) at(8,0){$v$};
\draw[fedge,femph] (a)--node[flabel,above] {$(-1)$}(x);
\draw[fedge,femph] (x)--node[flabel,above] {$(1,3)$}(b);
\draw[fedge,dashed] (a) to[bend right=23] node[flabel,below] {$(3)$}(b);
\draw[fedge] (b) to[bend right=23] node[flabel,above] {$(1)$}(a);
\draw[fedge] (b) edge[loop right] node[right] {$(-1)$}(b);
\node[below=4pt] at(x.south){$\deg^+(x)=1$};
\node[align=center] at(10.6,-.1){$u,v\in B$\\$x\notin B$};
\draw[fedge,line width=1pt] (4,-1.45)--(4,-2.1);
\node[ftitle,anchor=west] at(-.4,-2.5){\figtext{After contraction}{縮約後}};
\node[fnode] (aa) at(0,-3.9){$u$};\node[fnode] (bb) at(8,-3.9){$v$};
\draw[fedge,femph] (aa)--node[flabel,above] {$(-1,1,3)$}(bb);
\draw[fedge,dashed] (aa) to[bend right=23] node[flabel,below] {$(3)$}(bb);
\draw[fedge] (bb) to[bend right=23] node[flabel,above] {$(1)$}(aa);
\draw[fedge] (bb) edge[loop right] node[right] {$(-1)$}(bb);
\node[align=center] at(10.6,-3.9){$\ell=3$};
\node[align=center,text width=12cm] at(5,-5.6)
 {\figtext{Concatenate every letter, in order; keep the total elapsed length.}{すべての文字を順に連結し、経過した総語長を保持する。}};
\end{tikzpicture}

\caption{正確な縮約の例。頂点\(x\)の出辺は一本であり、\(u\)と\(v\)はそれぞれ二本の出辺を持つためマクロ頂点である。語\((-1)\)と\((1,3)\)は\((-1,1,3)\)に連結され、その語長は三のまま保たれる。他の辺はすべて保持される。線種で対応する辺を示す。}
\label{fig:word-contraction}
\end{figure}
% END PUBLIC ADDITION figure:word-contraction

% Markdown AST block 97: 4.4 Lifting by a new prime
\Needspace{6\baselineskip}
\subsection{新しい素数による持ち上げ}\label{sec:4.4}

% Markdown AST block 98: Para
\(p\) を \(p\nmid b\) を満たす素数とする。\emph{持ち上げ} \(\operatorname{Lift}_p(G)\) の頂点は \((u,q)\) であり、\(u\in V(G)\) および \(q\in\mathbb F_p^\times\) を満たす。辺 \((u,w,v)\in E(G)\) について、\(w=(c_0,\ldots,c_{\ell-1})\) および \(q_0=q\) として、
\[
q_{i+1}=b^{-1}(aq_i+c_i)\bmod p\qquad(0\le i<\ell).
\]
を計算する。辺 \(((u,q_0),w,(v,q_\ell))\) が存在することと、\(q_0,q_1,\ldots,q_\ell\) がすべて非ゼロであることは同値である。終点の剰余 \(q_\ell\) に関する条件も含める。\(q_\ell\) は次の辺の開始剰余だからである。どの辺にも接続しない頂点は削除してよい。出る辺をもたない頂点は無限ウォーク上に現れず、入る辺をもたない頂点は無限ウォークの始点にしか現れないので、開始点を1辺先へずらせば、これらの削除後も被覆が保存される。

% Markdown AST block 99: Proposition 17 (tail preservation under lifting).
\begin{samepage}
\begin{proposition}[持ち上げによる尾部保存]\leavevmode\label{thm:17}
\(p\nmid b\)、\(G\models n_0\) とし、\(p\nmid Q_n\) がすべての \(n\ge n'\) に対して成り立つと仮定する。このとき、
\[
\operatorname{Lift}_p(G)\models n_1
\]
が \(n_1=t_{i_0}\) に対して成り立つ。ここで、\(i_0\) は \(t_{i_0}\ge\max(n_0,n')\) を満たす最小の添字である。
\label{thm-end:17}
\end{proposition}
\end{samepage}

% Markdown AST block 100: Proof.
\begin{proof}\leavevmode\label{proof:17}
\(i\ge i_0\) に対して、
\[
q^{(i)}=Q_{t_i}\bmod p\ne0.
\]
とおく。\(bQ_{n+1}=aQ_n+c_n\) を法 \(p\) で還元すると、
\[
Q_{n+1}\equiv b^{-1}(aQ_n+c_n).
\]
を得る。

% Markdown AST block 101: Para
したがって、
\(w_i=(c_{t_i},\ldots,c_{t_{i+1}-1})\)
に沿い、\(q^{(i)}\) から始まる前向きの反復は、
\[
Q_{t_i},Q_{t_i+1},\ldots,Q_{t_{i+1}}
\]
 の剰余を生成する。これらはすべて非ゼロであり、最後は \(q^{(i+1)}\) となる。

% Markdown AST block 102: Para
したがって、
\[
((v_i,q^{(i)}),w_i,(v_{i+1},q^{(i+1)}))
\]
は持ち上げの辺であり、\((v_i,q^{(i)})_{i\ge i_0}\) は被覆を与えるウォークである。\qedhere
\end{proof}

% Markdown AST block 103: Para
開始剰余 \(q\in\mathbb F_p^\times\) をすべて列挙することにより、この議論は初期整数の大きさに依存しなくなる。単元 \(Q_{t_i}\bmod p\) がどの値であっても、対応する頂点が存在するからである。

% Markdown AST block 104: Lemma 18 (root formula).
\begin{samepage}
\begin{lemma}[根の公式]\leavevmode\label{thm:18}
\(p\nmid ab\) とし、
\[
C_0=0,\qquad C_{i+1}=aC_i+b^ic_i
\]
を語 \(w=(c_0,\ldots,c_{\ell-1})\) に対して定める。このとき(8)が成り立つ。さらに、
\[
R_i=-C_ia^{-i}\bmod p,
\]
とおく。このとき \(R_0=0\) であり、
\[
R_{i+1}=R_i-c_ib^ia^{-i-1}.
\]
が成り立つ。開始剰余 \(q\) に対して、前向きの反復で \(q_0,\ldots,q_\ell\) がすべて非ゼロとなるための必要十分条件は、
\[
q\notin\{R_0,\ldots,R_\ell\},
\]
である。その場合、
\[
\begin{aligned}
q_\ell&=\alpha q+\beta\qquad \text{with}\qquad \alpha=(ab^{-1})^\ell,\\
\beta_0&=0,\qquad \beta_{i+1}=ab^{-1}\beta_i+b^{-1}c_i.
\end{aligned}
\]
が成り立つ。
\label{thm-end:18}
\end{lemma}
\end{samepage}

% Markdown AST block 105: Proof.
\begin{proof}\leavevmode\label{proof:18}
\(q\) の任意の整数への持ち上げに(8)を適用すると、
\[
b^iq_i\equiv a^iq+C_i.
\]
となる。\(p\nmid b\) なので、\(q_i\equiv0\) と \(a^iq+C_i\equiv0\) は同値であり、これは \(q\equiv-C_ia^{-i}=R_i\)（\(p\nmid a\)）ということである。

% Markdown AST block 106: Para
\(R_i\) の漸化式は
\[
-(aC_i+b^ic_i)a^{-i-1}=R_i-c_ib^ia^{-i-1},
\]
であり、\(q_\ell\) のアフィン表示は、
\[
q_{i+1}=ab^{-1}q_i+b^{-1}c_i.\qedhere
\]
を展開したものである。
\end{proof}

% Markdown AST block 107: Para
したがって、禁止される開始剰余による列挙と前向きの反復は、同じ辺集合を与える。実装Aは前者を、実装Bは後者を用いる（第9節）。\(p\mid b\) の場合、\(b\) の逆元は存在せず、
\[
0\equiv aQ_n+c_n\pmod p,
\]
となるので、\(Q_n\bmod p\) は \(c_n\) のみで定まり、条件 \(p\nmid Q_n\) は文字集合に関する条件 \(c_n\not\equiv0\pmod p\) と同値になる。\(p\mid a\) の場合には、
\[
Q_{n+1}\equiv b^{-1}c_n,
\]
となり、開始剰余に依存しない。文字集合のすべての文字が \(p\) と互いに素であれば（補題23と24の文字集合では、各文字が \(ab\) と互いに素である）、持ち上げは何も除去せずに頂点を \(p-1\) 倍にする。一方、語が \(p\) で割り切れる文字を含む辺は、すべての開始剰余に対して除去される。第5節では、このようにして素数 \(2\) と \(5\) を扱う。

% BEGIN PUBLIC ADDITION figure:prime-lift
素数による持ち上げで語全体を検査するには、端点だけでなく、すべての中間剰余が必要である。

\begin{figure}[htbp]
\centering
\begin{tikzpicture}[paperfig]
\node[align=center] at(5.2,1.5)
 {$a=5,\quad b=2,\quad p=7,\quad w=(-1,1)$\\[5pt]
 $q_{i+1}=4(5q_i+c_i)\pmod7$};
\node[ftitle,anchor=west] at(-1.3,0){\figtext{Keep}{保持}};
\node[fnode] (a) at(1,0){$1$};\node[fnode] (b) at(5,0){$2$};\node[fnode] (c) at(9,0){$2$};
\draw[fedge] (a)--node[flabel,above] {$c_0=-1$}(b);
\draw[fedge] (b)--node[flabel,above] {$c_1=1$}(c);
\node[below=5pt] at(5,-.3){$q_0,q_1,q_2\ne0$};
\node[ftitle,anchor=west] at(-1.3,-2){\figtext{Discard}{除去}};
\node[fnode] (d) at(1,-2){$3$};\node[fnode,fill=black!12,double] (e) at(5,-2){$0$};\node[fnode] (f) at(9,-2){$4$};
\draw[fedge,dashed] (d)--node[flabel,above] {$c_0=-1$}(e);
\draw[fedge,dashed] (e)--node[flabel,above] {$c_1=1$}(f);
\node[below=5pt] at(5,-2.3){$q_0,q_2\ne0,\qquad q_1=0$};
\node[fbox,text width=11.5cm] at(5,-3.85)
 {\figtext{Endpoint checks alone would retain the second edge incorrectly.}{端点だけの検査では、二つ目の辺を誤って保持してしまう。}\\
 \figtext{Inspect all intermediate residues, including the last one.}{最後の剰余を含む、すべての中間剰余を検査する。}};
\end{tikzpicture}

\caption{\(a=5,b=2,p=7\)で同じ語\((-1,1)\)を持ち上げた二例。\(2^{-1}=4\pmod7\)である。開始剰余\(1\)からは\(1,2,2\)となり辺が残る。開始剰余\(3\)からは\(3,0,4\)となり、両端が非零でも辺は除去される。表示した剰余はすべて法七の整数演算で求まる。}
\label{fig:prime-lift}
\end{figure}
% END PUBLIC ADDITION figure:prime-lift

% Markdown AST block 108: 4.5 The finite success theorem
\Needspace{6\baselineskip}
\subsection{有限成功定理}\label{sec:4.5}

% Markdown AST block 109: Para
\textbf{仮定(H)。} 有限素数集合 \(S_0\) は
\[
M_0=\prod_{p\in S_0}p
\]
を満たし、文字集合 \(\mathcal C\) とともに次の条件を満たす。任意の \(\xi>0\) と任意の \(n\) に対して、
\[
\gcd(Q_n,M_0)=\gcd(Q_{n+1},M_0)=1
\]
ならば \(c_n\in\mathcal C\) である。

% Markdown AST block 110: Theorem 19 (word-labelled finite success).
\begin{samepage}
\begin{theorem}[語ラベル付き有限成功]\leavevmode\label{thm:19}
\(a>b\ge2\)、\(\gcd(a,b)=1\)、\(K\ge1\) とし、\(S_0,\mathcal C\) が(H)を満たすとする。\(p_1,\ldots,p_s\) を相異なる素数で、\(p_k\nmid ab\) および \(p_k\notin S_0\) を満たすものとする。\(G_0^{\rm in}\) を、頂点 \(\{0,\ldots,K-1\}\) と辺 \((j,(c),k)\) からなるグラフとする。辺は、すべての \(c\in\mathcal C\) と \(0\le j,k<K\) のうち
\[
aj-b(k+1)<cK<a(j+1)-bk.
\]
を満たすものに対して与える。次のように定義する。
\[
\begin{aligned}
G_k^{\rm out}&=P(R(G_k^{\rm in}))\ \ (0\le k\le s),\\
G_k^{\rm in}&=\operatorname{Lift}_{p_k}(G_{k-1}^{\rm out})\ \ (1\le k\le s).
\end{aligned}
\]
\(G_s^{\rm out}=\varnothing\) ならば、\(M=M_0\prod_{k=1}^sp_k\) に対して、
\[
\forall\xi>0\ \forall N\in\mathbb N\ \exists n\ge\max(N,1):\
\gcd(Q_n,M)>1.
\tag{20}\label{eq:20}
\]
が成り立つ。
\label{thm-end:19}
\end{theorem}
\end{samepage}

% Markdown AST block 111: Proof.
\begin{proof}\leavevmode\label{proof:19}
\(\xi>0\) を固定し、(20)が成り立たないと仮定する。すなわち、\(n_0\ge1\) が存在し、
\[
\gcd(Q_n,M)=1
\]
がすべての \(n\ge n_0\) に対して成り立つ。(H)により、\(c_n\in\mathcal C\) がすべての \(n\ge n_0\) に対して成り立つ。

% Markdown AST block 112: Para
セル
\(j_n=\lfloor K\theta_n\rfloor\)
は \(\{0,\ldots,K-1\}\) に属する。また、命題5を \(M=1\) に対して適用すると（この場合、(E2)は自明である）、
\[
j_n\xrightarrow{c_n}j_{n+1}
\]
は \(G_0^{\rm in}\) の辺であるとわかる。したがって、
\[
G_0^{\rm in}\models n_0
\]
が成り立ち、すべての語の長さは1である。

% Markdown AST block 113: Para
命題13と命題16から、
\[
G_0^{\rm out}\models n_1.
\]
を得る。

% Markdown AST block 114: Para
\(k\ge1\) に対して、\(p_k\mid M\) より \(p_k\nmid Q_n\) が \(n\ge n_0\) に対して成り立つので、命題17から
\[
G_k^{\rm in}\models n'_k,
\]
を得る。さらに再び命題13と16により、
\[
G_k^{\rm out}\models n_{k+1}.
\]
を得る。

% Markdown AST block 115: Para
有限回の操作の後に、
\[
G_s^{\rm out}\models n_{s+1},
\]
となる。しかし、空グラフには無限ウォークが存在しない。\qedhere
\end{proof}

% Markdown AST block 116: Para
(20)の量化は(18)と同じであり、グラフは \(\xi\) を入力として受け取らない。\(S_0\) の素数では持ち上げを行わず、これらは文字集合を通じてのみ関与する。(H)は算術的な主張であり、各底について手計算で確認する必要がある（補題23および24）。

% Markdown AST block 117: Corollary 20.
\begin{samepage}
\begin{corollary}\leavevmode\label{thm:20}
定理19の仮定のもとで、
\[
P_{\max}=\max(S_0\cup\{p_1,\ldots,p_s\}),
\]
とする。任意の \(\xi>0\) と任意の \(N\) に対して、ある \(n\ge\max(N,1)\) が存在し、\(Q_n\) は合成数となる。
\label{thm-end:20}
\end{corollary}
\end{samepage}

% Markdown AST block 118: Proof.
\begin{proof}\leavevmode\label{proof:20}
(7)により、\(Q_n>P_{\max}\) が十分大きなすべての \(n\) に対して成り立つ。その閾値以降かつ \(N\) 以降で(20)を適用する。素数 \(p\) が \(\gcd(Q_n,M)\) を割るとき、
\[
p\le P_{\max}<Q_n,
\]
を満たすので、\(Q_n=pv\) と書け、\(v>1\) である。\qedhere
\end{proof}

% Markdown AST block 119: 4.6 The all-edge test and commuting closed words
\Needspace{6\baselineskip}
\subsection{全辺判定と可換な閉語}\label{sec:4.6}

% Markdown AST block 120: Theorem 21.
\begin{samepage}
\begin{theorem}\leavevmode\label{thm:21}
\(C\) を語ラベル付きグラフの巡回的SCCとし、\(\rho\in C\) とする。次の条件は同値である。

% Markdown AST block 121: OrderedList
\begin{enumerate}
\def\labelenumi{(\alph{enumi})}
\item
  整数値関数 \(h\) が \(C\) 上に存在し、\(\lambda\) とともに(19)を満たす。さらに、ある \(h\) に対してこれが成り立つならば、\(h(v)\) を、\(\rho\) から \(v\) へ \(C\) 内で選んだ任意の道の総語長としても成り立つ。
\item
  \(C\) 内の任意の無限ウォークの出力は最終的周期的である。
\item
  \(\rho\) における空でない任意の二つの閉出力語（\(\rho\) から \(\rho\) へ \(C\) 内で戻る空でない閉ウォークの出力）は可換である。
\end{enumerate}
\label{thm-end:21}
\end{theorem}
\end{samepage}

% Markdown AST block 122: Proof.
\begin{proof}\leavevmode\label{proof:21}
(a)\(\Rightarrow\)(b)は補題12(b)である。

% Markdown AST block 123: Para
(b)\(\Rightarrow\)(c)。閉語 \(u,v\) が \(\rho\) において存在し、\(uv\ne vu\) を満たすと仮定する。\(U,V\) を対応する閉ウォークとする。\(UV\) と \(VU\) は、それぞれ出力 \(x_0=uv\) と \(x_1=vu\) をもち、長さが同じ \(L\) である閉ウォークであり、\(x_0\ne x_1\) である。\(\varepsilon\in\{0,1\}^{\mathbb N}\) に対して、ウォーク \(X_{\varepsilon_0}X_{\varepsilon_1}\cdots\) は \(C\) 内にとどまり、出力は
\[
z_\varepsilon=x_{\varepsilon_0}x_{\varepsilon_1}\cdots.
\]
となる。

% Markdown AST block 124: Para
\(z_\varepsilon\) を長さ \(L\) のブロックに分けて読むと、\(\varepsilon\) が復元される。もし \(z_\varepsilon\) が最終的周期的で、周期 \(P\) をもつならば、\(PL\) も周期となり、ブロックの復号によって \(\varepsilon\) も最終的周期的となる。\(\varepsilon\) を最終的周期的でないもの（例えば \(0\,1\,00\,1\,000\,1\cdots\)）に選ぶと、(b)に矛盾する。

% Markdown AST block 125: Para
(c)\(\Rightarrow\)(a)。可換な二つの空でない語は、共通の語の冪である（長い方の語の先頭から短い方の語を取り除く、長さに関する帰納法による）。また、空でない語の原始根は一意である。したがって、\(\rho\) におけるすべての閉語は、一つの原始語 \(z\) の冪となる。\(p=|z|\) とおく。各 \(v\in C\) に対して、道 \(\pi_v\) を \(\rho\) から \(v\) へ固定し、
\[
h(v)=|\operatorname{out}(\pi_v)|
\]
とする。また、戻りの道 \(\sigma_v\) を \(v\) から \(\rho\) へ固定する。\(v=\rho\) における道は空でもよい。

% Markdown AST block 126: Para
内部辺 \(e=(u,w,v)\) に対して、語
\[
\begin{aligned}
\omega_1&=\operatorname{out}(\pi_u)\,w\,\operatorname{out}(\sigma_v)\qquad \text{and}\\
\omega_2&=\operatorname{out}(\pi_v)\operatorname{out}(\sigma_v)
\end{aligned}
\]
はいずれも閉語であり、\(\omega_1\ne\varepsilon\) である。どちらも \(z\) の冪である。\(\omega_2\) が空の場合には0乗も許す。したがって、\(p\) は
\[
|\omega_1|-|\omega_2|=h(u)+|w|-h(v);
\]
を割る。よって、\(p\mid\delta(e)\) がすべての内部辺について成り立ち、\(p\mid d\) である。次のように定義する。
\[
\lambda(s)=z_{s\bmod p}
\]
これは \(s\in\mathbb Z/d\) に対する定義であり、\(p\mid d\) より矛盾なく定まる。

% Markdown AST block 127: Para
\(\omega_1=z^k\) であるから、その位置 \(t\) にある文字は \(z_{t\bmod p}\) である。また、\(w_i\) は、位置 \(h(u)+i\) にある \(\omega_1\) の文字であるので、
\[
w_i=z_{(h(u)+i)\bmod p}=\lambda((h(u)+i)\bmod d).
\]
となる。これは根からの道の長さ \(h\) に対する(19)である。\qedhere
\end{proof}

% Markdown AST block 128: Corollary 22.
\begin{samepage}
\begin{corollary}\leavevmode\label{thm:22}
根 \(\rho\in C\) と道 \(\pi_v:\rho\to v\) を \(C\) 内で選び、
\[
h(v)=|\operatorname{out}(\pi_v)|.
\]
とおく。(19)が成り立つかどうかは、根およびこれらの道の選択に依存しない。このようなどの選択に対しても、
\[
d=g_C:=\gcd\bigl\{|\operatorname{out}(\gamma)|:
\gamma\text{ is a nonempty closed walk in }C\bigr\}.
\]
が成り立つ。特に、根からの道の長さを用いる二つの完全な判定は、どの成分が認証されるか、および \(d\) の値について一致する。
\label{thm-end:22}
\end{corollary}
\end{samepage}

% Markdown AST block 129: Proof.
\begin{proof}\leavevmode\label{proof:22}
定理21により、このようなどの選択に対しても、判定が成立することと、内部のすべての無限ウォークが最終的周期的な出力をもつことは同値である。この性質は根や道の選択に関係しない。

% Markdown AST block 130: Para
最大公約数の等しさを示す。閉ウォークに沿って \(\delta(e)\) を足し合わせると、その総語長を得る。したがって、\(d\mid g_C\) である。逆に、\(e=(u,w,v)\) に対して戻りの道 \(\sigma_v:v\to\rho\) を選ぶと、
\[
\delta(e)
=|\operatorname{out}(\pi_u e\sigma_v)|
 -|\operatorname{out}(\pi_v\sigma_v)|.
\]
となる。右辺の両項はいずれも \(\rho\) における閉ウォークの長さであり、それぞれ \(g_C\) で割り切れる。第2のウォークが空の場合の長さゼロもこれに含まれる。したがって、\(g_C\mid\delta(e)\) がすべての内部辺に対して成り立つので、\(g_C\mid d\) である。これにより、その成分が認証されるかどうかにかかわらず、\(d=g_C\) が示された。\qedhere
\end{proof}

% Markdown AST block 131: Para
\(h\) に対する制限は不可欠である。次の2辺からなる閉路を考える。
\[
u\xrightarrow{(2)}v,\qquad v\xrightarrow{(4)}u.
\]
その無限出力は \((2,4)^\infty\) の二つの巡回シフトである。根からの道の長さ \(h(u)=0\)、\(h(v)=1\) を用いると、
\[
(\delta(u\to v),\delta(v\to u))=(0,2),\qquad d=2,
\]
となり、\(\lambda(0)=2\)、\(\lambda(1)=4\) として(19)が成り立つ。一方、任意の整数値関数として \(h(u)=h(v)=0\) を選ぶと、
\[
(\delta(u\to v),\delta(v\to u))=(1,1),\qquad d=1.
\]
となる。この場合、両方の文字が位相ゼロに置かれるため、(19)は成り立たない。一般の整数値関数 \(h\) に対しては、望遠鏡和の議論から得られるのは \(d\mid g_C\) だけである。補題12は依然として健全性を保証するが、その判定に失敗しても、非周期的な出力が存在するとは限らない。

% Markdown AST block 132: Para
根からの道の長さを用いる場合には、定理21から逆向きも得られる。完全な判定(19)に失敗したならば、その成分には、出力が最終的周期的でない内部無限ウォークが存在する。すべての無限出力が最終的周期的であるという理由だけで成分を除去する健全な判定では、このような成分を同じ有限グラフから除去することはできない。さらに除去を進めるには、算術または実現可能性について追加の情報が必要となる。

% Markdown AST block 133: 5. The two rational bases
\Needspace{6\baselineskip}
\section{二つの有理数の底}\label{sec:5}

% Markdown AST block 134: 5.1 The ratio 7/5 by the one-step method
\subsection{1ステップ手法による比7/5}\label{sec:5.1}

% Markdown AST block 135: Para
次のパラメータを用いる。
\[
(a,b,M,K_0,f)=(7,5,4290,32,4).
\tag{21}\label{eq:21}
\]
初期グラフの頂点数は \(\varphi(4290)K_0=960\cdot32=30720\) である。第9.2節で説明する1ステップの二つの実装は、表1の4段階を生成する。\(K=2048\) において保持集合は空になる。

% Markdown AST block 136: Table 1.
% Caption is rendered at the start of the immediately following table.

% Markdown AST block 137: Table 1
\begingroup
\small
\setlength{\tabcolsep}{4pt}
\renewcommand{\arraystretch}{1.15}
\setlength{\LTcapwidth}{\textwidth}
\begin{longtable}{@{} r r r r r r@{}}
\caption{(21)の全段階。辺はラベルを含めて数える。「認証済み」の列は認証された巡回的SCCの数、最後の列は \(|U_t|\) である。}\label{tab:1}\\
\toprule
\(K_t\) & 頂点数 & 辺数 & \shortstack[r]{巡回的\\SCC} & \shortstack[r]{認証済み\\SCC} & \shortstack[r]{保持\\頂点} \\
\midrule
\endfirsthead
\toprule
\(K_t\) & 頂点数 & 辺数 & \shortstack[r]{巡回的\\SCC} & \shortstack[r]{認証済み\\SCC} & \shortstack[r]{保持\\頂点} \\
\midrule
\endhead
\bottomrule
\endlastfoot
32 & 30720 & 55440 & 17 & 16 & 1363 \\*
128 & 5452 & 9512 & 3 & 2 & 952 \\*
512 & 3808 & 6281 & 1 & 0 & 944 \\*
2048 & 3776 & 5951 & 15 & 15 & 0 \\
\end{longtable}
\endgroup

% Markdown AST block 138: Para
非巡回的な1頂点成分を含む全SCC数は、順に29312、4440、2865、3658である。最初の3段階では、認証されない巡回的SCCがそれぞれ一つあり、その大きさは順に1363、952、944である。内部辺数は2332、1550、1505であり、三つとも \(d=1\) を満たす。位相判定ではこれらの除去は正当化されないため、次の細分へ引き継ぐ。特に、次段階の頂点数は、これらの保持集合の大きさのちょうど4倍である。

% Markdown AST block 139: Para
表2では、認証された成分を \(K,d\) と正規化した原始語によってまとめる。成分の大きさは重複も含めて列挙している。\(K=512\) では認証された成分は存在しない。

% Markdown AST block 140: Table 2.
% Caption is rendered at the start of the immediately following table.

% Markdown AST block 141: Table 2
\begingroup
\small
\setlength{\tabcolsep}{4pt}
\renewcommand{\arraystretch}{1.15}
\setlength{\LTcapwidth}{\textwidth}
\begin{longtable}{@{}>{\raggedright\arraybackslash}p{\dimexpr0.055\dimexpr\linewidth-40pt\relax} >{\raggedright\arraybackslash}p{\dimexpr0.1\dimexpr\linewidth-40pt\relax} >{\raggedright\arraybackslash}p{\dimexpr0.19\dimexpr\linewidth-40pt\relax} >{\raggedright\arraybackslash}p{\dimexpr0.035\dimexpr\linewidth-40pt\relax} >{\raggedright\arraybackslash}p{\dimexpr0.035\dimexpr\linewidth-40pt\relax} >{\raggedright\arraybackslash}p{\dimexpr0.485\dimexpr\linewidth-40pt\relax}@{}}
\caption{認証された成分の全9群。パラメータ \(d\) は(14)であり、\(p\) は表示した原始語の長さである。}\label{tab:2}\\
\toprule
\(K\) & SCC数 & 成分の大きさ & \(d\) & \(p\) & 原始語（巡回シフトを同一視） \\
\midrule
\endfirsthead
\toprule
\(K\) & SCC数 & 成分の大きさ & \(d\) & \(p\) & 原始語（巡回シフトを同一視） \\
\midrule
\endhead
\bottomrule
\endlastfoot
32 & 2 & 1,\allowbreak{} 1 & 1 & 1 & \((2)\) \\
32 & 6 & 2,\allowbreak{} 2,\allowbreak{} 2,\allowbreak{} 2,\allowbreak{} 2,\allowbreak{} 2 & 2 & 2 & \((-2,\allowbreak{}4)\) \\
32 & 8 & 6,\allowbreak{} 6,\allowbreak{} 6,\allowbreak{} 6,\allowbreak{} 6,\allowbreak{} 6,\allowbreak{} 6,\allowbreak{} 6 & 6 & 2 & \((-2,\allowbreak{}4)\) \\
128 & 1 & 24 & 8 & 8 & \((-4,\allowbreak{}2,\allowbreak{}2,\allowbreak{}2,\allowbreak{}2,\allowbreak{}2,\allowbreak{}2,\allowbreak{}2)\) \\
128 & 1 & 39 & 13 & 13 & \((-4,\allowbreak{}2,\allowbreak{}2,\allowbreak{}2,\allowbreak{}2,\allowbreak{}2,\allowbreak{}-4,\allowbreak{}2,\allowbreak{}2,\allowbreak{}2,\allowbreak{}2,\allowbreak{}2,\allowbreak{}2)\) \\
2048 & 3 & 1,\allowbreak{} 1,\allowbreak{} 1 & 1 & 1 & \((2)\) \\
2048 & 3 & 4,\allowbreak{} 4,\allowbreak{} 8 & 4 & 4 & \((-2,\allowbreak{}-2,\allowbreak{}4,\allowbreak{}4)\) \\
2048 & 3 & 6,\allowbreak{} 6,\allowbreak{} 6 & 6 & 1 & \((2)\) \\
2048 & 6 & 12,\allowbreak{} 12,\allowbreak{} 12,\allowbreak{} 12,\allowbreak{} 24,\allowbreak{} 24 & 12 & 4 & \((-2,\allowbreak{}-2,\allowbreak{}4,\allowbreak{}4)\) \\
\end{longtable}
\endgroup

% Markdown AST block 142: Para
したがって、相異なる原始語が五つ現れる。例えば、\(d=6,p=1\) と \(d=12,p=4\) の行は、\(d\) を原始語の長さとして報告してはならない理由を示している。これらの語は外側グラフの成分内で可能な出力を表す。補題4により、真の無限軌道がこれらのいずれかの成分内にとどまることはできない。

% Markdown AST block 143: Para
手続きの追加確認として、表3では、同じ構成を \(K_0=32\) において \hyperref[ref-dn]{DN, Theorem 1} の三つの床関数に関する結果に適用する。いずれも最初の段階で停止する。これは既述の整除性の結論の再現であり、それら三つの底に関する新しい主張ではない。

% Markdown AST block 144: Table 3.
% Caption is rendered at the start of the immediately following table.

% Markdown AST block 145: Table 3
\begingroup
\small
\setlength{\tabcolsep}{4pt}
\renewcommand{\arraystretch}{1.15}
\setlength{\LTcapwidth}{\textwidth}
\begin{longtable}{@{} r r r r r r r@{}}
\caption{1ステップの二つの実装によって再現した既知の例。保持集合はいずれも空である。}\label{tab:3}\\
\toprule
\(a/b\) & \shortstack[r]{素数\\集合} & \(M\) & 頂点数 & 辺数 & \shortstack[r]{巡回的\\SCC} & \shortstack[r]{認証済み\\SCC} \\
\midrule
\endfirsthead
\toprule
\(a/b\) & \shortstack[r]{素数\\集合} & \(M\) & 頂点数 & 辺数 & \shortstack[r]{巡回的\\SCC} & \shortstack[r]{認証済み\\SCC} \\
\midrule
\endhead
\bottomrule
\endlastfoot
\(3/2\) & \(\{2,\allowbreak{}5,\allowbreak{}7,\allowbreak{}11\}\) & 770 & 7680 & 8640 & 6 & 6 \\*
\(4/3\) & \(\{2,\allowbreak{}3,\allowbreak{}5\}\) & 30 & 256 & 480 & 3 & 3 \\*
\(5/4\) & \(\{2,\allowbreak{}3,\allowbreak{}7,\allowbreak{}11,\allowbreak{}13\}\) & 6006 & 46080 & 82170 & 18 & 18 \\
\end{longtable}
\endgroup

% Markdown AST block 146: 5.2 The ratio 7/5 by the word method
\Needspace{6\baselineskip}
\subsection{語を用いる手法による比7/5}\label{sec:5.2}

% Markdown AST block 147: Lemma 23.
\begin{samepage}
\begin{lemma}\leavevmode\label{thm:23}
\(a=7\)、\(b=5\) とする。\(Q_n\) と \(Q_{n+1}\) が奇数であり、\(5\nmid Q_n\) ならば、
\[
c_n\in\{-4,-2,2,4,6\}.
\]
である。したがって、(H)は
\[
S_0=\{2,5\}\qquad \text{and}\qquad \mathcal C=\{-4,-2,2,4,6\}.
\]
に対して成り立つ。
\label{thm-end:23}
\end{lemma}
\end{samepage}

% Markdown AST block 148: Proof.
\begin{proof}\leavevmode\label{proof:23}
(6)により、\(-4\le c_n\le6\) である。\(5\equiv7\equiv1\pmod2\) より、
\[
c_n=5Q_{n+1}-7Q_n\equiv Q_{n+1}-Q_n\pmod2,
\]
であり、これは両方の項が奇数のとき偶数となるので、
\[
c_n\in\{-4,-2,0,2,4,6\}.
\]
である。\(c_n=0\) ならば \(5Q_{n+1}=7Q_n\) であり、\(\gcd(5,7)=1\) より \(5\mid Q_n\) となる。\qedhere
\end{proof}

% Markdown AST block 149: Para
素数 \(2\) による持ち上げも命題17により可能であるが、\(\mathbb F_2^\times=\{1\}\) であり、持ち上げの条件は「すべての文字が偶数」という、文字集合にすでに含まれる条件に帰着する。素数 \(5\) は \(b\) を割るため、これによる持ち上げは行えない。補題18の後の注意により、\(5\nmid Q_n\) は \(c_n\not\equiv0\pmod5\) と同値であり、範囲 \([-4,6]\) で文字が偶数の場合、これは \(c_n\ne0\) を意味する。したがって、文字集合は \(2\) と \(5\) を避けるための必要条件を表す。定理19が要求するのはこれだけである。

% Markdown AST block 150: Para
定理19を、\(K=2048\)、補題23の文字集合、および素数 \(p_1,p_2,p_3=3,11,13\) に対して適用する。値 \(K=2048\) は自由パラメータであり、表1の最終セル数に合わせて選んだ。健全性はこの値に依存しない。表4に4段階を示す。第9.3節の二つの実装は、8個のグラフのすべての頂点とすべての語ラベル付き辺について一致し、最後の段階の出力グラフは空となる。この4段階では、複数の縮約鎖が同じマクロ辺を生成することはないので、相異なる辺の数は重複度を含めて数えた数と等しい。

% Markdown AST block 151: Table 4.
% Caption is rendered at the start of the immediately following table.

% Markdown AST block 152: Table 4
\begingroup
\small
\setlength{\tabcolsep}{2.5pt}
\renewcommand{\arraystretch}{1.15}
\setlength{\LTcapwidth}{\textwidth}
\begin{longtable}{@{} r r r r r r r r r r@{}}
\caption{\(7/5\) に対する語を用いる手法。\(K=2048\)、\(\mathcal C=\{-4,-2,2,4,6\}\)、素数 \(3,11,13\) を用いる。入力グラフは定理19の \(G_k^{\rm in}\)、出力グラフは同定理の \(G_k^{\rm out}\) である。「未認証頂点」は認証されない巡回的SCCの頂点数を表す。}\label{tab:4}\\
\toprule
\shortstack[r]{\(k\)} & \shortstack[r]{\(p_k\)} & \shortstack[r]{入力\\頂点} & \shortstack[r]{入力\\辺} & \shortstack[r]{巡回的\\SCC} & \shortstack[r]{認証済み} & \shortstack[r]{未認証\\頂点} & \shortstack[r]{出力\\頂点} & \shortstack[r]{出力\\辺} & \shortstack[r]{最大\\語長\\（出力）} \\
\midrule
\endfirsthead
\toprule
\shortstack[r]{\(k\)} & \shortstack[r]{\(p_k\)} & \shortstack[r]{入力\\頂点} & \shortstack[r]{入力\\辺} & \shortstack[r]{巡回的\\SCC} & \shortstack[r]{認証済み} & \shortstack[r]{未認証\\頂点} & \shortstack[r]{出力\\頂点} & \shortstack[r]{出力\\辺} & \shortstack[r]{最大\\語長\\（出力）} \\
\midrule
\endhead
\bottomrule
\endlastfoot
0 & --- & 2048 & 8366 & 1 & 0 & 2048 & 2048 & 8366 & 1 \\*
1 & 3 & 4096 & 9010 & 10 & 9 & 691 & 593 & 1292 & 13 \\*
2 & 11 & 5928 & 11276 & 4 & 3 & 89 & 62 & 131 & 8 \\*
3 & 13 & 744 & 1342 & 2 & 2 & 0 & 0 & 0 & --- \\
\end{longtable}
\endgroup

% Markdown AST block 153: Para
4段階の全SCC数は1、3387、5840、704である。最終的な素数集合は(21)と同じ \(\{2,3,5,11,13\}\) となる。これは新しい底に関する定理ではなく、第4節の手法による(1)の第2の証明であるとともに、二つの手法の整合性の確認でもある。

% Markdown AST block 154: 5.3 The ratio 5/2
\Needspace{6\baselineskip}
\subsection{比5/2}\label{sec:5.3}

% Markdown AST block 155: Lemma 24.
\begin{samepage}
\begin{lemma}\leavevmode\label{thm:24}
\(a=5\)、\(b=2\) とする。\(Q_n\) が奇数ならば、
\[
c_n\in\{-1,1,3\}.
\]
である。したがって、(H)は
\[
S_0=\{2\}\qquad \text{and}\qquad \mathcal C=\{-1,1,3\}.
\]
に対して成り立つ。
\label{thm-end:24}
\end{lemma}
\end{samepage}

% Markdown AST block 156: Proof.
\begin{proof}\leavevmode\label{proof:24}
(6)により、\(-1\le c_n\le4\) であり、
\[
c_n=2Q_{n+1}-5Q_n\equiv Q_n\pmod2.\qedhere
\]
\end{proof}

% Markdown AST block 157: Para
用いたのは \(Q_n\) が奇数であることだけであり、\(Q_{n+1}\) の奇数性も \(5\) を避けることも必要ない。素数 \(5\) を(2)の法に含めない理由は、独立に三つある。まず、\(5\mid a\) なので補題18は適用できない。次に、補題18の後の注意により、\(5\) による持ち上げでは
\[
Q_{n+1}\equiv2^{-1}c_n=3c_n\pmod5,
\]
となるが、これは \(c_n\in\{-1,1,3\}\) に対して非ゼロ（剰余は \(2,3,4\)）なので、奇数の繰り上がりをもつ尾部では \(5\nmid Q_{n+1}\) が自動的に成り立ち、持ち上げによって何も除去されない。最後に、\(5\) を含まない主張の方が強い。

% Markdown AST block 158: Para
\(K=4\) とすると、初期グラフ \(G_0^{\rm in}\) は四つの頂点 \(j\in\{0,1,2,3\}\) と、12本の辺 \((j,(c),k)\) をもつ。これらは \(c\in\{-1,1,3\}\) および
\[
5j-2(k+1)<4c<5(j+1)-2k,
\]
を満たし、表5に列挙されている。定理19により、(2)の素数を避けるすべての尾部はこのグラフ内のウォークを与える。これはすべてのセルを覆うものであり、初期値の標本抽出ではない。また、逆向きの実現可能性は主張していない。

% BEGIN PUBLIC ADDITION figure:initial-graph
初期の四セルグラフは、キャリーラベルごとに分ければ全体を描ける。

\begin{figure}[htbp]
\centering
\begin{tikzpicture}[paperfig]
\begin{scope}[xshift=0.0cm]
\node[ftitle] at(1.35,2.4) {$c=-1$};
\node[fnode] (n0) at(0,1.1) {$j=0$};
\node[fnode] (n1) at(2.7,1.1) {$j=1$};
\node[fnode] (n2) at(0,-1.1) {$j=2$};
\node[fnode] (n3) at(2.7,-1.1) {$j=3$};
% exact edge: 0 -1 2
\draw[fedge] (n0) -- (n2);
% exact edge: 0 -1 3
\draw[fedge] (n0) -- (n3);
\end{scope}
\begin{scope}[xshift=4.3cm]
\node[ftitle] at(1.35,2.4) {$c=1$};
\node[fnode] (n0) at(0,1.1) {$j=0$};
\node[fnode] (n1) at(2.7,1.1) {$j=1$};
\node[fnode] (n2) at(0,-1.1) {$j=2$};
\node[fnode] (n3) at(2.7,-1.1) {$j=3$};
% exact edge: 0 1 0
\draw[fedge] (n0) edge[loop above] (n0);
% exact edge: 1 1 0
\draw[fedge] (n1) -- (n0);
% exact edge: 1 1 1
\draw[fedge] (n1) edge[loop above] (n1);
% exact edge: 1 1 2
\draw[fedge] (n1) -- (n2);
% exact edge: 2 1 3
\draw[fedge] (n2) -- (n3);
\end{scope}
\begin{scope}[xshift=8.6cm]
\node[ftitle] at(1.35,2.4) {$c=3$};
\node[fnode] (n0) at(0,1.1) {$j=0$};
\node[fnode] (n1) at(2.7,1.1) {$j=1$};
\node[fnode] (n2) at(0,-1.1) {$j=2$};
\node[fnode] (n3) at(2.7,-1.1) {$j=3$};
% exact edge: 2 3 0
\draw[fedge] (n2) -- (n0);
% exact edge: 2 3 1
\draw[fedge] (n2) -- (n1);
% exact edge: 3 3 1
\draw[fedge] (n3) -- (n1);
% exact edge: 3 3 2
\draw[fedge] (n3) -- (n2);
% exact edge: 3 3 3
\draw[fedge] (n3) edge[loop below] (n3);
\end{scope}
\node[align=center,text width=12.5cm] at(5.65,-2.1)
{$5j-2(k+1)<4c<5(j+1)-2k$\\[4pt]
\figtext{The three panels together contain all 12 labelled edges.}{三つの図を合わせると、ラベル付きの12辺すべてを含む。}};
\end{tikzpicture}

\caption{\(K=4\)における\(5/2\)初期グラフの全十二辺を、キャリー\(c\)ごとの三図に分けて示す。各図の四頂点は共通であり、三図を重ねるとグラフ全体になる。辺は図の下に記した厳密な整数不等式から列挙され、表5と一致する。}
\label{fig:initial-graph}
\end{figure}
% END PUBLIC ADDITION figure:initial-graph

% Markdown AST block 159: Table 5.
% Caption is rendered at the start of the immediately following table.

% Markdown AST block 160: Table 5
\begingroup
\small
\setlength{\tabcolsep}{4pt}
\renewcommand{\arraystretch}{1.15}
\setlength{\LTcapwidth}{\textwidth}
\begin{longtable}{@{}>{\raggedright\arraybackslash}p{\dimexpr0.19\dimexpr\linewidth-8pt\relax} >{\raggedright\arraybackslash}p{\dimexpr0.81\dimexpr\linewidth-8pt\relax}@{}}
\caption{\(5/2\)、\(K=4\) に対する12本の初期辺を、組（繰り上がり、終点セル）として表示する。}\label{tab:5}\\
\toprule
始点セル \(j\) & 辺 \((c,k)\) \\
\midrule
\endfirsthead
\toprule
始点セル \(j\) & 辺 \((c,k)\) \\
\midrule
\endhead
\bottomrule
\endlastfoot
0 & \((-1,\allowbreak{}2),\allowbreak{}(-1,\allowbreak{}3),\allowbreak{}(1,\allowbreak{}0)\) \\*
1 & \((1,\allowbreak{}0),\allowbreak{}(1,\allowbreak{}1),\allowbreak{}(1,\allowbreak{}2)\) \\*
2 & \((1,\allowbreak{}3),\allowbreak{}(3,\allowbreak{}0),\allowbreak{}(3,\allowbreak{}1)\) \\*
3 & \((3,\allowbreak{}1),\allowbreak{}(3,\allowbreak{}2),\allowbreak{}(3,\allowbreak{}3)\) \\
\end{longtable}
\endgroup

% Markdown AST block 161: Para
定理19を、\(K=4\)、補題24の文字集合、および九つの素数
\[
p_1,\ldots,p_9=3,7,11,13,17,19,23,29,31.
\]
に対して適用する。表6に10段階を示す。第9.3節の二つの実装は、20個のグラフのすべての頂点とすべての語ラベル付き辺について一致し、段階 \(9\) の出力グラフは空となる。すなわち、\(G_9^{\rm in}\) の1084個の巡回的SCCはすべて認証される。辺は三つ組の集合として数える（第4.1節）。縮約鎖が同じ三つ組を生成する重複度は証明内容に関与せず、提供された研究ノートとの比較のために付録Aへ記録する。

% Markdown AST block 162: Table 6.
% Caption is rendered at the start of the immediately following table.

% Markdown AST block 163: Table 6
\begingroup
\small
\setlength{\tabcolsep}{2.5pt}
\renewcommand{\arraystretch}{1.15}
\setlength{\LTcapwidth}{\textwidth}
\begin{longtable}{@{} r r r r r r r r r r@{}}
\caption{\(5/2\) に対する語を用いる手法。\(K=4\)、\(\mathcal C=\{-1,1,3\}\) を用いる。各列は表4と同じである。}\label{tab:6}\\
\toprule
\shortstack[r]{\(k\)} & \shortstack[r]{\(p_k\)} & \shortstack[r]{入力\\頂点} & \shortstack[r]{入力\\辺} & \shortstack[r]{巡回的\\SCC} & \shortstack[r]{認証済み} & \shortstack[r]{未認証\\頂点} & \shortstack[r]{出力\\頂点} & \shortstack[r]{出力\\辺} & \shortstack[r]{最大\\語長\\（出力）} \\
\midrule
\endfirsthead
\toprule
\shortstack[r]{\(k\)} & \shortstack[r]{\(p_k\)} & \shortstack[r]{入力\\頂点} & \shortstack[r]{入力\\辺} & \shortstack[r]{巡回的\\SCC} & \shortstack[r]{認証済み} & \shortstack[r]{未認証\\頂点} & \shortstack[r]{出力\\頂点} & \shortstack[r]{出力\\辺} & \shortstack[r]{最大\\語長\\（出力）} \\
\midrule
\endhead
\bottomrule
\endlastfoot
0 & --- & 4 & 12 & 1 & 0 & 4 & 4 & 12 & 1 \\
1 & 3 & 8 & 17 & 1 & 0 & 7 & 5 & 12 & 2 \\
2 & 7 & 30 & 58 & 4 & 2 & 16 & 9 & 19 & 4 \\
3 & 11 & 90 & 159 & 3 & 0 & 68 & 39 & 84 & 8 \\
4 & 13 & 462 & 812 & 4 & 1 & 342 & 250 & 524 & 15 \\
5 & 17 & 3985 & 7020 & 11 & 8 & 2913 & 2059 & 4058 & 38 \\
6 & 19 & 36843 & 57660 & 73 & 70 & 15933 & 8890 & 16696 & 56 \\
7 & 23 & 194242 & 293181 & 23 & 21 & 84824 & 40051 & 73286 & 98 \\
8 & 29 & 1101968 & 1542910 & 311 & 310 & 269245 & 100401 & 171256 & 210 \\
9 & 31 & 2856938 & 3256019 & 1084 & 1084 & 0 & 0 & 0 & --- \\
\end{longtable}
\endgroup

% Markdown AST block 164: Para
10段階の全SCC数は1、2、16、25、123、1060、20800、109409、832175、2850774である。入力グラフの最大語長が、その直前の出力グラフの最大語長より小さくなることがある（段階9では207であり、段階8では210である）。持ち上げは、すべての開始剰余に対してゼロ剰余を強制する語をもつ辺をすべて取り除くからである。第4.3節のアンカー規則が働くことは一度もなかった。(2)の法によるすべての剰余と四つのセルの直積を直接作ると、頂点数は
\[
4\varphi(40112098026)=30656102400
\]
となるが、段階的な構成ではこの直積を作らない。

% BEGIN PUBLIC ADDITION figure:stage-counts
頂点数の推移から、持ち上げによる拡大と各段階内での削減の様子が分かる。

\begin{figure}[htbp]
\centering
\begin{tikzpicture}[paperfig,x=1.16cm,y=.61cm]
\draw[fedge] (-.3,0)--(9.5,0) node[right] {$k$};
\draw[fedge] (0,-.2)--(0,7.1);
\node[anchor=west] at(.05,7.3) {$\log_{10}(N+1)$};
\draw[black!18] (0,1)--(9,1);
\node[anchor=east] at(-.12,1) {$ 1 $};
\draw[black!18] (0,2)--(9,2);
\node[anchor=east] at(-.12,2) {$ 2 $};
\draw[black!18] (0,3)--(9,3);
\node[anchor=east] at(-.12,3) {$ 3 $};
\draw[black!18] (0,4)--(9,4);
\node[anchor=east] at(-.12,4) {$ 4 $};
\draw[black!18] (0,5)--(9,5);
\node[anchor=east] at(-.12,5) {$ 5 $};
\draw[black!18] (0,6)--(9,6);
\node[anchor=east] at(-.12,6) {$ 6 $};
\draw (0,0)--(0,-.12);
\node[below] at(0,-.12) {$ 0 $};
\node[below] at(0,-.62) {$ -- $};
\draw (1,0)--(1,-.12);
\node[below] at(1,-.12) {$ 1 $};
\node[below] at(1,-.62) {$ 3 $};
\draw (2,0)--(2,-.12);
\node[below] at(2,-.12) {$ 2 $};
\node[below] at(2,-.62) {$ 7 $};
\draw (3,0)--(3,-.12);
\node[below] at(3,-.12) {$ 3 $};
\node[below] at(3,-.62) {$ 11 $};
\draw (4,0)--(4,-.12);
\node[below] at(4,-.12) {$ 4 $};
\node[below] at(4,-.62) {$ 13 $};
\draw (5,0)--(5,-.12);
\node[below] at(5,-.12) {$ 5 $};
\node[below] at(5,-.62) {$ 17 $};
\draw (6,0)--(6,-.12);
\node[below] at(6,-.12) {$ 6 $};
\node[below] at(6,-.62) {$ 19 $};
\draw (7,0)--(7,-.12);
\node[below] at(7,-.12) {$ 7 $};
\node[below] at(7,-.62) {$ 23 $};
\draw (8,0)--(8,-.12);
\node[below] at(8,-.12) {$ 8 $};
\node[below] at(8,-.62) {$ 29 $};
\draw (9,0)--(9,-.12);
\node[below] at(9,-.12) {$ 9 $};
\node[below] at(9,-.62) {$ 31 $};
\node[anchor=east] at(-.15,-.4) {$k$};
\node[anchor=east] at(-.15,-.9) {$p_k$};
\draw[solid,line width=.85pt] plot coordinates {(0,0.698970004) (1,0.954242509) (2,1.491361694) (3,1.959041392) (4,2.665580991) (5,3.600537294) (6,4.566366774) (7,5.288345377) (8,6.042169377) (9,6.455900968)};
\node[draw,circle,fill=white,inner sep=1.2pt] at(0,0.698970004){};
\node[draw,circle,fill=white,inner sep=1.2pt] at(1,0.954242509){};
\node[draw,circle,fill=white,inner sep=1.2pt] at(2,1.491361694){};
\node[draw,circle,fill=white,inner sep=1.2pt] at(3,1.959041392){};
\node[draw,circle,fill=white,inner sep=1.2pt] at(4,2.665580991){};
\node[draw,circle,fill=white,inner sep=1.2pt] at(5,3.600537294){};
\node[draw,circle,fill=white,inner sep=1.2pt] at(6,4.566366774){};
\node[draw,circle,fill=white,inner sep=1.2pt] at(7,5.288345377){};
\node[draw,circle,fill=white,inner sep=1.2pt] at(8,6.042169377){};
\node[draw,circle,fill=white,inner sep=1.2pt] at(9,6.455900968){};
\draw[solid,line width=.85pt] (.6,6.2)--(1.3,6.2);
\node[draw,circle,fill=white,inner sep=1.2pt] at(.95,6.2){};
\node[anchor=west] at(1.45,6.2) {\figtext{Input vertices}{入力頂点数}};
\draw[dashed,line width=.85pt] plot coordinates {(0,0.698970004) (1,0.903089987) (2,1.230448921) (3,1.838849091) (4,2.535294120) (5,3.464489547) (6,4.202324813) (7,4.928523868) (8,5.430149260) (9,0.000000000)};
\node[draw,rectangle,fill=white,inner sep=1.2pt] at(0,0.698970004){};
\node[draw,rectangle,fill=white,inner sep=1.2pt] at(1,0.903089987){};
\node[draw,rectangle,fill=white,inner sep=1.2pt] at(2,1.230448921){};
\node[draw,rectangle,fill=white,inner sep=1.2pt] at(3,1.838849091){};
\node[draw,rectangle,fill=white,inner sep=1.2pt] at(4,2.535294120){};
\node[draw,rectangle,fill=white,inner sep=1.2pt] at(5,3.464489547){};
\node[draw,rectangle,fill=white,inner sep=1.2pt] at(6,4.202324813){};
\node[draw,rectangle,fill=white,inner sep=1.2pt] at(7,4.928523868){};
\node[draw,rectangle,fill=white,inner sep=1.2pt] at(8,5.430149260){};
\node[draw,rectangle,fill=white,inner sep=1.2pt] at(9,0.000000000){};
\draw[dashed,line width=.85pt] (.6,5.65)--(1.3,5.65);
\node[draw,rectangle,fill=white,inner sep=1.2pt] at(.95,5.65){};
\node[anchor=west] at(1.45,5.65) {\figtext{After removal}{除去後の頂点数}};
\draw[densely dotted,line width=.85pt] plot coordinates {(0,0.698970004) (1,0.778151250) (2,1.000000000) (3,1.602059991) (4,2.399673721) (5,3.313867220) (6,3.948950610) (7,4.602624207) (8,5.001742364) (9,0.000000000)};
\node[draw,diamond,fill=white,inner sep=1.2pt] at(0,0.698970004){};
\node[draw,diamond,fill=white,inner sep=1.2pt] at(1,0.778151250){};
\node[draw,diamond,fill=white,inner sep=1.2pt] at(2,1.000000000){};
\node[draw,diamond,fill=white,inner sep=1.2pt] at(3,1.602059991){};
\node[draw,diamond,fill=white,inner sep=1.2pt] at(4,2.399673721){};
\node[draw,diamond,fill=white,inner sep=1.2pt] at(5,3.313867220){};
\node[draw,diamond,fill=white,inner sep=1.2pt] at(6,3.948950610){};
\node[draw,diamond,fill=white,inner sep=1.2pt] at(7,4.602624207){};
\node[draw,diamond,fill=white,inner sep=1.2pt] at(8,5.001742364){};
\node[draw,diamond,fill=white,inner sep=1.2pt] at(9,0.000000000){};
\draw[densely dotted,line width=.85pt] (.6,5.1)--(1.3,5.1);
\node[draw,diamond,fill=white,inner sep=1.2pt] at(.95,5.1){};
\node[anchor=west] at(1.45,5.1) {\figtext{After contraction}{縮約後の頂点数}};
\node[flabel,anchor=east] at(8.6,.75) {$N=0$};
\draw[fedge] (8.65,.6)--(8.98,.04);
\node[align=center,text width=12cm] at(4.5,-1.65)
{\figtext{A lift can enlarge the graph; removal and contraction reduce each stage.}{持ち上げでグラフは拡大し得る。各段階で除去と縮約を行う。}};
\end{tikzpicture}

\caption{表6の頂点数。縦軸は\(\log_{10}(N+1)\)であり、最終的な空の出力\(N=0\)を高さゼロに表示できる。実線と丸は入力、破線と四角は除去後、点線と菱形は縮約後の頂点数を表す。点を結ぶ線は推移を見やすくするためのもので、段階間の単調減少を意味しない。段階九では入力の巡回SCCすべてが認証される。}
\label{fig:stage-counts}
\end{figure}
% END PUBLIC ADDITION figure:stage-counts

% Markdown AST block 165: 5.4 Proofs of Theorem 1 and Corollary 2
\Needspace{6\baselineskip}
\subsection{定理1と系2の証明}\label{sec:5.4}

% Markdown AST block 166: Proof of Theorem 1.
\begin{proof}[定理1の証明]\leavevmode\label{proof:1}
(i) 有限計算により、パラメータ(21)に対し、\(s=3\) および \(U_3=\varnothing\) として、(17)のグラフ、保持集合、および細分を検証する。第9節で詳述するように、この検証は完全な辺集合と成分の条件を含む。要約された件数だけでは証明書にはならない。定理10により(1)が得られる。

% Markdown AST block 167: Para
別の証明として、表4により、定理19の仮定が、\(S_0=\{2,5\}\)、\(\mathcal C=\{-4,-2,2,4,6\}\)（補題23）、\(K=2048\)、素数 \(3,11,13\)、および \(G_3^{\rm out}=\varnothing\) に対して成り立つことを検証できる。このとき、
\[
M=10\cdot3\cdot11\cdot13=4290
\]
に対する(20)が(1)そのものである。

% Markdown AST block 168: OrderedList
\begin{enumerate}
\def\labelenumi{(\roman{enumi})}
\setcounter{enumi}{1}
\tightlist
\item
  表6により、定理19の仮定が、\(S_0=\{2\}\)、\(\mathcal C=\{-1,1,3\}\)（補題24）、\(K=4\)、九つの素数 \(3,7,11,13,17,19,23,29,31\)（相異なり、それぞれ \(10\) と互いに素で、\(S_0\) に属さない）、および \(G_9^{\rm out}=\varnothing\) に対して成り立つことを検証できる。このとき、

% Markdown AST block 169: OrderedList
(20)を
\[
M=2\cdot3\cdot7\cdot11\cdot13\cdot17\cdot19\cdot23\cdot29\cdot31
=40112098026
\]
に対して適用したものが(2)である。\qedhere
\end{enumerate}
\end{proof}

% Markdown AST block 170: Proof of Corollary 2.
\begin{proof}[系2の証明]\leavevmode\label{proof:2}
(7)により、最終的に \(Q_n>13\) が比 \(7/5\) の数列について成り立ち、\(Q_n>31\) が比 \(5/2\) の数列について成り立つ。添字の任意の下限を与えたとき、その下限と増大性の閾値の両方を超えたところで定理1を適用する。(1)の最大公約数の素因数は \(\{2,3,5,11,13\}\) に属し、(2)の最大公約数の素因数は \(\{2,3,7,11,13,17,19,23,29,31\}\) に属する。どちらの場合も、それは \(Q_n\) の真の約数である。したがって、任意に大きな添字で \(Q_n\) は合成数となる。これは、いずれの場合にも系20を適用したものである。\qedhere
\end{proof}

% Markdown AST block 171: 6. Quantitative waiting time for 7/5
\Needspace{6\baselineskip}
\section{7/5に対する待ち時間の定量的評価}\label{sec:6}

% Markdown AST block 172: 6.1 A finite version of nonperiodicity
\subsection{非周期性の有限版}\label{sec:6.1}

% Markdown AST block 173: Lemma 25.
\begin{samepage}
\begin{lemma}\leavevmode\label{thm:25}
\(u\in\mathbb N\)、\(1\le d\le\ell\) とし、
\[
c_{u+i+d}=c_{u+i}
\]
が \(0\le i<\ell-d\) に対して成り立つと仮定する。
\[
D_i=Q_{u+i+d}-Q_{u+i}.
\]
とおく。このとき、
\[
b^{\ell-d}\mid D_0.
\tag{22}\label{eq:22}
\]
が成り立つ。\(D_0>0\) ならば、\(\nu_b(D_0)=\max\{k\ge0:b^k\mid D_0\}\) として、
\[
\begin{aligned}
\ell&\le d+\nu_b(D_0),\\
b^{\ell-d}&\le D_0<r^d(Q_u+1)-Q_u.
\end{aligned}
\tag{23}\label{eq:23}
\]
が成り立つ。
\label{thm-end:25}
\end{lemma}
\end{samepage}

% Markdown AST block 174: Proof.
\begin{proof}\leavevmode\label{proof:25}
\(0\le i<\ell-d\) に対して、漸化式(6)の差を取ると、
\[
bD_{i+1}=(aQ_{u+i+d}+c_{u+i+d})-(aQ_{u+i}+c_{u+i})=aD_i.
\]
を得る。帰納法により、
\[
b^kD_k=a^kD_0
\]
が \(0\le k\le\ell-d\) に対して成り立ち、\(\gcd(a^k,b^k)=1\) から(22)が得られる。

% Markdown AST block 175: Para
\(D_0>0\) ならば、正の約数 \(b^{\ell-d}\) は \(D_0\) 以下であるから、\(\ell-d\le\nu_b(D_0)\) となる。

% Markdown AST block 176: Para
最後に、
\[
Q_{u+d}\le x_{u+d}=r^dx_u<r^d(Q_u+1).
\]
を得る。整数 \(b\) は素数である必要はない。\qedhere
\end{proof}

% Markdown AST block 177: Lemma 26.
\begin{samepage}
\begin{lemma}\leavevmode\label{thm:26}
\(x_n\ge b/(a-b)\) ならば、\(Q_{n+1}\ge Q_n+1\) である。特に、\(Q_u\ge b/(a-b)\) ならば、\(Q_n\) は \(n\ge u\) に対して狭義単調増加であり、補題25において
\[
D_0=Q_{u+d}-Q_u\ge d\ge1
\]
が成り立つ。
\label{thm-end:26}
\end{lemma}
\end{samepage}

% Markdown AST block 178: Proof.
\begin{proof}\leavevmode\label{proof:26}
\[
rx_n-x_n=x_n(a-b)/b\ge1,
\]
であるから、
\[
\lfloor rx_n\rfloor\ge\lfloor x_n+1\rfloor=Q_n+1.\qedhere
\]
\end{proof}

% Markdown AST block 179: Para
\(7/5\) では \(b/(a-b)=5/2\) であるため、\(Q_u\ge3\) で十分である。整数部分が \(0\) または \(1\) にとどまる場合に、差が正であると仮定してはならない。

% Markdown AST block 180: 6.2 From a rank certificate to a waiting time
\Needspace{6\baselineskip}
\subsection{ランク証明書から待ち時間へ}\label{sec:6.2}

% Markdown AST block 181: Para
\(H\) を、すべての辺の長さが1の有限ラベル付きグラフとする。さらに、\(V(H)\) の\emph{ブロック}への分割と、以下のデータを与える。

% Markdown AST block 182: BulletList
\begin{itemize}
\tightlist
\item
  (Q-rank) 整数 \(\rho(B)\in[1,\rho_{\max}]\) を各ブロック \(B\) に与え、異なるブロック間のすべての辺 \(u\to v\) が
  \[
  \rho(B(v))\ge\rho(B(u))+1.
  \]
  を満たすものとする。\(R=\rho_{\max}-1\) とおく。
\item
  (Q-block) 各ブロックは、内部辺をもたないか、または\emph{周期的}である。すなわち、\(d_B\in[1,D]\)、\(h_B:B\to\mathbb Z\)、\(\lambda_B:\mathbb Z/d_B\to\mathcal C\) が存在し、すべての内部辺 \(u\xrightarrow{c}v\) に対して
  \[
  \begin{aligned}
  &h_B(v)\equiv h_B(u)+1\pmod{d_B}\qquad\text{and}\\
  &c=\lambda_B(h_B(u)\bmod d_B).
  \end{aligned}
  \]
  が成り立つ。
\item
  (Q-cyc) 整数 \(\eta(B)\in[0,m]\) を各ブロックに与え、
  \[
  \begin{aligned}
  &\eta(B)\ge\mathbf 1[B\text{ periodic}],\qquad\text{and}\\
  &\eta(B(v))\ge\eta(B(u))+\mathbf 1[B(v)\text{ periodic}]
  \end{aligned}
  \]
  が、異なるブロック間のすべての辺に対して成り立つものとする。
\end{itemize}

% Markdown AST block 183: Para
内部辺をもつブロックが認証されずに残ることは許されない。ブロックはSCCである必要はない。SCC分解を用いる場合、その非巡回商グラフを通じてこのようなデータが得られる。\(\rho\) はそのブロックで終わる最長のブロック列の長さ、\(\eta\) はそのような列に含まれる周期的ブロック数の最大値とする。

% Markdown AST block 184: Para
(Q-rank)により、\(H\) の任意の有限ウォークは、異なるブロック間の辺を高々 \(R\) 本しか通らず、各ブロックを訪れる時刻は一つの連続区間をなす。また、(Q-cyc)により、訪れる周期的ブロックは高々 \(m\) 個である。

% Markdown AST block 185: Para
\emph{長さ \(T\) の軌道区間が \(H\) 内にある}とは、軌道の状態 \(\sigma(0),\ldots,\sigma(T)\)（\(H\) の座標で表したもの）が \(H\) の頂点であり、\(\sigma(n)\xrightarrow{c_n}\sigma(n+1)\) が \(H\) の辺であることが \(0\le n<T\) に対して成り立つことをいう。

% Markdown AST block 186: Theorem 27 (waiting time from a rank certificate).
\begin{samepage}
\begin{theorem}[ランク証明書から得られる待ち時間]\leavevmode\label{thm:27}
\[
\begin{aligned}
&\kappa=\log_br>0,\qquad \gamma\ge1+\kappa,\\
&P=Q_0\ge b/(a-b),\qquad\text{and}\qquad A\ge\log_b(P+2).
\end{aligned}
\]
とする。長さ \(T\) の軌道区間が \(H\) 内にあるならば、
\[
T\le R\gamma^m+(A+\gamma D)\sum_{j=0}^{m-1}\gamma^j .
\tag{24}\label{eq:24}
\]
が成り立つ。\(m=0\) の場合、和は空であり、\(T\le R\) となる。
\label{thm-end:27}
\end{theorem}
\end{samepage}

% Markdown AST block 187: Proof.
\begin{proof}\leavevmode\label{proof:27}
\emph{第1段階：一つの周期的ブロック内での滞在。} 軌道区間が周期的ブロック \(B\) に時刻 \(u\) で入り、時刻 \(u,\ldots,u+\ell\) の状態が \(B\) に属するとする。このとき、時刻 \(u,\ldots,u+\ell-1\) の辺は内部辺である。(Q-block)を用いた帰納法により、
\[
c_{u+i}=\lambda_B((h_B(\sigma(u))+i)\bmod d_B),
\]
を得る。したがって、
\[
c_{u+i+d_B}=c_{u+i}
\]
が \(0\le i<\ell-d_B\) に対して成り立つ。

% Markdown AST block 188: Para
\(\ell\ge d_B\) ならば補題25を適用でき、補題26（\(Q_u\ge P\ge b/(a-b)\) として）から \(D_0>0\) が得られるので、
\[
\begin{aligned}
&b^{\ell-d_B}\le D_0<r^{d_B}(Q_u+1)\qquad\text{and}\\
&\ell-d_B<\log_b(Q_u+1)+d_B\kappa.
\end{aligned}
\]
となる。

% Markdown AST block 189: Para
\[
Q_u\le x_u=r^u\xi<r^u(P+1),
\]
であるから、
\[
Q_u+1<r^u(P+1)+1\le r^u(P+2),
\]
が成り立ち、
\[
\log_b(Q_u+1)<u\kappa+A.
\]
となる。以上を合わせると、
\[
\ell<A+\kappa u+(1+\kappa)d_B\le A+(\gamma-1)u+\gamma D.
\]
を得る。

% Markdown AST block 190: Para
\(\ell<d_B\) の場合にも、
\[
\ell<d_B\le D\le\gamma D.
\]
であるから、同じ上界は自明に成り立つ。したがって、退出時刻（または軌道区間の終了時刻）は、
\[
u+\ell<\gamma u+A+\gamma D.
\]
を満たす。

% Markdown AST block 191: Para
\emph{第2段階：量の整理。} \(B_0,\ldots,B_s\)（\(s\le R\)）を訪れるブロックとし、それぞれへの進入時刻を \(u_0=0<u_1<\cdots<u_s\) とする。\(\ell_k\) を \(B_k\) 内で通る内部辺の本数とすると、
\[
u_{k+1}=u_k+\ell_k+1\qquad \text{and}\qquad T=u_s+\ell_s;
\]
となる。辺のないブロックでは \(\ell_k=0\) である。

% Markdown AST block 192: Para
\(k_1<\cdots<k_{m'}\)（\(m'\le m\)）を周期的ブロックの添字とし、\(k_0=0\)、\(x_0=0\) とおく。
\[
x_j=u_{k_j}+\ell_{k_j}
\]
を \(j\) 番目の周期的ブロックからの退出時刻とし、
\[
\begin{aligned}
&\delta_j=k_j-k_{j-1}\qquad \text{for}\qquad 1\le j\le m',\qquad\text{and}\\
&\delta_{m'+1}=s-k_{m'}.
\end{aligned}
\]
とおく。周期的ブロックの間には辺のないブロックしか現れないので、時刻はブロック間の辺1本につき1だけ進む。すなわち、
\[
\begin{aligned}
&u_{k_j}=x_{j-1}+\delta_j,\qquad T=x_{m'}+\delta_{m'+1},\qquad\text{and}\\
&\sum_{j=1}^{m'+1}\delta_j=s\le R.
\end{aligned}
\]
である。

% Markdown AST block 193: Para
\emph{第3段階：展開。} 第1段階により、
\[
x_j<\gamma u_{k_j}+A+\gamma D=\gamma(x_{j-1}+\delta_j)+A+\gamma D.
\]
である。帰納法から、
\[
x_{m'}<\sum_{j=1}^{m'}\gamma^{m'-j+1}\delta_j+(A+\gamma D)\sum_{j=0}^{m'-1}\gamma^j .
\]
を得る。\(\gamma\ge1\) より、\(\gamma^{m'-j+1}\le\gamma^{m'}\) および \(1\le\gamma^{m'}\) が成り立つので、
\[
\begin{aligned}
T&=x_{m'}+\delta_{m'+1}\\
&<\gamma^{m'}\sum_{j=1}^{m'+1}\delta_j+(A+\gamma D)\sum_{j<m'}\gamma^j\\
&\le R\gamma^{m}+(A+\gamma D)\sum_{j<m}\gamma^j .
\end{aligned}
\]
となる。\(m'=0\) の場合には、
\[
T=s\le R.\qedhere
\]
\end{proof}

% Markdown AST block 194: Para
ブロック間の各辺の寄与は、その後に訪れる周期的ブロックごとに一度ずつ増幅され、その倍率は全体で高々 \(\gamma^m\) である。この証明は有限ウォークだけを扱い、無限ウォークも \(\xi\) の標本も用いない。認証されない巡回的ブロックが一つでもあれば、ウォークはその中に永久にとどまりうるため、この定理は適用できない。

% BEGIN PUBLIC ADDITION figure:rank-certificate
ランクは経路上のブロック数や周期ブロック数を数え、周期はブロック内の出力を表す。

\begin{figure}[htbp]
\centering
\begin{tikzpicture}[paperfig]
\node[ftitle] at(6,1.85){\figtext{Illustrative condensation DAG (not the full computed graph)}{SCC凝縮DAGの概念例（実計算グラフ全体ではない）}};
\node[fbox] (a) at(0,0){$A$\\$\rho=1,\ \eta=0$};
\node[fbox,double] (b) at(4,0.6){$B$\\$\rho=2,\ \eta=1$\\$d_B=2$};
\node[fbox] (c) at(4,-1.3){$C$\\$\rho=2,\ \eta=0$};
\node[fbox,double] (d) at(8,0.6){$D$\\$\rho=3,\ \eta=2$\\$d_D=3$};
\node[fbox] (e) at(12,0){$E$\\$\rho=4,\ \eta=2$};
\draw[fedge] (a)--(b);\draw[fedge] (a)--(c);\draw[fedge] (b)--(d);\draw[fedge] (c)--(d);\draw[fedge] (d)--(e);
\draw[fedge,dashed] (c) to[bend right=12](e);
\node[align=center,text width=12.7cm] at(6,-2.55)
 {\figtext{Double border: periodic block. Single border: no internal edge.}{二重枠：周期ブロック。一重枠：内部辺なし。}\\[4pt]
 $\rho$: \figtext{maximum number of blocks}{最大ブロック数}\qquad
 $\eta$: \figtext{maximum number of periodic blocks}{最大周期ブロック数}\\[4pt]
 $d_B$: \figtext{internal output period}{内部出力の周期}};
\node[fbox,text width=12cm] at(6,-4.3)
 {\figtext{Computed $7/5$ certificate (Table 7)}{実計算の $7/5$ 証明書（表7）}\\[5pt]
 $491\,275$ \figtext{blocks}{ブロック},\quad $35$ \figtext{cyclic blocks}{巡回ブロック}\\[4pt]
 $\max\rho=65,\quad\max\eta=6,\quad\max d_B=13$};
\end{tikzpicture}

\caption{SCC凝縮DAGの概念例で三つの量を区別する。\(\rho\)は当該ブロックで終わる経路上の最大ブロック数、\(\eta\)はそのような経路上の最大周期ブロック数、\(d_B\)は内部出力の周期である。五ブロックの図は模式図である。別枠には表7の実際の\(7/5\)証明書の値を記しており、これらの値は五ブロックの例の値ではない。}
\label{fig:rank-certificate}
\end{figure}
% END PUBLIC ADDITION figure:rank-certificate

% Markdown AST block 195: 6.3 The certificate for 7/5
\Needspace{6\baselineskip}
\subsection{7/5の証明書}\label{sec:6.3}

% Markdown AST block 196: Lemma 28 (reduction of the modulus).
\begin{samepage}
\begin{lemma}[法の縮小]\leavevmode\label{thm:28}
\(\gcd(Q_n,4290)=1\) が \(0\le n\le T\) に対して成り立つと仮定する。このとき、\(c_n\in\mathcal C=\{-4,-2,2,4,6\}\) が \(0\le n<T\) に対して成り立つ。また、\(M=429=3\cdot11\cdot13\) および \(K=2048\) とすると、長さ \(T\) の軌道区間は、グラフ
\[
H=G(429,2048)\ \text{restricted to labels in }\mathcal C .
\]
内にある。
\label{thm-end:28}
\end{lemma}
\end{samepage}

% Markdown AST block 197: Proof.
\begin{proof}\leavevmode\label{proof:28}
補題23から文字集合が得られる。\(\gcd(Q_n,429)=1\) なので、状態 
\[
(Q_n\bmod429,\lfloor2048\theta_n\rfloor)
\]
 は \(V(429,2048)\) に属する。また、命題5により、連続する状態はラベル \(c_n\) の辺で隣接する。\qedhere
\end{proof}

% Markdown AST block 198: Para
\(\gcd(5,429)=1\) であるから、合同式(E2)は一意な解
\[
z\equiv5^{-1}(7q+c)\pmod{429};
\]
をもつ。これは法を \(4290\) から \(429\) へ縮小した結果である（縮小前の法では \(5\) は単元ではない）。\(2\) と \(5\) を避ける条件は文字集合が担っており、素数 \(7\) を避けることは仮定していない。

% Markdown AST block 199: Para
グラフ \(H\) の頂点数は
\[
\varphi(429)\cdot2048=240\cdot2048=491520
\]
である。そのSCC分解と、非巡回商グラフ上で計算したランク \(\rho,\eta\) から、表7のデータを得る。第9.3節の二つの実装は、\(H\) のすべての頂点と辺を列挙し、ブロック分割、各ブロックの属性、および各周期的ブロックの位相ラベルについて一致する。また、すべての辺に対して(Q-rank)、(Q-block)、(Q-cyc)の不等式を検証する。

% Markdown AST block 200: Table 7.
% Caption is rendered at the start of the immediately following table.

% Markdown AST block 201: Table 7
\begingroup
\small
\setlength{\tabcolsep}{4pt}
\renewcommand{\arraystretch}{1.15}
\setlength{\LTcapwidth}{\textwidth}
\begin{longtable}{@{}>{\raggedright\arraybackslash}p{\dimexpr0.65\dimexpr\linewidth-8pt\relax} >{\raggedright\arraybackslash}p{\dimexpr0.35\dimexpr\linewidth-8pt\relax}@{}}
\caption{\(7/5\) のランク証明書。\(M=429\)、\(K=2048\) を用いる。}\label{tab:7}\\
\toprule
項目 & 値 \\
\midrule
\endfirsthead
\toprule
項目 & 値 \\
\midrule
\endhead
\bottomrule
\endlastfoot
許容頂点 & 491520 \\
ラベル付き辺 & 891990 \\
ブロック（SCC） & 491275 \\
巡回的ブロック & 35 \\
認証済み巡回的ブロック & 35 \\
未認証の巡回的ブロック & 0 \\
\(\max\rho\) & 65 \\
\(\max\eta\) & 6 \\
\(\max d_B\) & 13 \\
\end{longtable}
\endgroup

% Markdown AST block 202: Para
したがって、(Q-rank)は \(\rho_{\max}=65\)、\(R=64\) として、(Q-cyc)は \(m=6\) として、(Q-block)は \(D=13\) として成り立つ。\(7^4=2401<3125=5^5\) より \((7/5)^4<5\)、\(\kappa=\log_5(7/5)<1/4\) であり、\(\gamma=5/4\) を採用できる。
\[
A=L(P)=\min\{h\ge0:5^h\ge P+2\}\ge\log_5(P+2),
\]
とすると、(24)の定数は、有理数演算によって厳密に次のように求まる。
\[
\begin{aligned}
\sum_{j=0}^{5}\Bigl(\tfrac54\Bigr)^j&=\frac{11529}{1024},\\
64\Bigl(\tfrac54\Bigr)^6+\tfrac54\cdot13\cdot\frac{11529}{1024}
&=\frac{15625}{64}+\frac{749385}{4096}\\
&=\frac{1749385}{4096},
\end{aligned}
\]
\[
\text{and}\quad
1749385<428\cdot4096=1753088,\qquad 11529<12\cdot1024=12288.
\]
したがって、長さ \(T\) の軌道区間が \(H\) 内にあり、\(P\ge3\) であるならば、
\[
T\le\frac{1749385}{4096}+\frac{11529}{1024}\,L(P)<428+12L(P).
\tag{25}\label{eq:25}
\]
となる。

% Markdown AST block 203: Proof of Theorem 3.
\begin{proof}[定理3の証明]\leavevmode\label{proof:3}
\(P=\lfloor\xi\rfloor>13\) および \(T_0=428+12L(P)\) とする。\(\gcd(Q_n,4290)=1\) がすべての \(0\le n\le T_0\) に対して成り立つならば、補題28により、長さ \(T_0\) の軌道区間は \(H\) 内にある。表7の証明書を用いて定理27を適用すると（\(P>13\ge3\) より補題26を適用できる）、
\[
T_0<428+12L(P)=T_0,
\]
となり、矛盾する。

% Markdown AST block 204: Para
したがって、ある \(0\le n\le T_0\) が \(\gcd(Q_n,4290)>1\) を満たす。

% Markdown AST block 205: Para
このような \(n\) に対し、素数 \(p\in\{2,3,5,11,13\}\) が
\[
Q_n\ge Q_0=P>13\ge p,
\]
を割るので、\(Q_n=pv\) は \(v>1\) を満たし、合成数である。

% Markdown AST block 206: Para
最後の主張については、結果を \(\xi'=\xi(7/5)^N\) に適用する。この場合、
\[
\lfloor\xi'(7/5)^n\rfloor=Q_{N+n}
\]
および \(\lfloor\xi'\rfloor=Q_N\) が成り立つ。グラフ \(H\) は \(\xi\) に依存せず、\(P\) は上界 \(Q_u<r^u(P+1)\) を通じてのみ関与する。\qedhere
\end{proof}

% Markdown AST block 207: Para
上界(25)自体は \(P\ge3\) に対して成り立ち、仮定 \(P>13\) は合成数性の結論を導くためにのみ用いる。\(L(P)\) を整数で定義することにより、丸めの問題を避けられる。この値は、\(\lceil\log_5(P+2)\rceil\) に等しく、この等式はすべての \(P\ge0\) に対して成り立つ。

% Markdown AST block 208: 7. Scope of the certificate method
\Needspace{6\baselineskip}
\section{証明書手法の適用範囲}\label{sec:7}

% Markdown AST block 209: Para
本節と次節では、第3--4節の手続きで許される操作を厳密に定め、その適用範囲を調べる。以下では一貫して
\[
\Sigma=\{1-b,\ldots,a-1\},
\]
とし、\(\operatorname{rad}M\) は \(M\) を割る相異なる素数の積とする。

% Markdown AST block 210: 7.1 The general carry alphabet
\Needspace{6\baselineskip}
\subsection{一般の繰り上がりアルファベット}\label{sec:7.1}

% Markdown AST block 211: Para
既約な底 \(a/b\) に対して、
\[
\mathcal C(a,b)=\{c\in\mathbb Z:1-b\le c\le a-1,\ c\equiv b-a\pmod2,\
\gcd(c,ab)=1\}.
\tag{26}\label{eq:26}
\]
と定める。

% Markdown AST block 212: Lemma 29.
\begin{samepage}
\begin{lemma}\leavevmode\label{thm:29}
\(\operatorname{rad}(2ab)\mid M\) とする。\((q,j)\xrightarrow{e}(z,k)\) を \(G(M,K)\) の任意の辺とすると、\(e\in\mathcal C(a,b)\) が成り立つ。
特に、\(c_n\in\mathcal C(a,b)\) は \(\gcd(Q_nQ_{n+1},2ab)=1\) のとき常に成り立つ。逆に、このような \(M\) に対して、\(2ab\) を割る素数が課すラベル条件は、ちょうど \(e\in\mathcal C(a,b)\) である。
\label{thm-end:29}
\end{lemma}
\end{samepage}

% Markdown AST block 213: Proof.
\begin{proof}\leavevmode\label{proof:29}
\(p\mid a\) とする。このとき \(p\nmid b\) であり、(E2) より
\[
bz\equiv e\pmod p,
\]
を得る。したがって、\(p\mid e\) と \(p\mid z\) は同値であるが、\(z\) は単元なので、これは起こらない。

% Markdown AST block 214: Para
\(p\mid b\) とする。このとき
\[
0\equiv aq+e\pmod p,
\]
であるから、\(p\mid e\) と \(p\mid q\) は同値であり、これも起こらない。

% Markdown AST block 215: Para
\(a,b\) が奇数で \(2\mid M\) ならば、
\[
z\equiv q+e\pmod2
\]
である。ここで \(q,z\) は奇数なので、\(e\) は偶数、すなわち \(e\equiv b-a\pmod2\) である。\(a,b\) の一方が偶数ならば、偶奇条件 \(e\equiv b-a\pmod2\) は \(\gcd(e,ab)\) が奇数であることと同値であり、これはすでに示した。逆は、同じ計算を逆にたどればよい。\qedhere
\end{proof}

% Markdown AST block 216: Para
\(7/5\) では
\[
\mathcal C(7,5)=\{-4,-2,2,4,6\},
\]
であり、\(5/2\) では
\[
\mathcal C(5,2)=\{-1,1,3\}:
\]
である。したがって、補題23および24のアルファベットは、ちょうどこの一般形に一致する。
\(\mathcal C(a,b)\) によって \(2ab\) を割るすべての素数の回避を表す場合、証明書から得られる結論は、\(2ab\) を割るすべての素数を素数集合に含む整除性の命題となる。第5節の個別の還元では、これらの素数の一部を使わずに済む（\(7\) は \(7/5\) の場合、\(5\) は \(5/2\) の場合）。

% Markdown AST block 217: 7.2 The procedure and its three soundness notions
\Needspace{6\baselineskip}
\subsection{手続きと三つの健全性の概念}\label{sec:7.2}

% Markdown AST block 218: Para
\textbf{展開を伴う語ラベル付きグラフ。} 本節では、語ラベル付きグラフの頂点集合は、ある \(V(M,K)\) に含まれるものとする。また、各辺 \(u\xrightarrow{w}v\) には、\emph{展開道}
\[
u=u_0\xrightarrow{w_0}u_1\cdots\xrightarrow{w_{|w|-1}}u_{|w|}=v
\]
を記録する。この道は \(G(M,K)\) の辺からなる。1ステップグラフは、語長が1の場合に当たる。
\(M\mid M'\) および \(K'=fK\) に対し、\emph{射影}を
\[
\pi(q,J)=(q\bmod M,\lfloor J/f\rfloor).
\]
と定める。

% Markdown AST block 219: Para
\textbf{擬軌道。} \emph{右向きの \(M\)-擬軌道}とは、列 \((q_n,\theta_n,c_n)_{n\ge n_0}\) であって、\(q_n\in(\mathbb Z/M)^\times\)、\(\theta_n\in[0,1)\)、\(c_n\in\mathbb Z\) および
\[
a\theta_n-b\theta_{n+1}=c_n\ \text{(exactly, in }\mathbb R),\qquad
bq_{n+1}\equiv aq_n+c_n\pmod M .
\]
を満たすものをいう。\emph{左向き}の擬軌道は、\(n\le n_0\) に対して同じ条件を満たすものとする。
尾部が \(M\) の素因数を回避する真の軌道からは、\(q_n=Q_n\bmod M\) とすることで右向きの擬軌道が得られる。また、擬軌道からは、任意の \(K\) に対して \((q_n,\lfloor K\theta_n\rfloor)\) という \(G(M,K)\) のウォークが得られる。これは、実数上の恒等式と半開区間のセルだけを用いる命題5の証明による。
\(M\) を法とする擬軌道を、\(M\) の任意の約数を法として還元しても擬軌道になる。真の軌道にあって擬軌道にはないものは、剰余と小数部分を結び付ける一つの整数列 \(Q_n\) である。

% Markdown AST block 220: Para
\textbf{操作。} \emph{実行}とは、以下の操作によって構成される語ラベル付きグラフの有限列 \(\Gamma_0,\ldots,\Gamma_s\) をいう。各グラフのパラメータは \((M_t,K_t)\) であり、\(M_t\mid M_{t+1}\) および \(K_t\mid K_{t+1}\) を満たす。

% Markdown AST block 221: BulletList
\begin{itemize}
\tightlist
\item
  \begin{enumerate}
  \def\labelenumi{(\Roman{enumi})}
  \tightlist
  \item
    \emph{初期グラフ。} \(\Gamma_0=G(M_0,K_0)\) とし、すべての許容状態を頂点とし、\(\Sigma\) のすべてのラベルを使う。
    または、最終的な法が \(\operatorname{rad}(2ab)\) の倍数になる場合には \(\Gamma_0=G'(R_0,K_0)\) とする。後者は、頂点 \((q,j)\)、\(q\in(\mathbb Z/R_0)^\times\)
    （\(R_0=1\) も許し、\(\gcd(R_0,2ab)=1\) とする）をもち、辺のラベルは \(e\in\mathcal C(a,b)\) であって
    \[
    bz\equiv aq+e\pmod{R_0}
    \]
    および (E3) を満たすグラフである。この形では \(M_0=\operatorname{rad}(2ab)R_0\) と定め、\(2ab\) の素因数を法とする剰余はアルファベットによって表現される（補題29および命題36）。
  \end{enumerate}
\item
  \begin{enumerate}
  \def\labelenumi{(\Alph{enumi})}
  \setcounter{enumi}{15}
  \tightlist
  \item
    \emph{位相による除去。} \(\Gamma_t\) を、認証されていない巡回的SCCとその内部辺の和集合で置き換える。各成分では、共通の根からの道の総語長を用い、すべての内部辺とその各文字に (19) を適用する（第4.2節、語長1の場合は補題7）。系22により、結果はこれらの選択に依存しない。
  \end{enumerate}
\item
  \begin{enumerate}
  \def\labelenumi{(\Alph{enumi})}
  \setcounter{enumi}{2}
  \tightlist
  \item
    \emph{縮約。} 第4.3節の操作 \(P\) を行う。展開道は連結される。
  \end{enumerate}
\item
  \begin{enumerate}
  \def\labelenumi{(\Alph{enumi})}
  \setcounter{enumi}{17}
  \tightlist
  \item
    \emph{セルの細分。} \(K_{t+1}=fK_t\) とする。頂点 \((q,j)\) を \((q,fj+i)\)、\(0\le i<f\) で置き換え、各辺を、その展開道のすべての細分で置き換える。ただし、各ステップで \(G(M_t,K_{t+1})\) の辺条件を満たすものに限る。
  \end{enumerate}
\item
  \begin{enumerate}
  \def\labelenumi{(\Alph{enumi})}
  \setcounter{enumi}{11}
  \tightlist
  \item
    \emph{法の拡大。} \(M_{t+1}=hM_t\) とする。各辺を、その展開道に沿った \(M_{t+1}\) を法とする剰余列のすべてで置き換える。ただし、各剰余が単元であり、各ステップで (E2) を満たすことを課す（\(b\) の逆元の存在は仮定せず、すべての解を取る）。素数 \(p\nmid ab\) による持ち上げ（第4.4節）は \(h=p\) の場合であり、(I) の第2の形では、新たな素数は \(2ab\) を割らない。
  \end{enumerate}
\item
  \begin{enumerate}
  \def\labelenumi{(\Alph{enumi})}
  \tightlist
  \item
    \emph{繰り上がりの制限。} \(\operatorname{rad}(2ab)\mid M_t\) のとき、ラベルを \(\mathcal C(a,b)\) に制限する。
  \end{enumerate}
\end{itemize}

% Markdown AST block 222: Para
これらの操作だけを許す手続きを \(\mathfrak M_0\) とする。実行が\emph{成功する}とは、\(\Gamma_s=\varnothing\) となることをいう。
第5節の実行はその具体例である。表1は (I)、(P)、(R) を使い、表4および6は \(R_0=1\) とした (I) の第2の形から始め、次いで (P)、(C)、(L) を使う。

% Markdown AST block 223: Para
\textbf{健全性。} \(\omega\) を \(G(M_*,K_*)\) の無限ウォーク（右向きまたは左向き）とし、\(M_t\mid M_*\) および \(K_t\mid K_*\) を満たし、その出力はウォークの向きに最終的周期的でないものとする。
\emph{\(\omega\) の尾部が \(\Gamma_t\) のウォークになる}とは、\(\Gamma_t\) の無限ウォークが存在して、その連続する展開道が \(\omega\) の状態列を \((M_t,K_t)\) に射影したものに一致し、その出力が \(\omega\) の繰り上がり列の尾部となることをいう。操作の健全性を次のように定める。

% Markdown AST block 224: BulletList
\begin{itemize}
\tightlist
\item
  この性質を、そのようなすべての \(\omega\) に対して保存する操作を、\emph{ウォークに関して健全}という。
\item
  最終的周期的でない繰り上がりをもつすべての \(M\)-擬軌道（\(M_t\mid M\)）を \(\omega\) の代わりに用いたとき、この性質を保存する操作を、\emph{擬軌道に関して健全}という。
\item
  最終的な法の素因数を回避する真の軌道のすべての尾部に対して、この性質を保存する操作を、\emph{軌道に関して健全}という。
\end{itemize}

% Markdown AST block 225: Para
上記の包含関係から、ウォークに関する健全性は擬軌道に関する健全性を導き、擬軌道に関する健全性は軌道に関する健全性を導く。各操作の軌道に関する健全性が、定理10および19の内容である。

% Markdown AST block 226: Lemma 30 (preservation).
\begin{samepage}
\begin{lemma}[保存]\leavevmode\label{thm:30}
操作 (P)、(C)、(R)、(L)、(A) は、右向きおよび左向きのウォークに関して健全である。
\label{thm-end:30}
\end{lemma}
\end{samepage}

% Markdown AST block 227: Proof.
\begin{proof}\leavevmode\label{proof:30}
(R) 以下の補題31により、\(\omega\) の \((M_t,K_{t+1})\) への射影は \(G(M_t,K_{t+1})\) のウォークであり、それをさらに \((M_t,K_t)\) に射影すると現在の展開道になる。したがって、これは細分された辺の一つである。

% Markdown AST block 228: Para
\textbf{(L)} 同様に、\(\omega\) の \(M_{t+1}\) を法とする剰余は、(E2) を満たす単元であるから、列挙されるすべての解の中に含まれる。

% Markdown AST block 229: Para
\textbf{(A)} 補題29により、\(\omega\) のすべてのラベルが \(\mathcal C(a,b)\) に属するという性質は、\(\operatorname{rad}(2ab)\mid M_t\mid M_*\) となった時点で成り立つ。

% Markdown AST block 230: Para
\textbf{(P)} \(\Gamma_t\) の右向きの無限ウォークは、最終的に一つの巡回的SCC \(C\) にとどまる（補題6）。左向きのウォークについては、いったん離れたSCCには再び入れないので、各SCC内にいる時刻の集合は区間となる。有限個の区間が \(\mathbb Z_{\le n_0}\) を覆うため、最も左の区間は下に非有界であり、そのSCCは内部辺を含む。
もし \(C\) が認証されていれば、尾部の出力は周期的になる。右向きのウォークでは、これは補題12(b)による。左向きのウォークでは、補題12(a)の合同式を逆向きに伝播させると、時刻 \(t\le-1\)（最後の頂点 \(v_0\) を基準とする）の文字は \(\lambda((h(v_0)+t)\bmod d)\) に等しいことが分かる。これは \(\omega\) の尾部の非周期性に矛盾する。したがって \(C\) は残され、尾部は残された部分グラフのウォークとなる。

% Markdown AST block 231: Para
\textbf{(C)} 右向きのウォークについては命題16である。左向きのウォークが、ある時点より左でマクロ頂点集合 \(B\) を回避すると仮定する。その部分の頂点はすべて出次数1をもつ。有限性からある頂点が繰り返し現れ、\(v_m=v_{m'}\) となる \(m<m'\) が存在する。このとき \(Z=\{v_m,\ldots,v_{m'-1}\}\) は、出次数1の頂点からなる、後続頂点について閉じたサイクルである。
そのSCCは \(Z\) 自身に等しく、すべての頂点が出次数1をもつ。よって、そのアンカーは \(B\cap Z\) に属し、\(B\) の回避に矛盾する。したがって、ウォークは左に向かって \(B\) を無限回訪れ、各訪問の間の区間は、命題16と同様にマクロ辺になる。\qedhere
\end{proof}

% Markdown AST block 232: Para
両方向で擬軌道に関して健全な操作は、第8節の障害定理に影響を与えることなく \(\mathfrak M_0\) に追加できる。
その一例である語全体の区間検査を、付録E.4（命題73）で扱う。この検査は擬軌道に関して健全だが、ウォークに関しては健全でなく、定理33の対象には含めない。

% Markdown AST block 233: 7.3 Projection and inheritance of certification
\Needspace{6\baselineskip}
\subsection{射影と認証の継承}\label{sec:7.3}

% Markdown AST block 234: Lemma 31 (projection).
\begin{samepage}
\begin{lemma}[射影]\leavevmode\label{thm:31}
\(M\mid M'\) および \(K'=fK\) とする。射影 \(\pi\) は、\((q,J)\xrightarrow{e}(z,H)\) を \(G(M',K')\) の任意の辺とするとき、それを同じラベルをもつ \(G(M,K)\) の辺に写す。
\label{thm-end:31}
\end{lemma}
\end{samepage}

% Markdown AST block 235: Proof.
\begin{proof}\leavevmode\label{proof:31}
(E1) は変わらず、(E2) は法 \(M\) に還元される。(E3) については、\(J=fj+i\)、\(H=fk+\ell\)、\(0\le i,\ell<f\) と書く。このとき
\(eK'=efK\)
であり、
\[
\begin{aligned}
f\bigl(aj-b(k+1)\bigr)&\le afj+ai-bfk-b(\ell+1)\\
&<efK<afj+a(i+1)-bfk-b\ell\\
&\le f\bigl(a(j+1)-bk\bigr),
\end{aligned}
\]
が成り立つ。両端の不等式では、\(i\ge0\)、\(\ell+1\le f\) および \(i+1\le f\)、\(\ell\ge0\) を用いた。\(f\) で割れば、\((j,k)\) に対する (E3) を得る。\qedhere
\end{proof}

% Markdown AST block 236: Lemma 32 (inheritance).
\begin{samepage}
\begin{lemma}[継承]\leavevmode\label{thm:32}
補題31の設定のもと、\(W\subseteq V(M,K)\) および \(W'\subseteq V(M',K')\) は \(\pi(W')\subseteq W\) を満たすとする。また、\(C'\) を \(G_{K'}(W')\) の巡回的SCCとし、根を \(\rho'\)、根からの道の長さを \(h'\) とする。
このとき \(\pi(C')\) は、ある巡回的SCC \(C\)（\(G_K(W)\) の成分）に含まれる。その根からの道の長さを \(h\)、最大公約数を \(d\) とすると、\(d\mid d'\) が成り立つ。さらに、\(C\) が認証されていれば \(C'\) も認証される。
\label{thm-end:32}
\end{lemma}
\end{samepage}

% Markdown AST block 237: Proof.
\begin{proof}\leavevmode\label{proof:32}
補題31および \(\pi(W')\subseteq W\) により、\(C'\) の各内部辺の射影は、誘導部分グラフ \(G_K(W)\) の辺となる。したがって、\(C'\) の像は \(G_K(W)\) 内で強連結であり、内部辺を含むため、一つの巡回的SCC \(C\) に含まれる。

% Markdown AST block 238: Para
選んだ \(\rho'\) から \(x\) への道を射影し、その道に沿って (15) を用いると、
\[
h(\pi x)\equiv h(\pi\rho')+h'(x)\pmod d.
\]
を得る。

% Markdown AST block 239: Para
\(u'\to v'\) を \(C'\) の内部辺とすると、
\[
h'(u')+1-h'(v')\equiv h(\pi u')+1-h(\pi v')
\equiv0\pmod d,
\]
であるから、\(d\mid d'\) が成り立つ。

% Markdown AST block 240: Para
\(C\) が認証されているならば、\(u'\to v'\) のラベルは
\[
\lambda_{h(\pi u')\bmod d}
=\lambda_{(h(\pi\rho')+h'(u'))\bmod d},
\]
であり、これは \(h'(u')\bmod d'\) だけに依存する。したがって \(C'\) も認証される。\qedhere
\end{proof}

% Markdown AST block 241: 7.4 Normal form of finite certificates
\Needspace{6\baselineskip}
\subsection{有限証明書の標準形}\label{sec:7.4}

% Markdown AST block 242: Theorem 33 (contraction does not enlarge the class).
\begin{samepage}
\begin{theorem}[縮約はクラスを拡大しない]\leavevmode\label{thm:33}
\(\Gamma_0,\ldots,\Gamma_s\) を \(\mathfrak M_0\) の成功する実行とし、\(M_*=M_s\)（(I) の第2の形では \(\operatorname{rad}(2ab)\mid M_*\)）および \(K_*=K_s\) とする。
このとき、完全な1ステップグラフ \(G(M_*,K_*)\) のすべての巡回的SCCは、第3.2節の位相検査に合格する。逆に、定理10（または補題11）のような1ステップ証明書は、(C) を使わない \(\mathfrak M_0\) の実行である。
\label{thm-end:33}
\end{theorem}
\end{samepage}

% Markdown AST block 243: Proof.
\begin{proof}\leavevmode\label{proof:33}
\(G(M_*,K_*)\) に、出力が最終的周期的でない無限ウォーク \(\omega\) が存在すると仮定する。その \((M_0,K_0)\) への射影は \(\Gamma_0\) のウォークである。(I) の第1の形では補題31による。第2の形では、補題29によりラベルが \(\mathcal C(a,b)\) に属し、\(R_0\) を法とする剰余が (E2) を満たすことによる。

% Markdown AST block 244: Para
補題30を帰納的に用いると、\(\omega\) の尾部はすべての \(\Gamma_t\) のウォークとなり、\(\Gamma_s=\varnothing\) のウォークにもなる。これは矛盾である。

% Markdown AST block 245: Para
したがって、\(G(M_*,K_*)\) のすべての無限ウォークは最終的周期的な出力をもつ。特に、各巡回的SCCの内部のすべての無限ウォークがそうである。定理21の (b)\(\Rightarrow\)(a) により、すべての巡回的SCCが認証される。逆は明らかである。\qedhere
\end{proof}

% Markdown AST block 246: Para
したがって、\(K\) と素数集合を固定したとき、素数を加える順序は、実行が成功し得るかどうかに影響せず、中間段階の大きさだけに影響する。「\(G(M_*,K_*)\) のウォークで、出力が最終的周期的でないものは存在しない」という結論自体には、定理21を用いていない。これを位相検査として言い換えるときにのみ、同定理を用いる。

% Markdown AST block 247: Theorem 34 (projection monotonicity).
\begin{samepage}
\begin{theorem}[射影の単調性]\leavevmode\label{thm:34}
操作 (P)、(R)、(L) を用いる1ステップの実行を \(W_0=V(M_0,K_0)\) から始める。\(U_t=R_{K_t}(W_t)\) とし、\(W_{t+1}\) は \(U_t\) の完全な拡大（補題9および11）であるとする。\(U_s=\varnothing\) ならば、
\[
R_{K_s}(V(M_s,K_s))=\varnothing:
\]
が成り立つ。すなわち、最終的なパラメータにおける完全なグラフは、1段階で認証される。
したがって、\(M_s\mid M'\) および \(K_s\mid K'\) ならば、
\[
R_{K'}(V(M',K'))=\varnothing
\]
も成り立つ。特に、対 \((s!,s!)\)、\(s=1,2,\ldots\) を調べれば、1ステップ証明書を見落とすことはない。
ただし、\((M',K')\) における証明書が示すのは、\(\gcd(Q_n,M')>1\) が無限回成り立つことだけであり、元の素数集合についての結論は回復されない。
\label{thm-end:34}
\end{theorem}
\end{samepage}

% Markdown AST block 248: Proof.
\begin{proof}\leavevmode\label{proof:34}
\[
\pi_t(R_{K_s}(V(M_s,K_s)))\subseteq U_t
\]
を \(t\) に関する帰納法で示す。ここで、\(\pi_t\) は \((M_t,K_t)\) への射影である。
\(t=0\) の場合は、\(W_0\) が頂点集合全体なので補題32を適用できる。

% Markdown AST block 249: Para
帰納段階では、\(v\) が最終グラフの認証されていない巡回的SCC \(C'\) に属するとする。帰納法の仮定により
\[
\pi_t(C')\subseteq U_t,
\]
であるから、
\[
\pi_{t+1}(C')\subseteq W_{t+1}.
\]
である。集合 \(\pi_{t+1}(C')\) は \(G_{K_{t+1}}(W_{t+1})\) 内で強連結であり（誘導部分グラフでは辺が保持される）、一つの巡回的SCC \(C\) に含まれる。もし \(C\) が認証されていれば、補題32により \(C'\) も認証される。よって \(C\subseteq U_{t+1}\) であり、\(\pi_{t+1}(v)\in U_{t+1}\) を得る。

% Markdown AST block 250: Para
\(t=s\) では、\(\pi_s\) は恒等写像である。

% Markdown AST block 251: Para
最後の主張は、\((M',K')\) から \((M_s,K_s)\) への射影に対し、\(W=V(M_s,K_s)\)、\(W'=V(M',K')\) として補題32を適用すれば従う。\qedhere
\end{proof}

% Markdown AST block 252: Para
これは相対的な完全性にとどまる。対 \((s!,s!)\) にわたる探索は、何らかのパラメータで証明書が存在すれば、それを見つける。しかし、得られる結論はより弱い \(\gcd(Q_n,s!)>1\) であり、探索の停止は主張していない。

% Markdown AST block 253: Proposition 35 (prime powers).
\begin{samepage}
\begin{proposition}[素数冪]\leavevmode\label{thm:35}
任意の \(M\) および \(K\) に対して、\(G(M,K)\) と \(G(\operatorname{rad}M,K)\) の無限出力語の集合は、付随するセル列も含めて一致する。
\label{thm-end:35}
\end{proposition}
\end{samepage}

% Markdown AST block 254: Proof.
\begin{proof}\leavevmode\label{proof:35}
一方の包含は補題31である。逆に、\((\bar q_n,j_n)_{n\ge0}\) を \(G(\operatorname{rad}M,K)\) の無限ウォークとし、その出力を \((c_n)\) とする。
各 \(p^e\parallel M\) に対し、単元 \(q_n\) を法 \(p^e\) で構成して、
\[
q_n\equiv\bar q_n\pmod p\qquad \text{and}\qquad bq_{n+1}\equiv aq_n+c_n\pmod{p^e}.
\]
を満たすようにする。

% Markdown AST block 255: Para
\(p\nmid b\) の場合は、\(q_0\) を \(\bar q_0\) の任意の持ち上げとし、
\[
q_{n+1}=b^{-1}(aq_n+c_n)\bmod p^e;
\]
とする。法 \(p\) では、これは初期値も漸化式も元の列と同じなので、\(q_n\equiv\bar q_n\) である。

% Markdown AST block 256: Para
\(p\mid b\) の場合（したがって \(p\nmid a\)）は、
\[
q_n=-\sum_{j=0}^{e-1}b^ja^{-j-1}c_{n+j}\bmod p^e.
\]
とおく。このとき
\[
\begin{aligned}
bq_{n+1}-aq_n&=-\sum_{j=1}^{e}b^ja^{-j}c_{n+j}+\sum_{j=0}^{e-1}b^ja^{-j}c_{n+j}\\
&=c_n-b^ea^{-e}c_{n+e}\\
&\equiv c_n\pmod{p^e},
\end{aligned}
\]
であり、法 \(p\) では
\[
q_n\equiv-a^{-1}c_n\equiv\bar q_n
\]
が成り立つ。これは、元のウォークが \(0\equiv a\bar q_n+c_n\pmod p\) を満たすからである。

% Markdown AST block 257: Para
中国剰余定理によってこれらを合わせればよい。(E3) はセルとラベルだけに関わる。\qedhere
\end{proof}

% Markdown AST block 258: Proposition 36 (primes dividing \(2ab\)).
\begin{samepage}
\begin{proposition}[\(2ab\) を割る素数]\leavevmode\label{thm:36}
\(M_0=\operatorname{rad}(2ab)\) および \(\gcd(R,M_0)=1\) とする。\(G(M_0R,K)\) と、(I) で定義した \(G'(R,K)\) の無限出力語の集合は、セル列も含めて一致する。
\label{thm-end:36}
\end{proposition}
\end{samepage}

% Markdown AST block 259: Proof.
\begin{proof}\leavevmode\label{proof:36}
\(G(M_0R,K)\) のウォークのラベルは補題29により \(\mathcal C(a,b)\) に属し、\(R\) を法とする剰余は (E2) を満たす。これにより \(G'(R,K)\) のウォークが得られる。

% Markdown AST block 260: Para
逆に、ウォーク \((\bar q_n,j_n)\) が \(G'(R,K)\) に与えられ、出力が \(c_n\in\mathcal C(a,b)\) であるとする。各素数 \(p\mid M_0\) について剰余を定める。
\(p\mid a\) のときは、\(q_0\) を任意の単元とし、
\[
q_{n+1}\equiv b^{-1}c_n
\]
とする（このとき \(bq_{n+1}\equiv c_n\equiv aq_n+c_n\)）。
\(p\mid b\) のときは、
\[
q_n\equiv-a^{-1}c_n
\]
とする（このとき \(bq_{n+1}-aq_n\equiv c_n\)）。
\(p=2\) で \(a,b\) が奇数のときは、\(q_n=1\) とする（このとき \(b\equiv a+c_n\pmod2\) である。\(c_n\) は偶数だからである）。\(p\nmid c_n\) により、これらはすべて単元である。中国剰余定理によって、これらと \(\bar q_n\) を合わせればよい。\qedhere
\end{proof}

% Markdown AST block 261: Para
\emph{有限に認証可能な底}のクラスを
\[
\begin{aligned}
\mathcal F_{\rm cert}&=\Bigl\{\tfrac ab:\ a>b\ge2,\ \gcd(a,b)=1,\\\
\exists K\ \exists S\text{ finite},\ p\nmid2ab\ (p&\in S):\\\
&\text{every cyclic SCC of }G'\bigl(\textstyle\prod_{p\in S}p,K\bigr)\\
&\text{ satisfies Theorem 21(c)}\Bigr\}.
\end{aligned}
\]
と定義する。

% Markdown AST block 262: Theorem 37 (normal form).
\begin{samepage}
\begin{theorem}[標準形]\leavevmode\label{thm:37}
既約な底 \(a/b\) について、次は同値である。

% Markdown AST block 263: OrderedList
\begin{enumerate}
\def\labelenumi{(\arabic{enumi})}
\item
  \(\mathfrak M_0\) の成功する実行が存在する。
\item
  ある法 \(M\) とあるセル数の列に対し、成功する1ステップの実行（操作 (I)、(P)、(R)、(L)、第1の形）が存在する。
\item
  \(K\) と、素数の有限集合 \(S\) が存在し、そのどの素数も \(2ab\) を割らず、\(G'(\prod_{p\in S}p,K)\) のすべての巡回的SCCが位相検査に合格する。
\item
  \(a/b\in\mathcal F_{\rm cert}\)。
\end{enumerate}
\label{thm-end:37}
\end{theorem}
\end{samepage}

% Markdown AST block 264: Proof.
\begin{proof}\leavevmode\label{proof:37}
(3)\(\Leftrightarrow\)(4) は定理21による。

% Markdown AST block 265: Para
(3)\(\Rightarrow\)(2)：命題36および定理21により、\(G(M_0R,K)\) のすべての巡回的SCCが認証される。これは \(s=0\) とした1ステップの実行である。

% Markdown AST block 266: Para
(2)\(\Rightarrow\)(1) は明らかである。

% Markdown AST block 267: Para
(1)\(\Rightarrow\)(3)：定理33により、\(G(M_*,K_*)\) のすべての巡回的SCCが認証される。
補題31および定理21により、同じことが \(G(M_0M_*,K_*)\) についても成り立つ。ここで \(M_0=\operatorname{rad}(2ab)\) である。実際、細かい側のグラフの各無限ウォークは、最終的周期的な出力をもつウォークに射影される。

% Markdown AST block 268: Para
命題35により、同じことが
\[
G(\operatorname{rad}(M_0M_*),K_*)=G(M_0R,K_*),
\]
についても成り立つ。ここで \(R\) は、\(M_*\) の素因数のうち \(2ab\) を割らないものの積である。命題36および定理21により、\(G'(R,K_*)\) についても成り立つ。\qedhere
\end{proof}

% Markdown AST block 269: Para
この同値性は、何らかの証明書の\emph{存在}についてのものであり、結論に現れる素数集合についてのものではない。
(2) の証明書は \(\gcd(Q_n,M)>1\) が無限回成り立つことを示す。一方、(3) の証明書は \(\gcd(Q_n,2abR)>1\) が無限回成り立つことを示し、この結論の素数集合は \(2ab\) を割るすべての素数を含む。
たとえば (1) では \(M=4290\) を使っており、これは \(7\) を含まないが、\(7/5\) の標準形では \(\{2,5,7\}\cup\{3,11,13\}\) を使う。
\(3/2,4/3,5/4,7/5,5/2\) が \(\mathcal F_{\rm cert}\) に属することは、第5節の証明書（表1、3、4、6）から従う。それ以外の底がこのクラスに属するとは主張しない。また、このクラスは \(a,b\) に関する単純な不等式や合同式で特徴付けられているわけではない。

% Markdown AST block 270: 8. Obstructions and the external-prime budget
\Needspace{6\baselineskip}
\section{障害と外部素数の積の大きさ}\label{sec:8}

% Markdown AST block 271: Para
本節の結果は、次の共通の形をもつ。与えられた法 \(M\) に対し、\emph{\(\mathfrak M_0\) の実行で、すべての法 \(M_t\) が \(M\) を割るものは成功しない}。
\(\mathfrak M_0\) に、両方向で擬軌道に関して健全な操作を加えた場合も同じである。これらは、手続きが停止しないことについての命題であって、
\[
\exists\,\xi>0\ \exists\,n_0\ \forall n\ge n_0:\ \gcd(\lfloor\xi r^n\rfloor,M)=1,
\]
を導くものでは\emph{ない}。また、合成数項に関する命題の反例でもない。
障害を示す証拠となるのは擬軌道であり、そこには剰余と小数部分を結び付ける整数列がない。補題41によれば、このような証拠の有限窓は、どれも真の軌道によって実現される。ただし、その量化の形は \(\forall N\ \exists\xi_N\) であり、\(\exists\xi\ \forall N\) に交換することはできない。
真の軌道だけに利用できる情報（\(Q_n\) の整数性、前向きの実現可能性、または第10節の再帰性の仮定）を用いる手続きは、これらの定理の対象ではない。

% Markdown AST block 272: Lemma 38 (type of the obstruction).
\begin{samepage}
\begin{lemma}[障害の種類]\leavevmode\label{thm:38}
\(M\ge2\) とする。右向きまたは左向きの \(M\)-擬軌道で、その繰り上がり列がその向きに最終的周期的でなく、(I) の第2の形を用いる場合には繰り上がりが \(\mathcal C(a,b)\) に属するものが存在するとする。
このとき、\(\mathfrak M_0\)（または擬軌道に関して健全な操作を加えた拡張）の実行で、すべての段階で \(M_t\mid M\) を満たすものは成功しない。
\label{thm-end:38}
\end{lemma}
\end{samepage}

% Markdown AST block 273: Proof.
\begin{proof}\leavevmode\label{proof:38}
この擬軌道から、任意の \(K\) に対して、同じ繰り上がりをもつ \(G(M,K)\) のウォークが得られる。その \((M_0,K_0)\) への射影は \(\Gamma_0\) のウォークである。(I) の第2の形では、仮定によりラベルが \(\mathcal C(a,b)\) に属し、\(R_0\) を法とする剰余は (E2) を満たす。
補題30を帰納的に適用すると、その尾部は各 \(\Gamma_t\) のウォークとなる。したがって \(\Gamma_s\ne\varnothing\) である。\qedhere
\end{proof}

% Markdown AST block 274: 8.1 Large numerators
\Needspace{6\baselineskip}
\subsection{分子が大きい場合}\label{sec:8.1}

% Markdown AST block 275: Theorem 39.
\begin{samepage}
\begin{theorem}\leavevmode\label{thm:39}
\(M\ge2\)、\(R=\operatorname{rad}M\) とし、
\[
a\ge2R.
\]
を仮定する。このとき、\(\mathfrak M_0\)（または擬軌道に関して健全な操作を加えた拡張）の実行で、すべての段階で \(M_t\mid M\) を満たし、\(\Gamma_0\) が (I) の形であるものは成功しない。
したがって、実行の成功には
\[
a<2\operatorname{rad}(M).
\tag{27}\label{eq:27}
\]
が必要である。
\label{thm-end:39}
\end{theorem}
\end{samepage}

% Markdown AST block 276: Proof.
\begin{proof}\leavevmode\label{proof:39}
\emph{ステップ1：左向きの非周期的な繰り上がり列。} \(\theta_0=0\) とおく。\(\theta_{n+1}\in[0,1)\) が与えられたとき、
\[
\begin{aligned}
c_n&\in[-b\theta_{n+1},\,a-b\theta_{n+1})\cap(b-a+R\mathbb Z),\\
\theta_n&=\frac{b\theta_{n+1}+c_n}{a}.
\end{aligned}
\]
と選ぶ。この半開区間の長さは \(a\ge2R\) なので、この等差数列の項を少なくとも \(\lfloor a/R\rfloor\ge2\) 個含む。このとき
\[
0\le b\theta_{n+1}+c_n<a
\]
より \(\theta_n\in[0,1)\) が得られ、恒等式
\[
a\theta_n-b\theta_{n+1}=c_n
\]
が厳密に成り立つ。また、\(-b<c_n<a\) より \(1-b\le c_n\le a-1\) である。
各ステップで、順序を付けた二つの選択肢だけを使うと、写像 \(\{0,1\}^{\mathbb N}\to\Sigma^{\mathbb Z_{\le0}}\) は単射となる。実際、二つの選択列が初めて異なる添字では、\(\theta_{n+1}\) の値が一致し、\(c_n\) が異なる。したがって、非可算個の列が得られる。
一方、列 \((c_n)_{n\le0}\) が\emph{左向きに最終的周期的}であるとは、ある \(N\le0\) と \(L\ge1\) が存在し、\(c_{n-L}=c_n\) がすべての \(n\le N\) で成り立つことをいう。このような列は \((N,L,c_{N-L+1},\ldots,c_0)\) によって定まるため、可算個しかない。
左向きに最終的周期的でない列を一つ選ぶ。法 \(R\) では、\(q_n=1\) は \(bq_{n+1}\equiv aq_n+c_n\) を満たす。これは \(c_n\equiv b-a\pmod R\) だからである。

% Markdown AST block 277: Para
\emph{ステップ2：素数冪を法とする単元。} 各 \(p^k\parallel M\) に対して、\(q_n\equiv1\pmod p\) を
\[
bq_{n+1}\equiv aq_n+c_n\pmod{p^k}.
\]
が成り立つように定める。
\(p\nmid a\) の場合は、\(q_0=1\) とし、
\[
q_n=a^{-1}(bq_{n+1}-c_n)\bmod p^k
\]
を \(n<0\) に対して用いる。\(c_n\equiv b-a\pmod p\) なので、\(q_{n+1}\equiv1\) から
\[
q_n\equiv a^{-1}(b-(b-a))=1.
\]
が従う。

% Markdown AST block 278: Para
\(p\mid a\) の場合（したがって \(b\) は法 \(p^k\) で可逆）は、\(L\) を \(p^k\mid a^L\) となるように選び、
\[
q_n=\sum_{j=0}^{L-1}a^jb^{-j-1}c_{n-1-j}\bmod p^k.
\]
とおく。このとき
\[
\begin{aligned}
bq_{n+1}-aq_n&=\sum_{j=0}^{L-1}a^jb^{-j}c_{n-j}-\sum_{j=1}^{L}a^jb^{-j}c_{n-j}\\
&=c_n-a^Lb^{-L}c_{n-L}\\
&\equiv c_n\pmod{p^k},
\end{aligned}
\]
であり、法 \(p\) では
\[
q_n\equiv b^{-1}c_{n-1}\equiv b^{-1}(b-a)\equiv1.
\]
となる。用いた添字はすべて \(\le-1\) である。

% Markdown AST block 279: Para
中国剰余定理により、
\(\gcd(q_n,M)=1\)
および
\[
bq_{n+1}\equiv aq_n+c_n\pmod M.
\]
を得る。

% Markdown AST block 280: Para
\emph{ステップ3。} \((q_n,\theta_n,c_n)_{n\le0}\) は左向きの \(M\)-擬軌道であり、その繰り上がり列は左向きに最終的周期的でない。
\(\operatorname{rad}(2ab)\mid M\) の場合、合同式 \(c_n\equiv b-a\pmod p\) がすべての \(p\mid2ab\) について成り立つことから、\(c_n\in\mathcal C(a,b)\) を得る。実際、\(p\mid a\) ならば \(c_n\equiv b\)、\(p\mid b\) ならば \(c_n\equiv-a\) であり、\(p=2\) で \(a,b\) が奇数ならば \(c_n\) は偶数である。
よって補題38を適用できる。\qedhere
\end{proof}

% Markdown AST block 281: Para
初期グラフに関する仮定は不可欠である。真の軌道の性質から導かれる別の健全な初期制限から始める手続きは \(\mathfrak M_0\) に属さず、この証拠となる擬軌道がその制限を通過するとは限らない。

% Markdown AST block 282: Corollary 40.
\begin{samepage}
\begin{corollary}\leavevmode\label{thm:40}
(a) 有限の素数集合 \(\mathcal P\) を固定する。既約な底 \(a/b\) のうち、\(\mathfrak M_0\) の成功する実行で、その法が \(\mathcal P\) の素数だけからなるもの（指数は任意）が存在する底は、有限個である。実際、
\[
a<2\prod_{p\in\mathcal P}p
\]
および \(b<a\) が必要である。

% Markdown AST block 283: OrderedList
\begin{enumerate}
\def\labelenumi{(\alph{enumi})}
\setcounter{enumi}{1}
\tightlist
\item
  底 \(r_t=(2t+7)/(2t+5)\)、\(t\ge0\) は既約である。\(r_t\) に対し、最終的な法が \(M_t\) である実行が成功するには、
  \[
  \operatorname{rad}(M_t)>(2t+7)/2,
  \]
  すなわち
  \[
  \operatorname{rad}(M_t)\ge t+4.
  \]
  が必要である。(1) の素数 \(\{2,3,5,11,13\}\) だけを使う実行は、\(t\ge4287\) では成功しない。
  同じことは、\((14k+1)/(10k+1)\)（\(7/5\) に収束する）についても、\(k\ge613\) のとき成り立つ。
\end{enumerate}
\label{thm-end:40}
\end{corollary}
\end{samepage}

% Markdown AST block 284: Proof.
\begin{proof}\leavevmode\label{proof:40}
(a) は (27) である。(b) 各対は互いに素である。第1の族では、両者の差は \(2\) であり、両者とも奇数である。第2の族では、公約数は
\[
5(14k+1)-7(10k+1)=-2
\]
を割り、やはり両者とも奇数である。閾値は、次の同値関係から得られる。
\[
2t+7\ge8580\iff t\ge4287
\]
（\(2\cdot4286+7=8579\) に注意）。また、
\[
14k+1\ge8580\iff k\ge613.\qedhere
\]
\end{proof}

% Markdown AST block 285: Para
(b) は、\(7/5\) の証明書を同じ素数集合のまま近くの底へ移せないことを述べている。これらの底における合成数項については何も述べていない。

% Markdown AST block 286: 8.2 Finite windows are realized by true orbits
\Needspace{6\baselineskip}
\subsection{有限窓は真の軌道によって実現される}\label{sec:8.2}

% Markdown AST block 287: Lemma 41.
\begin{samepage}
\begin{lemma}\leavevmode\label{thm:41}
整数 \(q_i\) は \(\gcd(q_i,M)=1\) を満たすとする。実数 \(\theta_i\in[0,1)\)（\(0\le i\le L\)）および整数 \(c_i\)（\(0\le i<L\)）が
\[
\begin{aligned}
&a\theta_i-b\theta_{i+1}=c_i\qquad\text{and}\\
&bq_{i+1}\equiv aq_i+c_i\pmod M.
\end{aligned}
\]
を満たすならば、任意の \(T_0\) に対して、ある \(\xi>T_0\) が存在し、\(\lfloor\xi r^i\rfloor\equiv q_i\pmod M\) および \(\{\xi r^i\}=\theta_i\) が \(0\le i\le L\) で成り立つ。
\label{thm-end:41}
\end{lemma}
\end{samepage}

% Markdown AST block 288: Proof.
\begin{proof}\leavevmode\label{proof:41}
\(Q_i=q_i+Mt_i\) と書き、整数 \(t_i\) を求める。また、
\[
d_i=(aq_i+c_i-bq_{i+1})/M\in\mathbb Z.
\]
とおく。

% Markdown AST block 289: Para
条件 \(bQ_{i+1}=aQ_i+c_i\) は
\[
bt_{i+1}=at_i+d_i.
\]
と書ける。\(S_0=0\)、\(S_{j+1}=aS_j+b^jd_j\) とおけば、帰納法により
\[
b^jt_j=a^jt_0+S_j
\]
である。

% Markdown AST block 290: Para
\(\gcd(a,b^L)=1\) なので、
\[
t_0\equiv-a^{-L}S_L\pmod{b^L}.
\]
と選ぶ。\(j\le L\) に対し、
\[
S_L\equiv a^{L-j}S_j\pmod{b^j},
\]
であるから、
\[
a^jt_0+S_j\equiv-a^ja^{-L}a^{L-j}S_j+S_j\equiv0\pmod{b^j},
\]
となる。したがって、
\[
t_j=(a^jt_0+S_j)/b^j
\]
は整数であり、\(bt_{j+1}=at_j+d_j\) は有理数の恒等式として成り立ち、したがって整数の恒等式としても成り立つ。

% Markdown AST block 291: Para
\(b^LT\) を \(t_0\) に加えると、\(a^jb^{L-j}T\) が \(t_j\) に加わる。\(T\) を大きくすれば、すべての \(Q_i\) を正にでき、しかも望むだけ大きくできる。

% Markdown AST block 292: Para
最後に、\(x_i=Q_i+\theta_i\) は
\[
bx_{i+1}=aQ_i+c_i+b\theta_{i+1}=ax_i,
\]
を満たすので、\(\xi=x_0\) とすれば \(x_i=\xi r^i\) および \(\lfloor x_i\rfloor=Q_i\) を得る。\qedhere
\end{proof}

% Markdown AST block 293: 8.3 The two-carry obstruction
\Needspace{6\baselineskip}
\subsection{二つの繰り上がりによる障害}\label{sec:8.3}

% Markdown AST block 294: Para
\(\Delta=a-b\) とおく（位相の最大公約数 \(d\) と混同しないこと）。

% Markdown AST block 295: Theorem 42.
\begin{samepage}
\begin{theorem}\leavevmode\label{thm:42}
\(M\ge2\) とし、\(R\) を、\(M\) の素因数のうち \(2ab\) を割らないものの積とする（\(R=1\) も許す）。整数 \(m\ge1\) が存在し、
\[
h=2Rm\le a-1,\qquad c=\Delta-h,\qquad \gcd(c,ab)=1 .
\tag{28}\label{eq:28}
\]
を満たすと仮定する。このとき、\(\mathfrak M_0\)（または擬軌道に関して健全な操作を加えた拡張）の実行で、すべての段階で \(M_t\mid M\) を満たすものは成功しない。
\label{thm-end:42}
\end{theorem}
\end{samepage}

% Markdown AST block 296: Proof.
\begin{proof}\leavevmode\label{proof:42}
\emph{繰り上がりの条件。} 次が成り立つ。
\[
\begin{aligned}
&1-b\le c<\Delta\le a-1,\qquad c\equiv\Delta\pmod2\quad(h\text{ is even}),\\
&c\equiv\Delta\pmod R,\\
&\text{and}\quad\gcd(\Delta,ab)=\gcd(a-b,a)\gcd(a-b,b)=1;
\end{aligned}
\]
したがって、\(c\) と \(\Delta\) はどちらも \(\mathcal C(a,b)\) に属する。（\(c=0\) は、\(\gcd(c,ab)=1\) によって除外される。）

% Markdown AST block 297: Para
\emph{不動点と不変区間。} 次のようにおく。
\[
\begin{aligned}
\psi_\Delta(x)&=(bx+\Delta)/a=1+\tfrac ba(x-1)\qquad \text{and}\\
\psi_c(x)&=(bx+c)/a=\psi_\Delta(x)-h/a;
\end{aligned}
\]
これらは逆向き写像 \(\theta_{n+1}\mapsto\theta_n\) であり、それぞれ繰り上がり \(\Delta\) および \(c\) に対応する。

% Markdown AST block 298: Para
語 \(W_L=\Delta^Lc\)（前向きの時間順序）に対する逆向きの合成は
\[
\begin{aligned}
\Phi_L&=\psi_\Delta^L\circ\psi_c,\\
\Phi_L(x)&=1+\Bigl(\tfrac ba\Bigr)^{L+1}(x-1)-\tfrac ha\Bigl(\tfrac ba\Bigr)^L,
\end{aligned}
\]
であり、不動点は
\[
x_L=1-hb^L/(a^{L+1}-b^{L+1}).
\]
である。\(a>b\) より、\(x_L<x_{L+1}\) および \(x_L\uparrow1\) が成り立ち、常に \(x_L<1\) である。\(c\ge1-b\) より \(\max(0,-c/b)<1\) であるから、最小の \(L\) で
\[
\max(0,-c/b)<x_L;
\]
を満たすものが存在する。\(I=[x_L,x_{L+1}]\) とおく。
写像 \(\Phi_L,\Phi_{L+1}\) は、傾きが \((0,1)\) に属する単調増加なアフィン写像であり、その不動点は、それぞれ \(I\) の左端点と右端点である。したがって、どちらも \(I\) を \(I\) に写す。

% Markdown AST block 299: Para
\(x\in I\) に対し、\(\psi_c(x)>0\) が成り立つ。これは \(x>-c/b\) による。また、
\[
\psi_c(x)<\psi_\Delta(x)<1
\]
が成り立つ。これは \(c<\Delta\) および \(x<1\) による。
\(\psi_\Delta\) は \((0,1)\) を \((\Delta/a,1)\) に写す。
したがって、\(W_L\) または \(W_{L+1}\) を \(I\) から逆向きに適用したときの中間値は、すべて \((0,1)\) に属する。

% Markdown AST block 300: Para
\emph{無限連結。} \(U=W_L\)、\(V=W_{L+1}\) とし、\(B_1B_2\cdots\) を \(\{U,V\}\) のブロックからなる任意の列とする。ブロック \(i\) は時刻 \(n_i\)（\(n_1=0\)）に始まるものとする。
\(k\ge i\) に対し、集合
\[
I_i^{(k)}=\Phi_{B_i}\circ\cdots\circ\Phi_{B_k}(I)
\]
は \(I\) の閉部分区間であり、\(k\) に関して減少する。その長さは高々
\[
(b/a)^{(L+1)(k-i+1)}|I|\to0;
\]
であるから、それらの共通部分は一点 \(\theta_{n_i}\in I\) となる。構成により、
\[
\Phi_{B_i}(\theta_{n_{i+1}})=\theta_{n_i}
\]
が成り立つ。

% Markdown AST block 301: Para
各ブロックの内部の値 \(\theta_n\) は、\(\theta_{n_{i+1}}\) からの逆向き漸化式によって定める。その合成は \(\Phi_{B_i}\) である。
したがって \(0<\theta_n<1\) であり、
\[
a\theta_n-b\theta_{n+1}=c_n
\]
がすべての \(n\ge0\) について厳密に成り立つ。

% Markdown AST block 302: Para
\emph{非周期的符号化。} 語 \(UV\) と \(VU\) は同じ長さ \(2L+3\) をもち、位置 \(L\) で異なる（文字はそれぞれ \(c\) と \(\Delta\)）。
最終的周期的でない二進列（たとえば \(1\,0\,1\,00\,1\,000\cdots\)）を、\(0\mapsto UV\)、\(1\mapsto VU\) によって符号化する。
繰り上がり列が最終的に周期 \(P\) をもつならば、\(P(2L+3)\) も周期となり、ブロックを復号することで二進列も最終的周期的になってしまう。

% Markdown AST block 303: Para
\emph{法 \(M\) における単元剰余。} \(M\) の各素因数は \(2abR\) を割る。
\(p^e\parallel M\) に対し、\(p\mid a\) の場合は、\(q_0=1\) とし、
\[
q_{n+1}=b^{-1}(aq_n+c_n)\bmod p^e;
\]
とおく。法 \(p\) では
\[
q_{n+1}\equiv b^{-1}c_n\ne0
\]
である。これは \(c_n\in\{c,\Delta\}\) が \(a\) と互いに素だからである。

% Markdown AST block 304: Para
\(p\mid b\) の場合は、
\[
q_n=-\sum_{j=0}^{e-1}b^ja^{-j-1}c_{n+j}\bmod p^e;
\]
とおく。命題35の証明と同様に、
\[
\begin{aligned}
&bq_{n+1}-aq_n\equiv c_n\pmod{p^e},\qquad\text{and}\\
&q_n\equiv-a^{-1}c_n\ne0\pmod p.
\end{aligned}
\]
が成り立つ。

% Markdown AST block 305: Para
\(p\mid R\) の場合は、\(q_0\equiv-1\pmod{p^e}\) とし、
\[
q_{n+1}=b^{-1}(aq_n+c_n)\bmod p^e
\]
とおく（\(b\) は法 \(p^e\) で単元である。\(p\nmid2ab\) だからである）。法 \(p\) に還元すると、\(c_n\equiv\Delta\) により
\[
q_{n+1}\equiv b^{-1}(-a+a-b)=-1,
\]
となるので、各 \(q_n\) は単元である。

% Markdown AST block 306: Para
\(p=2\) で \(a,b\) が奇数の場合は、\(c_n\) は偶数である。\(q_0\) を奇数に取れば、
\[
q_{n+1}\equiv q_n+c_n\equiv q_n\pmod2.
\]
となる。中国剰余定理によってこれらを合わせる。

% Markdown AST block 307: Para
こうして、右向きの \(M\)-擬軌道で、繰り上がりが \(\mathcal C(a,b)\) に属し、その列が最終的周期的でないものを得る。よって補題38を適用できる。\qedhere
\end{proof}

% Markdown AST block 308: Corollary 43 (external primes are necessary).
\begin{samepage}
\begin{corollary}[外部素数の必要性]\leavevmode\label{thm:43}
任意の既約な底 \(a>b\ge2\) に対して、\(\mathfrak M_0\) の実行で、法に現れる素数が \(2ab\) の素因数だけであるものは成功しない。
\label{thm-end:43}
\end{corollary}
\end{samepage}

% Markdown AST block 309: Proof.
\begin{proof}\leavevmode\label{proof:43}
定理42で \(R=1\) とする。\(\Delta\) が奇数で \(\ge3\) ならば、
\[
c=1,\qquad h=\Delta-1\ge2.
\]
とおく。

% Markdown AST block 310: Para
\(\Delta\) が偶数で \(\ge4\) ならば、\(a,b\) は奇数であり、\(c=2\) は \(ab\) と互いに素である。このとき
\[
h=\Delta-2\ge2.
\]
とおく。

% Markdown AST block 311: Para
\(\Delta=1\) ならば、
\[
c=-1\ge1-b,\qquad h=2\le a-1
\]
とおく（\(a\ge3\) による）。

% Markdown AST block 312: Para
\(\Delta=2\) ならば、\(a,b\) は奇数であり、\(b\ge3\)、\(a\ge5\) なので、
\[
c=-2\ge1-b,\qquad h=4\le a-1.
\]
とおく。

% Markdown AST block 313: Para
どの場合も \(h\le a-1\) は正の偶数であるから、\(m=h/2\) として (28) が成り立つ。\qedhere
\end{proof}

% Markdown AST block 314: Para
\(7/5\) で \(M=4290\) の場合には、\(R=3\cdot11\cdot13=429\) であり、\(h=858m\le6\) は不可能である。これは定理1(i)と整合する。

% Markdown AST block 315: 8.4 The external-prime budget
\Needspace{6\baselineskip}
\subsection{外部素数の積の大きさ}\label{sec:8.4}

% Markdown AST block 316: Para
\(J(N)\) をJacobsthal関数とする。すなわち、正整数 \(J\) のうち、任意の連続する \(J\) 個の整数の中に \(N\) と互いに素なものが含まれるという性質をもつ最小のものである。
\(J(1)=1\) とする。同値な定義として、\(J(N)\) は、\(N\) と互いに素な整数を順に並べたときの最大の間隔である。
\[
N_0=\prod\{p\mid ab:p\text{ odd}\}
\]
とおき、その素因数の個数を \(k=\omega(N_0)\) とする。\(a>b\ge2\) は互いに素なので、\(N_0\ge3\) である。

% Markdown AST block 317: Theorem 44.
\begin{samepage}
\begin{theorem}\leavevmode\label{thm:44}
\(R\) を定理42と同じものとする。\(\mathfrak M_0\) の実行で、すべての段階で \(M_t\mid M\) を満たすものが成功するならば、
\[
a\le2R\,J(N_0).
\tag{29}\label{eq:29}
\]
である。
さらに、任意の奇数かつ平方因子をもたない \(N\) で \(\omega(N)=k\) を満たすものについて、
\[
J(N)\le3^k;
\]
である。したがって、
\[
R\ge\frac{a}{2\cdot3^{k}},\qquad 2^kk!\le N_0<a^2,
\tag{30}\label{eq:30}
\]
となる。よって \(3^k=a^{o(1)}\) であり、
\[
\liminf\log R/\log a\ge1
\]
が、\(a\to\infty\) となる成功した証明書の任意の列に沿って成り立つ。
\(R\) に使える素数が、積を \(R_0\) とする固定の有限集合に制限されているならば、
\[
a\le2R_0\,3^{k_0-1},
\]
である。ここで \(k_0\ge4\) は
\[
2^{k_0}k_0!>4R_0^2\,9^{k_0};
\]
を満たす任意の整数である。特に、その素数集合で成功する実行をもつ既約な底は有限個しかない。
\label{thm-end:44}
\end{theorem}
\end{samepage}

% Markdown AST block 318: Proof.
\begin{proof}\leavevmode\label{proof:44}
定理42により、実行が成功するならば、
\[
\gcd(\Delta-2Rm,ab)>1
\]
が、すべての
\[
1\le m\le H:=\lfloor(a-1)/(2R)\rfloor.
\]
に対して成り立つ。

% Markdown AST block 319: Para
\(2R\) は法 \(N_0\) で可逆なので、
\[
u\equiv\Delta(2R)^{-1}\pmod{N_0};
\]
とおく。奇素数 \(p\mid ab\) に対して、\(p\mid\Delta-2Rm\) と \(m\equiv u\pmod p\) は同値である。
素数 \(2\) が \(\Delta-2Rm\) を割ることは、\(2\mid ab\) の場合には決して起こらず（このとき \(\Delta\) は奇数）、それ以外の場合には、その素数は \(ab\) の素因数ではない。

% Markdown AST block 320: Para
したがって、
\[
\gcd(\Delta-2Rm,ab)=1\iff\gcd(m-u,N_0)=1.
\]
が成り立つ。

% Markdown AST block 321: Para
\(H\ge J(N_0)\) ならば、連続する \(H\) 個の整数 \(1-u,\ldots,H-u\) の中に \(N_0\) と互いに素なものが存在し、障害を示す \(m\) が得られる。
よって、実行の成功には \(H\le J(N_0)-1\)、すなわち
\[
(a-1)/(2R)<J(N_0),
\]
が必要である。これは (29) である。

% Markdown AST block 322: Para
\(J\) の上界を示す。連続する \(H\) 個の整数からなる窓の中で、\(N\) と互いに素なものの個数は
\[
\sum_{t\mid N}\mu(t)\,\#\{x:t\mid x\},
\qquad\text{and}\qquad
\#\{x:t\mid x\}=H/t+\epsilon_t
\]
である。ここで \(|\epsilon_t|<1\)、\(\epsilon_1=0\) なので、その個数は少なくとも
\[
H\varphi(N)/N-(2^k-1).
\]
である。

% Markdown AST block 323: Para
\(N\) は奇数であるから、
\[
\varphi(N)/N=\prod(1-1/p)\ge(2/3)^k,
\]
である。\(H=3^k\) とすれば、この個数は \(\ge1\) となる。

% Markdown AST block 324: Para
(30) の第1の不等式は (29) から、第2の不等式は
\[
N_0\ge\prod_{i\le k}(2i+1)>2^kk!\qquad \text{and}\qquad N_0\le ab<a^2.
\]
から従う。

% Markdown AST block 325: Para
\[
k!<a^2\qquad \text{and}\qquad k!\ge(k/e)^k
\]
から
\[
k\log k\ll\log a,
\]
を得る。したがって \(k=O(\log a/\log\log a)\) であり、
\[
3^k=a^{O(1/\log\log a)}.
\]
である。

% Markdown AST block 326: Para
有効な上界については、比 \(2^kk!/(4R_0^2\,9^k)\) は、次の添字に進むごとに \(2(k+1)/9>1\) 倍になる（\(k\ge4\) の場合）。したがって、この比は \(1\) を、すべての \(k\ge k_0\) に対して超える。

% Markdown AST block 327: Para
\(k\ge k_0\) ならば、
\[
a^2>N_0\ge2^kk!>4R_0^2\,9^k\ge(2R_0\,3^k)^2\ge a^2
\]
となる。ここでは (29) および \(J(N_0)\le3^k\) を用いた。これは矛盾であるから、\(k\le k_0-1\) であり、次を得る。
\[
a\le2R_0\,3^{k_0-1}.\qedhere
\]
\end{proof}

% Markdown AST block 328: Para
条件 (29) は必要条件であって、十分条件ではない。すべての \(m\le H\) に対応する障害を排除しても、証明書が得られるとは限らない。

% Markdown AST block 329: Corollary 45.
\begin{samepage}
\begin{corollary}\leavevmode\label{thm:45}
(a) \(a=p\) が奇素数で \(b=2^s<p\) ならば、\(N_0=p\)、\(J(p)=2\) であり、実行の成功には
\[
R\ge p/4.
\]
が必要である。
(b) \(a,b\) が相異なる奇素数ならば、\(J(ab)=3\) であり、実行の成功には
\[
R\ge a/6.
\]
が必要である。
\label{thm-end:45}
\end{corollary}
\end{samepage}

% Markdown AST block 330: Proof.
\begin{proof}\leavevmode\label{proof:45}
連続する二つの整数のうち、\(p\) で割り切れるものは高々一つなので、\(J(p)=2\) である。
連続する三つの整数のうち、\(a\) で割り切れるものは高々一つ、\(b\) で割り切れるものも高々一つなので、\(J(ab)\le3\) である。一方、中国剰余定理により、連続する整数で \(x\equiv0\pmod a\)、\(x+1\equiv0\pmod b\) を満たすものが存在するから、\(J(ab)=3\) である。(29) を適用すればよい。\qedhere
\end{proof}

% Markdown AST block 331: Para
これらは、挙げた底に対する任意の証明書が満たすべき制約であり、その底の合成数項に関する命題ではない。

% Markdown AST block 332: Remark (Iwaniec’s bound).
\begin{remark}[Iwaniecの上界]
\(C(k)\) を、任意に選んだ \(k\) 個の素数の少なくとも一つで割り切れる整数が、連続して並び得る最大の個数とする。Iwaniecは
\[
C(k)\ll k^2\log^2k
\]
を証明した（\hyperref[ref-iw78]{IW78, 系, p.226}）。
\hyperref[ref-cw13]{CW13, §1}も参照されたい。そこでは同じ結果を、\(g(n)\le X(k\log k)^2\)（\(k=\omega(n)\)）と書いている。
\(J(N)\le C(\omega(N))+1\) は \(\omega(N)\ge2\) のときに成り立ち、\(J(N)=2\) は \(\omega(N)=1\) のときに成り立つので、
\[
J(N_0)\ll(k\log2k)^2
\]
を得る。ここで暗黙の定数は絶対定数である。(30) の \(k\log2k\ll\log a\) と合わせると、条件 (29) から
\[
R\gg\frac{a}{(\log a)^2}.
\]
が従う。
これは (30) を定量的に改善する。この評価は外部の既知結果を用いたものであり、本論文の結果ではない。定数は絶対定数だが、ここではその値を指定しない。また、ここまでのどの定理も、この評価には依存しない。具体的な有限入力に対しては、(29) および定理42の最大公約数検査を直接用いる。
\end{remark}

% Markdown AST block 333: Para
底 \(101/2\) および \(101/100\) に対する具体的な証明書、残り続ける帰還語を除去するための素数の選択、ならびに語全体の区間検査を、付録E.4に示す。

% Markdown AST block 334: 9. Verification and reproducibility
\Needspace{6\baselineskip}
\section{検証と再現性}\label{sec:9}

% BEGIN PUBLIC ADDITION availability
論文ソース、計算コード、Lean形式証明、照合用データ、および再現手順は、次の公開リポジトリから入手できる。

\url{https://github.com/ixixi/rational-floor-certificates}

収録内容はリポジトリのREADMEに、再現手順と必要な環境の構築方法はREPRODUCE.mdおよびENVIRONMENT.mdに記載している。主な実行コマンドは付録Aに示す。
% END PUBLIC ADDITION availability

% Markdown AST block 335: 9.1 What the finite computations establish
\subsection{有限計算が確立すること}\label{sec:9.1}

% Markdown AST block 336: Para
第5節の各表および表7の数値は概要であり、証明書そのものではない。有限計算によって確立され、二つの実装によって要素ごとに検査された数学的な要件は、次のとおりである。

% Markdown AST block 337: OrderedList
\begin{enumerate}
\def\labelenumi{\arabic{enumi}.}
\tightlist
\item
  \emph{初期被覆。} 初期頂点集合は、許容状態の全体である（表1と表7では \(V(M,K_0)\) 全体、表4と表6では \(K\) 個のセルすべて）。
\item
  \emph{辺集合の完全性。} (E1)--(E3)、または定理19の初期不等式を満たすすべての辺が含まれる。各辺の持ち上げは、すべての始点剰余に対して計算される。
\item
  \emph{すべての中間時刻における剰余。} 持ち上げでは、\(\ell+1\) 個の剰余、すなわち語長 \(\ell\) の語に沿った終点を含むすべての剰余が非零である。
\item
  \emph{全文字に対する位相条件。} 認証された各成分について、(19) がすべての内部辺のすべての文字位置で成り立ち、各位相に対応するラベルは一つである。
\item
  \emph{順位の単調性。} 表7では、(Q-rank) および (Q-cyc) の不等式が異なるブロック間のすべての辺で成り立ち、上界 \(\rho\le65\)、\(\eta\le6\)、\(d_B\le13\) が成り立つ。
\item
  \emph{保持。} 認証されていない巡回的成分のすべての頂点と内部辺が保持され、それ以外は保持されない。
\item
  \emph{出力を保存する縮約。} マクロ頂点は、出次数が2以上の頂点とアンカーの全体にちょうど一致する。出次数1の頂点からなる各鎖はマクロ頂点で終わり、マクロ語は鎖上の語の連結である。
\item
  \emph{最終的な空性。} 最後の保持集合（表1）または最後の出力グラフ（表4と表6）は空である。表7では、認証されていない巡回的ブロックは存在しない。
\end{enumerate}

% Markdown AST block 338: Para
強連結成分そのものは、十分条件によって検証する。提案されたブロックは空でなく、互いに交わらず、すべての頂点を覆う。各ブロック内では前向きにも後ろ向きにも到達可能であり、異なるブロック間のすべての辺に対して厳密に単調な整数順位が存在する（商グラフが非巡回的である）。これらの条件によってSCC分割が定まる。
件数やハッシュの一致も記録するが、それを判定基準とはしない。判定基準は、完全なデータの要素ごとの比較である。

% Markdown AST block 339: 9.2 The one-step computation for 7/5
\Needspace{6\baselineskip}
\subsection{7/5の1ステップ計算}\label{sec:9.2}

% Markdown AST block 340: Para
表1--3の1ステップ計算は、Pythonの任意精度整数演算を使う二つの実装で行った。いずれも標準ライブラリだけを用い、証明用の計算では浮動小数点を使用していない。
実装Aは (12) によって到着セルを列挙し、反復的なKosaraju走査でSCCを計算する。実装Bはすべての \(k\in\{0,\ldots,K-1\}\) に対して (E3) を直接検査し、得られたリストを \((j,e)\) ごとにキャッシュする。SCCはTarjanのアルゴリズムの反復形式で計算する（\hyperref[ref-t]{T, §4}）。
両実装とも、(E2) を解くために、すべての単元剰余を \(bz\bmod M\) によって分類する。
Bは数学的仕様と共通の出力スキーマから開発され、Aおよび提供された基準プログラムを読んだりインポートしたりしていない。その設計、ソース、出力は、Aとの比較を行う前に固定した。

% Markdown AST block 341: Para
検証では、すべての頂点とラベル付き辺を検査する。最短距離、すべての辺の差分、位相ごとのラベル集合、認証フラグ、および原始語を再計算する。また、保持集合が認証されていない巡回的成分の和集合に正確に一致することと、(16) のすべての子セルが次段階に存在することを検査する。
位相検査に不合格となる場合には、出発位相が同じでラベルが異なる二つの内部辺が、検査可能な証拠を与える。
主要な4段階と表3の三つの初段を合わせると、比較の対象は97772頂点、168474本のラベル付き辺、および94006個のSCCであり、科学的データの不一致は見つからなかった。
記録された実行では、Linux上でAにCPython 3.11.13、BにCPython 3.10.12を用い、各ケースに実経過時間600秒、アドレス空間4 GiBの上限を設けた。上限に達した場合は \texttt{inconclusive} と記録し、反例や最小性の根拠とはしない。採用したケースはすべて完了した。
プログラムは \nolinkurl{supplement/computations/implementation_a/} および \nolinkurl{supplement/computations/implementation_b/} に同梱している。
公開版の検証手順と全要素比較は \href{supplement/COMPUTATION.md}{計算ガイド} および \href{REPRODUCE.md}{再現ガイド} に記載している。

% Markdown AST block 342: 9.3 The word-labelled computations
\Needspace{6\baselineskip}
\subsection{語ラベル付きグラフの計算}\label{sec:9.3}

% Markdown AST block 343: Para
表4、6、7の計算は、科学的中核を共有しない二つの実装によって行った。

% Markdown AST block 344: Para
\emph{実装A}（\href{supplement/COMPUTATION.md}{\nolinkurl{word_a}}）は、本論文の基礎となる研究ノートに同梱されたプログラムの、来歴を追跡できる作業用コピーである。構成処理はPythonから駆動されるC++プログラムで行い、表6の最初の5段階にはPythonによる第3の検査も備える。これに、資源上限を固定し、入力と出力を記録し、グラフを正規化されたテキスト形式で書き出すドライバを加えている。
提供されたプログラムには「A」と「B」と名付けられた二つの構成処理が含まれる。これらは起源、ファイル形式、規約を共有しているため、両者の一致は一つの系統内での整合性確認であり、次に述べる独立検証ではない。
科学的中核は公開ソースにも保持している。来歴と配布のための変更は \href{supplement/COMPUTATION.md}{計算ガイド} に記載し、\href{MANIFEST.json}{ソースマニフェスト} で各同梱ソースのSHA-256を示す。

% Markdown AST block 345: Para
\emph{実装B}（\href{supplement/COMPUTATION.md}{\nolinkurl{word_b}}）は、別のプログラマが、第4節と同値な数学的仕様と比較規約
\href{supplement/computations/word-schema.md}{\nolinkurl{word-schema.md}}
から記述した。この規約が固定するのは、定義、テキスト形式、正規化だけである。提供されたプログラム、実装A、研究ノートの表は読んでいない。
Bは、C++17の本計算用プログラムと、別個に記述した逐語的なPython参照プログラムからなる。後者は、すべての \((j,c,k)\) に対する不等式の直接検査、すべての始点剰余に対する前向き反復、Kosarajuのアルゴリズム、および非巡回性を調べるKahnのアルゴリズムを用いる。両プログラムの正規化出力は、すべてのグラフでバイト単位に一致する。
Bは、最初に見つかった根からの道の語長として \(h\) を計算し、補題18の代わりに前向き反復を用いる。また、第9.1節の十分条件によってSCC分割を再検証する。
Bの出力は、Aの出力を開く前に固定し、そのサイズとハッシュを記録した。

% Markdown AST block 346: Para
\emph{比較。} Bの比較器は、29組のファイル対（表6について入力10グラフと出力10グラフ、表4についてそれぞれ4グラフ、表7について1グラフ）ごとに、ヘッダのパラメータ、頂点集合、および語ラベル付き三つ組 \((u,w,v)\) の集合を比較する。さらに、ソート順序、重複がないこと、および各ファイルの整合性を検査する。
表7については、ブロック分割を所属頂点集合の族として比較し、ブロック番号には依存させない。また、各ブロックの属性（サイズ、内部辺、\(d_B\)、\(\rho\)、\(\eta\)）、位相ラベル列、および認証されていない巡回的ブロックがないことを比較する。
29組すべてが要素ごとに一致した。比較手順と参照データは \href{supplement/COMPUTATION.md}{計算ガイド} および \href{checks/reference.json}{checks/reference.json} に記載している。公開版を再現する入口は付録Aに示す。
両実装とも、各成分の内部で共通の根からの道の総語長を用いて \(h\) を計算する。したがって、根や道が異なっていても、系22により、報告された \(d_B\) の値は一致する。

% Markdown AST block 347: Para
\emph{環境と上限。} AはLinux上でPython 3.10.12およびg++ 11.4.0を使用した。実経過時間とアドレス空間の上限は、表6では提供された各構成処理につき600 sと4 GiB（全体で1500 s）、それ以外では300 sと4 GiBとした。
BのC++プログラムの上限は、表6で3600 sと8 GiB、表4で300 sと4 GiB、表7で600 sと4 GiBとした。Python参照プログラムには1800 sと8 GiB、比較には1800 sと8 GiBの上限を設けた。
どの上限にも達しなかった。達した場合には \texttt{inconclusive} と記録することになっていた。実経過時間は実行記録に残しているが、性能に関する主張には用いない。

% Markdown AST block 348: Para
\emph{独立性の意味。} 独立性が対象とするのは、二つの実装の科学的中核である。辺の列挙、持ち上げ、成分分解、認証、および縮約を、仕様から別個に実装した。
この独立性は、共通に用いる数学的補題や共有のテキスト形式、人間による査読にまで及ぶものではなく、法律上の意味でのクリーンルーム開発を意味するものでもない。
二つの実装の一致は、それらが同じ有限対象を計算した証拠である。一方、操作の健全性（第4節）と、すべての \(\xi>0\) への移行（定理19）は、すでに上で証明した数学的命題であり、プログラムの出力ではない。

% Markdown AST block 349: 9.4 The Lean formalization and its boundary
\Needspace{6\baselineskip}
\subsection{Leanによる形式化とその範囲}\label{sec:9.4}

% Markdown AST block 350: Para
Leanによる形式化は、定理1の両部分および系2を対象とする。ただし、有限評価の信頼基盤は両者で異なる。
Lean 4.29.1と、コミット
\nolinkurl{5e932f97dd25535344f80f9dd8da3aab83df0fe6}
のMathlibを使用し、推移的依存も \nolinkurl{supplement/lean/lake-manifest.json} で固定している。

% Markdown AST block 351: Para
\textbf{1ステップ証明書。} 1ステップの形式化は、補題4の整数差分による証明、軌道の包含（命題5）、順位と位相を用いた尾部保存、および定理10の有限反復に沿っている。

% Markdown AST block 352: Para
形式的な証明書では、SCCアルゴリズムそのものの形式的証明を避けている。各段階で、各頂点をブロックに割り当て、各ブロックに自然数の順位を与える。
異なるブロック間のすべての辺では、順位が厳密に減少しなければならない。これは第6.2節の (Q-rank) における厳密に増加する \(\rho\) とは逆向きの規約であるが、どちらの規約からも整礎な順序が得られる。
各ブロックを、次のいずれかに指定する。

% Markdown AST block 353: OrderedList
\begin{enumerate}
\def\labelenumi{\arabic{enumi}.}
\tightlist
\item
  \textbf{保持：} ブロックのすべての頂点が保持集合に含まれる。
\item
  \textbf{辺なし：} ブロックに内部辺が存在しない。
\item
  \textbf{周期的：} 正整数 \(d\)、各頂点に割り当てた \(\{0,\ldots,d-1\}\) の位相、および各位相のラベルが、すべての内部辺について位相の増分条件と出発ラベル条件を満たす。
\end{enumerate}

% Markdown AST block 354: Para
これらは補題8の十分条件である。無限ウォークに沿って順位が無限回厳密に減少することはできないため、ウォークは最終的に一つのブロックにとどまる。
そのブロックに内部辺がないということはあり得ない。周期的なブロックならば、そのラベル列は補題4に矛盾する。したがって、そのブロックは保持されなければならない。
SCC分解は、その非巡回的な商グラフを通じて、このようなブロックと順位を与える。しかし、形式的な含意を示すうえでは、各ブロックがSCCであることを特定する必要はない。

% Markdown AST block 355: Para
抽象構造 \texttt{FiniteSuccess} は、すべての許容状態の初期被覆、現在の状態間で (E1)--(E3) を満たす\textbf{すべての}辺に対する順位・位相条件、保持ブロックの全頂点の被覆、次段階におけるすべての細分状態の包含、および最終保持集合の空性を要求する。
以下の宣言は、主要な形式的依存関係を示す。すべての名前空間の接頭辞は \texttt{MathPaper.} である。

% Markdown AST block 356: Table 8.
% Caption is rendered at the start of the immediately following table.

% Markdown AST block 357: Table 8
\begingroup
\small
\setlength{\tabcolsep}{4pt}
\renewcommand{\arraystretch}{1.15}
\setlength{\LTcapwidth}{\textwidth}
\begin{longtable}{@{}>{\raggedright\arraybackslash}p{\dimexpr0.44\dimexpr\linewidth-8pt\relax} >{\raggedright\arraybackslash}p{\dimexpr0.56\dimexpr\linewidth-8pt\relax}@{}}
\caption{主要なLean宣言。すべての名前の接頭辞は
\texttt{MathPaper.} であり、\texttt{Pipe.} は実行可能な語ラベル付きパイプラインを表す。}\label{tab:8}\\
\toprule
数学的内容 & Lean宣言 \\
\midrule
\endfirsthead
\toprule
数学的内容 & Lean宣言 \\
\midrule
\endhead
\bottomrule
\endlastfoot
増大性 (7) & \nolinkurl{floor_orbit_eventually_gt} \\
補題4 & \nolinkurl{floor_carry_not_eventually_periodic} \\
命題5および許容軌道状態 & \nolinkurl{orbit_edge},\allowbreak{} \nolinkurl{orbit_state_legal} \\
位相に従う周期的ラベル & \nolinkurl{phase_labels_periodic} \\
順位・位相データからの尾部保存 & \nolinkurl{rank_certificate_retained_tail},\allowbreak{} \nolinkurl{certified_orbit_retained_tail} \\
補題9 & \nolinkurl{floor_refinement_bounds},\allowbreak{} \nolinkurl{refined_state_mem} \\
証明書形式の定理10 & \nolinkurl{finite_success_visits} \\
条件付きの合成数性の結論 & \nolinkurl{finite_success_composite} \\
符号付き区間の列挙 (12) & \nolinkurl{edge_interval_iff} \\
具体的な4段階証明書 & \texttt{mainFiniteSuccess} \\
定理1(i) & \nolinkurl{floor_seven_fifths_visits} \\
\(7/5\) に対する系2 & \nolinkurl{floor_seven_fifths_composite} \\
語ラベル付き被覆と位相出力、\allowbreak{} 補題12(b) & \texttt{Covers},\allowbreak{} \nolinkurl{phase_output_periodic} \\
語ラベル付き尾部保存、\allowbreak{} 命題13 & \nolinkurl{WordCert.stage_tail} \\
マクロ頂点への到達と出力の保存、\allowbreak{} 補題15および命題16の順方向 & \nolinkurl{macro_visited},\allowbreak{} \nolinkurl{chain_realized},\allowbreak{} \nolinkurl{contract_tail} \\
素数による持ち上げ、\allowbreak{} 命題17 & \nolinkurl{runWord_orbit},\allowbreak{} \nolinkurl{lift_tail} \\
証明書形式の定理19および系20 & \nolinkurl{word_finite_success_visits},\allowbreak{} \nolinkurl{word_finite_success_composite} \\
補題23--24 & \nolinkurl{carry_alphabet_75},\allowbreak{} \nolinkurl{carry_alphabet_52} \\
初期被覆、\allowbreak{} 段階検査、\allowbreak{} 持ち上げ、\allowbreak{} およびパイプラインの健全性 & \nolinkurl{Pipe.initialStage_sound},\allowbreak{} \nolinkurl{Pipe.checkStage_sound},\allowbreak{} \nolinkurl{Pipe.liftArr_sound},\allowbreak{} \nolinkurl{Pipe.wordPipeline_sound} \\
語ラベル付き有限評価 & \nolinkurl{five_halves_ok},\allowbreak{} \nolinkurl{seven_fifths_ok} \\
定理1(ii)および \(5/2\) に対する系2 & \nolinkurl{floor_five_halves_visits},\allowbreak{} \nolinkurl{floor_five_halves_composite} \\
定理1(i)、\allowbreak{} 語ラベル付きの別証明 & \nolinkurl{floor_seven_fifths_visits_word} \\
\end{longtable}
\endgroup

% Markdown AST block 358: Para
閉じた項 \texttt{mainFiniteSuccess\ :\ FiniteSuccess\ 7\ 5\ 4290} が、一般的な還元に必要な証明書引数を与える。
1ステップによる二つの最終定理が受け取るのは、実係数 \(\xi\)、\(\xi>0\) の証明、および任意の \(N\in\mathbb N\) だけであり、それぞれの結論を満たす \(n\ge\max(N,1)\) を返す。未証明の有限計算の仮定はない。
有限検査器は、Pythonが書き出した辺リストの完全性を仮定しない。証明済みの列挙補題によって、(E1)--(E3) を満たすすべての辺が、三つの繰り上がり候補、五つの剰余の持ち上げ、および三つの到着セル候補の中に入ることを示す。検査器は、これらの辺の外側近似に対して順位・位相条件を検査する。
真偽値の検査は \texttt{decide\ +kernel} によって証明し、小さな命題へ分割したうえで、証明済みの合成補題でつないでいる。
合成数性の証明では、最終的に成り立つ評価 \(Q_n>M\) を用い、合成数性を \(Q_n=uv\)（整数 \(u,v>1\)）と定義する。素因数分解には依存しない。

% Markdown AST block 359: Para
完全な1ステップ証明書、主モジュール、および全体ビルドは成功した。
\texttt{mainFiniteSuccess} と二つの最終定理を含む、指定された30宣言の公理監査では、依存する公理はLeanの標準公理 \texttt{propext}、\texttt{Classical.choice}、\texttt{Quot.sound} に限られていた。
\texttt{sorryAx}、独自の公理、外部計算の公理は現れない
（\href{supplement/FORMALIZATION.md}{形式化ガイド}）。
独立監査では、1ステップのプロジェクトの全79モジュールを新しいプロジェクトキャッシュから再ビルドし、60宣言の型と公理依存を検査した。その内訳は、プロジェクトの58宣言と、定理1(i)およびその系を独立に書き直した二つの宣言である。そこでも同じ三つの標準公理だけが確認された。
明示した検査範囲と、型・公理検査を再実行するコマンドについては、\href{supplement/FORMALIZATION.md}{形式化ガイド} を参照されたい。
この再ビルドでは、固定したMathlibとその依存ライブラリのコンパイル済みキャッシュを使い、それらのリビジョンをロックファイルと照合した。すべての依存ライブラリをソースから再ビルドしたわけではない。
この1ステップ証明の信頼基盤は、固定したLeanカーネル、その標準公理、および固定したライブラリからなる。
宣言ごとの対応範囲は \href{supplement/FORMALIZATION.md}{形式化ガイド} にまとめ、公開版のビルド手順は \href{REPRODUCE.md}{再現ガイド} に記載している。

% Markdown AST block 360: Para
\textbf{語ラベル付き証明書。} \nolinkurl{MathPaper/Word/} に追加した形式化は、ブロック・順位の手法を非空の辺語へ拡張する。
その被覆関係には、ウォークと厳密に増加する境界時刻を記録し、各辺語が繰り上がり列の対応する区間に等しいことを要求する。
周期的ブロックでは、位相は語長全体の分だけ進み、各文字は対応する位相ラベルに一致しなければならない。
保持ブロックでは、検査器はマクロ頂点集合と、各非マクロ頂点に対する残りの語および到着先を検査する。整合性条件によって、非マクロ頂点からなる道に沿って残りの語長が減少し、マクロ辺への移行時には連結された出力が保存される。
したがって、この形式的証明は命題16の順方向を用いている。出力グラフのすべてのウォークが入力グラフのウォークへ展開できることを示す必要はない。

% Markdown AST block 361: Para
カーネルで検査された定理 \nolinkurl{Pipe.wordPipeline_sound} は、
\texttt{Pipe.wordPipeline\ a\ b\ K\ alphabet\ primes\ =\ true} から
\texttt{WordFiniteSuccess\ a\ b\ K\ alphabet\ primes} が従うことを証明する。
このパイプラインは、厳密不等式 (E3) によって初期セル辺を列挙する。各段階で、補助アルゴリズムがブロック、順位、位相、縮約データ、および出力辺の候補を計算する。その正しさは仮定しない。証明済みの検査器が、すべての入力辺に対して必要な条件を検証する。
証明済みの持ち上げ関数は、すべての非零の始点剰余とすべての文字を検討し、最終終点の剰余も検査する。そこで使う \(b\) の逆元は、新たに追加した素数を法として検査される。
頂点の番号の振り直しには、任意の頂点写像のもとでの被覆の保存を用いるため、未証明の単射性の仮定は不要である。
パイプラインが受理するのは、最後の出力が空になったときだけである。
その健全性定理は、証明済みの二つの繰り上がりアルファベット補題とともに、一般的な訪問定理と合成数性定理に必要な有限証明書の引数を与える。

% Markdown AST block 362: Para
受理された二つの有限評価は、次のとおりである。

% Markdown AST block 363: reproduction command
\begin{Verbatim}[fontsize=\small,breaklines=true,breakanywhere=true]
Pipe.wordPipeline 5 2 4 [-1, 1, 3] [3, 7, 11, 13, 17, 19, 23, 29, 31] = true
Pipe.wordPipeline 7 5 2048 [-4, -2, 2, 4, 6] [3, 11, 13] = true
\end{Verbatim}

% Markdown AST block 364: Para
これらは \nolinkurl{native_decide} によって証明される。固定したLeanの版では、この仕組みは各定理に対して生成される公理を通じてネイティブ評価の結果を表明する。\texttt{Lean.ofReduceBool} は使用しない。
実際に追加される公理は、名前空間接頭辞を \texttt{MathPaper.} として、それぞれ
\nolinkurl{five_halves_ok._native.native_decide.ax_1_1} と
\nolinkurl{seven_fifths_ok._native.native_decide.ax_1_1}
である。
したがって、\(5/2\) の最終定理は、三つの標準公理に加えて第1の生成公理に依存し、\(7/5\) の別証明による定理は第2の生成公理に依存する。
語ラベル付き理論と検査器の健全性は、標準公理だけに依存する。記録された依存関係に \texttt{sorryAx} は現れない。
追加される信頼の対象は、Leanパイプラインのコンパイラとネイティブ評価であって、PythonやC++の成功フラグ、書き出されたグラフ、ハッシュではない。
これらの最終定理が受け取るのは \(\xi>0\) と \(N\) だけであり、証明書引数は不要である。そして、定理1および系2の \(5/2\) の部分と正確に同じ量化と法を与える。

% Markdown AST block 365: Para
追加モジュール、全体ビルド、および公理の表示は受理された。形式化の範囲と公理に関する信頼境界は \href{supplement/FORMALIZATION.md}{形式化ガイド} に記載している。別に行われた語ラベル付き証明の独立監査では、追加した全14個のWordモジュールを再ビルドし、二つのネイティブ有限評価を再実行した。さらに、プロジェクトの46宣言と、独立な言い換えまたは検査を行う五つの宣言の型と公理依存を調べ、上記の信頼範囲を確認した。
この監査では、コンパイル済みの1ステップのプロジェクト依存物と、固定したライブラリのキャッシュを再利用した。1ステップの全モジュールやライブラリ全体を新たに再コンパイルしたわけではない。
より小さなパイプラインの例で \texttt{decide\ +kernel} による評価を試みたが、不透明なハッシュ演算で停止した。したがって、語ラベル付きの二つの有限評価のどちらについても、カーネルだけによる有限計算とは主張しない。

% Markdown AST block 366: Para
\textbf{形式化の範囲。} この形式化は、上述の信頼基盤のもとで、定理1と系2の算術的結論を確立する。
報告したブロックがSCCであることを形式的に同定したり、\(h\) が最短距離であることや \(d\) が指定された最大公約数であることを証明したり、表に報告した件数や語を確立したりするものではない。
これらのデータはA/B計算によって検査される。Leanの検査器が用いるのは、順位、位相、および縮約に関する十分条件である。
第4節のうち、補題12(a)の最大公約数の正値性の議論、補題14、命題16の逆向きの対応、補題18の根の公式、および定理21と系22は形式化していない。検査器は正の周期を要求し、持ち上げには前向きの漸化式を用いる。
定理3、第6節（表7の待ち時間証明書を含む）、第7--8節、第10節の追加結果、および付録B--Eも、形式化の対象外である。
これらの根拠は、本論文中の証明と、有限計算が必要な場合には記録されたA/B比較からなる。

% Markdown AST block 367: 10. Further results deferred to the appendices
\Needspace{6\baselineskip}
\section{証明を付録に収めた追加の結果}\label{sec:10}

% Markdown AST block 368: Para
ここでは二群の追加の結果を証明なしで述べる。証明は付録BおよびDに示す。これらを裏付ける計算は9.3節の二つの実装で実行し、要素ごとに比較した（D.5節およびE.2節）。これらの結果は1--9節では用いない。

% Markdown AST block 369: Para
\textbf{再帰的な繰り上がり列。} 有限アルファベット上の無限語 \(w=w_0w_1\cdots\) が\emph{再帰的}であるとは、任意の \(L\ge1\) と任意の \(B\ge1\) に対し、ある \(n\ge B\) が存在して
\[
w_{n+i}=w_i
\]
が \(0\le i<L\) について成り立つことをいう。再帰的な語の任意の尾部も再帰的である。

% Markdown AST block 370: Theorem 46 (recurrent carries force composites; proof in Appendix B.3).
\begin{samepage}
\begin{theorem}[再帰的な繰り上がり列は合成数を生じさせる；証明は付録B.3]\leavevmode\label{thm:46}
\(a>b\ge2\)、\(\gcd(a,b)=1\)、\(\xi>0\) とし、\emph{軌道の繰り上がり列 \(c_0c_1\cdots\) が再帰的な尾部} \(c_{j_0}c_{j_0+1}\cdots\) \emph{をもつと仮定する}。このとき、\((Q_n)\) は合成数の項を無限に多く含む。さらに、\(j\ge j_0\) が
\(Q_j>1\) および \(\gcd(Q_j,ab)=1\) を満たすならば、\(Q_j\) はそれより後の無限に多くの項を割り切る。
\label{thm-end:46}
\end{theorem}
\end{samepage}

% Markdown AST block 371: Para
再帰性の仮定は、繰り上がり語そのものに関するものである。ウォークが一つの強連結成分内にとどまることから、その出力の再帰性は導かれない。

% Markdown AST block 371 (continued): Para
定理46は初期の剰余に戻ることを示す議論であり、Dubickasの定理 \hyperref[ref-dub09]{Dub09, Theorem 4} と同じ発想に基づく。この定理は、整数 \(x_n\) が整数 \(k\ne0\) と \(F\in\mathbb Z[z_1,\ldots,z_{d-1}]\) について \(x_{n+d}=kx_n+F(x_{n+1},\ldots,x_{n+d-1})\) を満たすならば、\(\gcd(k,q)=1\) を満たす任意の \(q\) を法として \((x_n)\) は純周期的であり、\(|x_n|\to\infty\) ならば合成数の項を無限に多く含む、というものである。Stephan \hyperref[ref-s26]{S26, Proposition 3.1} は、整数の底における最終的に周期的な桁の語に、その \(d=1\) の場合を適用している。非整数の底では、繰り上がり語が最終的に周期的になることはない（補題4）。定理46では、各文字がアフィン置換 \(x\mapsto b^{-1}(ax+c)\) によって \(\mathbb Z/M\) に作用し、語には再帰性だけを仮定する。そして、周期性の代わりに有限接頭語のコピー補題（補題49）が初期の剰余への回帰を生み出す（定理50）。

% Markdown AST block 372: Para
\textbf{有限高さの結果。} 底 \(6/5\) および \(9/7\) については、繰り上がり語に対する篩を厳密な区間判定および約数判定と組み合わせることで、明示した高さ以下のすべての初期整数を扱える。

% Markdown AST block 373: Theorem 47 (finite height; proof in Appendix D.4).
\begin{samepage}
\begin{theorem}[有限高さ；証明は付録D.4]\leavevmode\label{thm:47}
(a) \(6<\lfloor\xi\rfloor\le5^{40}\) ならば、\(\lfloor\xi(6/5)^j\rfloor\) は、ある \(0\le j\le40\) について合成数である。
(b) \(9<\lfloor\xi\rfloor\le7^{18}\) ならば、\(\lfloor\xi(9/7)^j\rfloor\) は、ある \(0\le j\le18\) について合成数である。
\label{thm-end:47}
\end{theorem}
\end{samepage}

% Markdown AST block 374: Para
\(6/5\) の場合、平行移動した数列
\(\lfloor\xi(6/5)^n\rfloor-1\) は
\hyperref[ref-dn]{DN, Theorem 3, p.637} で扱われている。定理47(a)は平行移動していない数列の有限範囲に関する結果であり、これとは異なる主張である。

% Markdown AST block 375: Para
付録Cでは族 \((6t+3)/(6t+1)\) を扱い、この族について全係数に対する主張が厳密な整数停止問題に帰着されることを示す（定理58）。付録Eでは、この族に沿った任意に長い素数接頭列（定理68。素数の狭い等差数列に関する既知の結果を外部入力として用いる）、\(9/7\) および \(6/5\) に対する明示的な帰還語の障害（系70）、ならびに8節から持ち越した事項を扱う。

% Markdown AST block 376: 11. Open problems
\Needspace{6\baselineskip}
\section{未解決問題}\label{sec:11}

% Markdown AST block 377: OrderedList
\begin{enumerate}
\def\labelenumi{\arabic{enumi}.}
\tightlist
\item
  \emph{他の底と無限族。} どのような他の有理数の底が有限証明書をもつかの決定と、無限族に対する一様な証明書の構成が課題として残る。定理44は外部素数の積に対する必要条件 \(R\ge a^{1-o(1)}\) を与えるが、これを達成する構成は与えない。特定の帰還語を排除すること（付録E.3）と、すべての再帰的成分を覆うことは異なる課題である。
\item
  \emph{有限高さと全係数。} 定理47では、述べた高さを超える初期整数は扱われていない。初期整数 \(s+kb^N\) が 最初の \(N\) 個について \(s\) と同じ繰り上がりをもつ場合、(8)により各項は \(Q_j(s)+ka^jb^{N-j}\) となり、\(k=0\) に対して見つかった約数が、他の \(k\) に対する項も割り切るとは限らない。
\item
  \emph{再帰性。} 定理46は、すべての項が素数である軌道に再帰的な繰り上がり尾部が存在することを排除する。しかし、非再帰的な全素数尾部が存在しうるかは未解決である（系51）。外側グラフの一つの強連結成分内にとどまることから再帰性は導かれない（付録B.1）。
\item
  \emph{追加の操作。} 全語区間判定（付録E.4）は擬軌道に対して健全だが、ウォークに対しては健全ではない。これを追加することでクラス \(\mathcal F_{\rm cert}\) が変わるかは未解決である。付録Cの停止プログラムがすべての \(t\) に対して停止するか（定理58）も未解決である。
\item
  \emph{最小性と部分列。} 証明書は有限素数集合を与えるが、その最小性は証明していない。指定した添字の等差数列や疎な部分列に沿って結論が成り立つか、また素数項の密度にどのような評価が成り立つかは、別の問題として残る。
\end{enumerate}

% Markdown AST block 378: Appendix A. Reproducing the computations and the formal proof
\appendix
\Needspace{6\baselineskip}\section{計算と形式証明の再現}\label{app:A}

% Public reproduction instructions.
論文ソース、科学的プログラム、およびLeanプロジェクトは、公開リポジトリ \href{https://github.com/ixixi/rational-floor-certificates}{rational-floor-certificates} に収録している。コマンドはダウンロードした配布物のルートから実行する。\href{REPRODUCE.md}{再現ガイド} に依存関係、コマンド、資源上限、出力先を示す。共通の入口は次のとおりである。

\begin{Verbatim}[fontsize=\small,breaklines=true,breakanywhere=true]
python3 -B reproduce.py verify
python3 -B reproduce.py check --output ../certificate-preflight
python3 -B reproduce.py recompute --output ../certificate-replay
\end{Verbatim}

配布物のソースは、\href{MANIFEST.json}{ソースマニフェスト} の \nolinkurl{source_sha256} フィールドで識別する。この値は、同梱するソースと検査用ファイルのSHA-256値から計算する。個々のファイルのハッシュは \nolinkurl{SHA256SUMS} に列挙し、比較に用いる有限の参照データは \href{checks/reference.json}{checks/reference.json} に記録する。生成するグラフとログは指定した出力先に保存する。資源上限に達した場合は \nolinkurl{inconclusive} と記録し、未完了の実行を受理済みの証明書とはしない。

\textbf{語ラベル付き計算、実装A。} 公開ドライバ \nolinkurl{supplement/computations/word_a/run_word.py} は、資源上限を適用し、提供された科学的構成を実行し、正規化グラフを書き出して、同系統内の事後検査を実行する。\nolinkurl{five-halves}、\nolinkurl{cross-75}、\nolinkurl{q75} の各ケースは、それぞれ表6、4、7のデータを再生成する。完全なコマンドは \href{REPRODUCE.md}{再現ガイド}、ソースの来歴と比較範囲は \href{supplement/COMPUTATION.md}{計算ガイド} に記載している。C++による構成とPythonによる正規化コードを同梱しているため、配布されたソースから再生成できる。

\textbf{辺の重複度。} 表4および6では、辺を相異なる三つ組として数えている。提供された研究ノートと同様に、縮約した鎖を重複度込みで数え、持ち上げに沿って重複度を伝播させると、表6の十段階の入力辺数は12, 17, 58, 167, 820, 7083, 60242, 315483, 1685458, 3754612、出力辺数は12, 12, 20, 85, 529, 4256, 18069, 80704, 201199, 0となる。表4では重複度はすべて1である。両実装がこれらの数を報告し、一致している。ただし、これらの数自体は証明内容を担わない。

\textbf{語ラベル付き計算、実装B。} 独立なC++プログラム、Python参照実装、および比較器は \nolinkurl{supplement/computations/word_b/} に収録している。次のコマンドに数学的入力を示す。出力先は配布物の外部にある、実行前には存在しないディレクトリとする。公開版の入口は資源上限を適用し、第9.3節に記載した完全な比較を実施する。

\begin{Verbatim}[fontsize=\small,breaklines=true,breakanywhere=true]
mkdir -p ../certificate-manual/build \
  ../certificate-manual/OUT52 ../certificate-manual/OUT75 \
  ../certificate-manual/OUTQ75 ../certificate-manual/REF52
g++ -O2 -std=c++17 -o ../certificate-manual/build/wgb supplement/computations/word_b/wgb.cpp
../certificate-manual/build/wgb pipeline --a 5 --b 2 --K 4    --carries -1,1,3      --primes 3,7,11,13,17,19,23,29,31 --out ../certificate-manual/OUT52
../certificate-manual/build/wgb pipeline --a 7 --b 5 --K 2048 --carries -4,-2,2,4,6 --primes 3,11,13 --out ../certificate-manual/OUT75
../certificate-manual/build/wgb q75      --a 7 --b 5 --M 429 --K 2048 --carries -4,-2,2,4,6 --out ../certificate-manual/OUTQ75
python3 supplement/computations/word_b/wgb_ref.py pipeline --a 5 --b 2 --K 4 --carries -1,1,3 --primes 3,7,11,13,17,19,23,29,31 --out ../certificate-manual/REF52
python3 supplement/computations/word_b/compare_wordgraph.py batch --out ../certificate-manual/CMPDIR LABEL=FILE_A:FILE_B ...
\end{Verbatim}

共有する正規化規約は \href{supplement/computations/comparison-schema.md}{comparison-schema.md} と \href{supplement/computations/word-schema.md}{word-schema.md} である。参照データと成果物のハッシュは \href{checks/reference.json}{checks/reference.json} および \href{MANIFEST.json}{ソースマニフェスト} に記録している。ハッシュは有限の成果物を識別するものであり、要素ごとの比較や数学的な完全性の議論の代わりにはならない。

\textbf{1ステップ計算。} 公開版の再現手順は、表1--3の四つのケースについても両方の1ステップ実装で再生成し、出力全体を比較する。独立検証器では検査にassertionを使うため、Pythonの最適化を無効にして実行する。ソースプログラム、参照検査データ、および生成済みのLean証明書ソースを同梱している。比較対象のデータは \href{supplement/COMPUTATION.md}{計算ガイド}、コマンドは \href{REPRODUCE.md}{再現ガイド} に示す。

\textbf{形式証明。} 生成済みのLeanデータと証明ソースを同梱しているため、1ステップ証明の検査に生成器の再実行は必要ない。語ラベル付きモジュールは有限処理系をLean内で定義・実行し、A/Bのグラフダンプを読み込まない。\href{REPRODUCE.md}{再現ガイド} に従って固定した依存ライブラリを取得し、新しいプロジェクトキャッシュを用いてモジュールを順にビルドする。1ステップの有限モジュールの後、最終の \nolinkurl{MathPaper} ターゲットと全体ビルドの前に、語ラベル付きモジュールを順に検査する。

\begin{Verbatim}[fontsize=\small,breaklines=true,breakanywhere=true]
python3 supplement/lean/tools/check_finite.py \
  MathPaper.Word.Graph MathPaper.Word.Cert MathPaper.Word.Lift \
  MathPaper.Word.Chain MathPaper.Word.Alphabet \
  MathPaper.Word.Pipeline.Data MathPaper.Word.Pipeline.Check \
  MathPaper.Word.Pipeline.Compute MathPaper.Word.Pipeline.LiftArr \
  MathPaper.Word.Pipeline.Initial MathPaper.Word.Pipeline.Main \
  MathPaper.Word.FiveHalves MathPaper.Word.SevenFifths MathPaper.Word
\end{Verbatim}

\nolinkurl{FiveHalves} および \nolinkurl{SevenFifths} ターゲットは、新しいプロジェクトキャッシュからビルドすると、二つの有限な \nolinkurl{native_decide} 評価を再実行する。その後、\href{REPRODUCE.md}{再現ガイド} に従って最終ターゲット、全体ビルド、および \nolinkurl{Audit.lean} による公理表示を実行し、依存関係を \href{supplement/FORMALIZATION.md}{形式化ガイド} に明示した公理の境界と照合する。追加の評価公理は第9.4節で説明した信頼境界であり、これらの評価はカーネルだけによる計算ではない。完了した独立再ビルドと型・公理検査の範囲も第9.4節に示している。

% Markdown AST block 397: Appendix B. Recurrent carries
\Needspace{6\baselineskip}\section{再帰的な繰り上がり列}\label{app:B}

% Markdown AST block 398: Para
本付録では定理46と、族 \((6t+3)/(6t+1)\) に対する整数コピー条件を証明する。ここでの結果は整数漸化式と有限置換に関する一般的な補題であり、3--9節には依存せず、それらの節でも用いない。

% Markdown AST block 399: B.1 Recurrent words
\Needspace{6\baselineskip}
\subsection{再帰語}\label{sec:B.1}

% Markdown AST block 400: Para
10節の定義を再掲する。有限アルファベット上の無限語 \(w=w_0w_1\cdots\) が\emph{再帰的}であるとは、任意の \(L\ge1\) および \(B\ge1\) に対して、ある \(n\ge B\) が存在し、
\[
w[n:n+L]=w[0:L],
\]
が成り立つことをいう。ここで
\[
w[i:j]=w_iw_{i+1}\cdots w_{j-1}.
\]
である。同値な条件として、任意の有限接頭語が無限回現れること、あるいは一度でも現れるすべての因子が無限回現れることを用いてもよい（因子はいずれかの接頭語に含まれる）。

% Markdown AST block 401: Lemma 48 (tails).
\begin{samepage}
\begin{lemma}[尾部]\leavevmode\label{thm:48}
\(w\) が再帰的ならば、\(w[k:]\) も、任意の \(k\ge0\) に対して再帰的である。
\label{thm-end:48}
\end{lemma}
\end{samepage}

% Markdown AST block 402: Proof.
\begin{proof}\leavevmode\label{proof:48}
\(L,B\ge1\) が与えられたとする。\(w\) の再帰性を、接頭語長 \(k+L\) と下限 \(B\) に適用すると、ある \(n\ge B\) が存在して
\[
w[n:n+k+L]=w[0:k+L].
\]
となる。位置 \(k\) 以降を比較すると
\[
w[n+k:n+k+L]=w[k:k+L],
\]
を得る。すなわち
\[
w'[n:n+L]=w'[0:L]
\]
が
\(w'=w[k:]\)
について成り立つ。\qedhere
\end{proof}

% Markdown AST block 403: Para
一様再帰性は必要ない。例えば、\(w_0=0\) とし、
\[
w_{k+1}=w_k\,1\,0^{k+1}\,1\,w_k;
\]
とおく。各 \(w_k\) は \(w_{k+1}\) の接頭語なので、極限語 \(w_\infty\) が存在する。その中では、すべての \(w_k\) が無限回現れる（\(w_{j+1}\) の位置 \(|w_j|+j+3\) に、任意の \(j\ge k\) に対して現れる）。一方、0の連続長には上限がなく、\(1\) が無限回現れるので、\(w_\infty\) は最終的周期的ではない。これに対し、語 \(1\,0\,1\,00\,1\,000\cdots\) では0の連続長が \(1,2,3,\ldots\) であり、再帰的な尾部は存在しない。実際、任意の尾部のうち文字 \(1\) を二つ含む接頭語には因子 \(1\,0^k\,1\) が含まれるが、連続長は狭義に増加するため、この因子は語全体で一度しか現れない。この語は二つの自己ループをもつ1頂点グラフで生成される。したがって、ウォークが一つの強連結成分内にとどまることから、その出力の再帰性は導かれない。

% Markdown AST block 404: Para
いずれの例についても、真の軌道の繰り上がり列であるとは主張していない。

% Markdown AST block 405: B.2 The finite-prefix copying lemma
\Needspace{6\baselineskip}
\subsection{有限接頭語のコピー補題}\label{sec:B.2}

% Markdown AST block 406: Lemma 49 (finite-prefix copying).
\begin{samepage}
\begin{lemma}[有限接頭語のコピー]\leavevmode\label{thm:49}
\(X\) を \(m\) 個の元をもつ集合とし、各文字 \(c\) が 置換 \(f_c\) として \(X\) に作用するとする。\(w\) を有限語、\(t_1,\ldots,t_m\) を正整数とし、\(S_0=0\)、\(S_k=t_1+\cdots+t_k\)、\(L\ge0\) が
\[
w[t_k:\,t_k+S_{k-1}+L]=w[0:\,S_{k-1}+L]\qquad(k=1,\ldots,m)
\tag{31}\label{eq:31}
\]
および \(|w|\ge S_m+L\) を満たすとする。
\[
G_n=f_{w_{n-1}}\circ\cdots\circ f_{w_0}
\]
とおく（\(G_0=\mathrm{id}\)。写像は右から作用するため、\(G_n\) は最初の文字を最初に適用する）。このとき、任意の \(x\in X\) に対し、ある
\(0\le i<j\le m\) が存在して、
\[
n=t_{i+1}+\cdots+t_j>0
\]
は
\[
G_n(x)=x\qquad \text{and}\qquad w[n:n+L]=w[0:L].
\]
を満たす。
\label{thm-end:49}
\end{lemma}
\end{samepage}

% Markdown AST block 407: Proof.
\begin{proof}\leavevmode\label{proof:49}
\emph{(1) 合成則。} \(0\le s\le S_{k-1}\) ならば
\[
G_{t_k+s}=G_s\circ G_{t_k}:
\]
である。実際、
\[
G_{t_k+s}=f_{w_{t_k+s-1}}\circ\cdots\circ f_{w_{t_k}}\circ G_{t_k},
\]
であり、(31)により、\(w_{t_k+i}=w_i\) が \(0\le i<s\) について成り立つので、左の因子は \(G_s\) である。

% Markdown AST block 408: Para
\emph{(2) コピーは合成できる。} 添字 \(i_1<\cdots<i_h\) および
\(n=t_{i_1}+\cdots+t_{i_h}\) について、
\[
\begin{aligned}
G_n&=G_{t_{i_1}}\circ G_{t_{i_2}}\circ\cdots\circ G_{t_{i_h}},\\
w[n:n+L]&=w[0:L].
\end{aligned}
\]
が成り立つ。\(h\) に関する帰納法で示す。\(h=0\) の場合は \(G_0=\mathrm{id}\) である。\(h\ge1\) の場合は
\[
n'=t_{i_1}+\cdots+t_{i_{h-1}}
\]
および \(k=i_h\) とおく。すべての
\(i_j<k\) より、\(n'\le S_{k-1}\) である。(1)により
\[
G_n=G_{t_k+n'}=G_{n'}\circ
G_{t_k},
\]
であり、帰納法の仮定により
\[
G_{n'}=G_{t_{i_1}}\circ\cdots\circ
G_{t_{i_{h-1}}}.
\]
である。

% Markdown AST block 409: Para
語については、\(n'+L\le S_{k-1}+L\) であるから、(31)により
\[
w[t_k+n':t_k+n'+L]=w[n':n'+L]=w[0:L].
\]
を得る。

% Markdown AST block 410: Para
\emph{(3) 鳩の巣原理。} \(H_0=\mathrm{id}\) とおき、
\[
H_k=G_{t_1}\circ\cdots\circ G_{t_k};
\]
とする。(2)を \(\{1,\ldots,k\}\) に適用すると、\(H_k=G_{S_k}\) となり、これは時刻 \(S_k\) における実際の状態を与える。

% Markdown AST block 411: Para
\(m+1\) 個の値 \(H_0(x),\ldots,H_m(x)\) は \(m\) 元集合 \(X\) に属するので、そのうち二つは一致する。すなわち、\(H_i(x)=H_j(x)\) となる \(i<j\) が存在する。

% Markdown AST block 412: Para
ここで
\[
B=G_{t_{i+1}}\circ\cdots\circ G_{t_j}
\]
とおくと、\(H_j=H_i\circ B\) なので、
\[
H_i(x)=H_i(B(x)),
\]
を得る。\(H_i\) は置換の合成であり単射だから、\(B(x)=x\) が従う。

% Markdown AST block 413: Para
(2)を \(\{i+1,\ldots,j\}\) に適用すると、
\[
\begin{aligned}
&B=G_n\qquad \text{with}\qquad n=t_{i+1}+\cdots+t_j>0\qquad \text{and}\\
&w[n:n+L]=w[0:L].
\end{aligned}\qedhere
\]
となる。
\end{proof}

% Markdown AST block 414: Para
単射性は(3)でただ一度だけ用いており、単射でない作用では結論は成り立たない。軌道上の状態に直接鳩の巣原理を適用しても、\emph{ある}状態への帰還が得られるだけで、初期状態への帰還は得られない。コピーによって合成の位置を揃えることで、共通する左因子を消去できる。

% Markdown AST block 415: Theorem 50 (residue recurrence).
\begin{samepage}
\begin{theorem}[剰余の再帰性]\leavevmode\label{thm:50}
整数列 \((Q_n)_{n\ge0}\)、\((c_n)_{n\ge0}\) が
\[
bQ_{n+1}=aQ_n+c_n,
\]
を満たし、語 \(c=c_0c_1\cdots\) は再帰的とする。また、\(M\ge2\) は
\(\gcd(M,ab)=1\) を満たすとする。このとき、任意の \(B\ge1\) および \(L\ge0\) に対し、ある
\(n\ge B\) が存在して
\[
Q_n\equiv Q_0\pmod M\qquad \text{and}\qquad c[n:n+L]=c[0:L].
\]
が成り立つ。
\label{thm-end:50}
\end{theorem}
\end{samepage}

% Markdown AST block 416: Proof.
\begin{proof}\leavevmode\label{proof:50}
\(X=\mathbb Z/M\) とし、
\[
f_c(x)=b^{-1}(ax+c);
\]
とおく。\(a,b\) は法 \(M\) で可逆であるから、各 \(f_c\) はアフィン置換であり、
\[
Q_n\equiv G_n(Q_0).
\]
が成り立つ。

% Markdown AST block 417: Para
コピー時刻を順に選ぶ。\(k=1,\ldots,M\) に対して、必要な接頭語長 \(S_{k-1}+L\) は有限である。したがって、再帰性により（この長さが正の場合。そうでない場合は任意の \(t_k\ge B\) を選べばよい）、(31)を満たす \(t_k\ge B\) が存在する。ここで、\(S_{k-1}\) は \(t_1,\ldots,t_{k-1}\) によって定まる。

% Markdown AST block 418: Para
補題49を \(m=M\) および \(x=Q_0\bmod M\) に適用すると、
\[
n=t_{i+1}+\cdots+t_j\ge t_j\ge B,
\]
と \(G_n(Q_0)\equiv Q_0\)、および接頭語の帰還が得られる。\qedhere
\end{proof}

% Markdown AST block 419: Para
仮定 \(\gcd(M,ab)=1\) は取り除けない。素数 \(p\mid M\) が \(a\) を割り切ると、\(f_c(x)\equiv b^{-1}c\pmod p\) は定数写像となって単射ではなく、\(Q_n\bmod p\) は \(c_{n-1}\) だけで定まる。\(Q_0\bmod p\) が、形 \(b^{-1}c\) に語のどの文字 \(c\) を用いても表せないとする（例えば、\(p\mid Q_0\) であり、どの文字も \(p\) で割り切れない場合）。このとき、語が再帰的であっても、\(Q_n\equiv Q_0\pmod p\) は \(n\ge1\) に対して決して成り立たない。

% Markdown AST block 420: B.3 Proof of Theorem 46
\Needspace{6\baselineskip}
\subsection{定理46の証明}\label{sec:B.3}

% Markdown AST block 421: Proof of Theorem 46.
\begin{proof}[定理46の証明]\leavevmode\label{proof:46}
\(c[j_0:]\) は再帰的とする。

% Markdown AST block 422: Para
\emph{場合1：ある \(j\ge j_0\) が \(Q_j>1\) および \(\gcd(Q_j,ab)=1\) を満たす。} 補題48により、\(c[j:]\) は再帰的である。定理50を数列 \(Q'_n=Q_{j+n}\)、\(c'_n=c_{j+n}\) に対し、\(M=Q_j\ge2\) および \(L=0\) として適用する。任意の \(B\) に対して、ある \(n\ge B\) が存在して
\[
Q_{j+n}\equiv Q_j\equiv0\pmod{Q_j}.
\]
となる。(7)により、
\[
Q_{j+n}>Q_j>1
\]
はすべての十分大きな \(n\) に対して成り立つ。したがって、\(Q_j\) は真の約数であり、\(Q_{j+n}\) は合成数である。このような \(n\) は無限に多く存在する。これにより定理の最後の主張も証明された。

% Markdown AST block 423: Para
\emph{場合2：そのような \(j\) は存在しない。} (7)により、\(Q_j\le1\) となる \(j\) は有限個しかない。したがって、
\[
\gcd(Q_j,ab)>1
\]
がすべての十分大きな \(j\) について成り立つ。素数 \(p\mid\gcd(Q_j,ab)\) は \(p\le ab\) を満たすので、\(Q_j>ab\) となれば
\(Q_j=pv\) かつ \(v>1\) である。\qedhere
\end{proof}

% Markdown AST block 424: Para
場合2では再帰性を用いないが、場合1の「\(Q_j\) がそれより後の無限に多くの項を割り切る」という結論では再帰性を用いている。一般には \(Q_j\mid Q_n\) が起こるとは限らない。このため、定理の主張には再帰性の仮定が含まれている。

% Markdown AST block 425: Corollary 51 (refutation form).
\begin{samepage}
\begin{corollary}[反証形]\leavevmode\label{thm:51}
項 \(Q_n\) が \(n\ge n_0\) に対してすべて素数ならば、繰り上がり語は再帰的な尾部をもたない。
\label{thm-end:51}
\end{corollary}
\end{samepage}

% Markdown AST block 426: Proof.
\begin{proof}\leavevmode\label{proof:51}
\(c[j_0:]\) が再帰的と仮定する。(7)により、\(j\ge\max(j_0,n_0)\) を \(P=Q_j>a\) を満たすように選ぶ。このとき、\(P\) は \(a>b\) より大きい素数なので、\(\gcd(P,ab)=1\) である。定理50を \(M=P\) として適用すると、\(n\) を任意に大きく取ることができ、それに対して \(P\mid Q_n\) および \(Q_n>P\) が成り立ち、したがって合成数となる。これは矛盾である。\qedhere
\end{proof}

% Markdown AST block 427: Para
条件 \(P>a\) の役割は、\(\gcd(P,ab)=1\) を保証することだけである。定理46は、ある時点以降のすべての項が素数である軌道に再帰的な繰り上がり尾部が存在することを排除する。しかし、繰り上がり列のすべての尾部が非再帰的である軌道は排除しない（B.1節で、このような語が語としては存在することを示した）。この残された場合は未解決である（11節）。

% Markdown AST block 428: B.4 The integer copying condition for the family \((6t+3)/(6t+1)\)
\Needspace{6\baselineskip}
\subsection{\texorpdfstring{族 \((6t+3)/(6t+1)\) に対する整数コピー条件}{族 (6t+3)/(6t+1) に対する整数コピー条件}}\label{sec:B.4}

% Markdown AST block 429: Para
ここで
\[
\begin{aligned}
&t\ge1,\qquad a=6t+3,\qquad b=6t+1,\\
&\delta=t\bmod2,\qquad H=3t-1-\delta
\end{aligned}
\]
（奇整数）とし、
\[
\begin{aligned}
T_t(Q)&=Q+2\Bigl\lfloor\frac{Q+3t+1}{b}\Bigr\rfloor,\\
F_t(z)&=\Bigl\lceil\frac{az-\delta}{b}\Bigr\rceil,\\
e&=bF_t(z)-az+\delta .
\end{aligned}
\tag{32}\label{eq:32}
\]
とおく。付録Cでは、すべての項が \(6\) と互いに素である真の軌道の尾部が、写像 \(T_t\) に従うことを示す。ここでは \(T_t\) を整数写像として調べる。

% Markdown AST block 430: Lemma 52 (conjugation and digits).
\begin{samepage}
\begin{lemma}[共役と数字]\leavevmode\label{thm:52}
奇整数 \(Q=2z+H\) に対して次が成り立つ。
(a) \(T_t(Q)=2F_t(z)+H\) であり、\(T_t\) は偶奇性を保存する。
(b)
\[
0\le e\le b-1\qquad \text{and}\qquad e\equiv\delta-az\pmod b;
\]
(c) 繰り上がり
\(c=bT_t(Q)-aQ\) は \(2(e-(3t-1))\) に等しい。
(d)
\[
F_t(s+bv)=F_t(s)+av
\]
がすべての整数 \(s,v\) について成り立つ。
\label{thm-end:52}
\end{lemma}
\end{samepage}

% Markdown AST block 431: Proof.
\begin{proof}\leavevmode\label{proof:52}
(a)
\[
Q+3t+1=2z+(b-1)-\delta,
\]
なので、
\[
\lfloor(Q+3t+1)/b\rfloor=1+\lfloor(2z-1-\delta)/b\rfloor.
\]
である。\(a=b+2\) とおくと、
\[
\begin{aligned}
&(az-\delta)/b=z+(2z-\delta)/b,\qquad\text{and}\\
&\lceil n/b\rceil=\lfloor(n-1)/b\rfloor+1
\end{aligned}
\]
が整数 \(n\) について成り立つことから、
\[
F_t(z)=z+1+\lfloor(2z-1-\delta)/b\rfloor.
\]
を得る。したがって、
\[
2F_t(z)+H=2z+H+2+2\lfloor(2z-1-\delta)/b\rfloor=T_t(Q).
\]
となる。\(T_t(Q)-Q\) は偶数なので、偶奇性は保存される。

% Markdown AST block 432: Para
\textbf{(b)} は天井関数の定義そのものである。

% Markdown AST block 433: Para
\textbf{(c)}
\[
c=2(bF_t(z)-az)+(b-a)H=2(e-\delta)-2H=2(e-(3t-1)).
\]

% Markdown AST block 434: Para
\textbf{(d)}
\[
\lceil(a(s+bv)-\delta)/b\rceil=\lceil(as-\delta)/b+av\rceil.\qedhere
\]
\end{proof}

% Markdown AST block 435: Proposition 53 (digit words and residues).
\begin{samepage}
\begin{proposition}[数字語と剰余]\leavevmode\label{thm:53}
\(L\ge1\) とする。整数 \(z_0\) に、最初の \(L\) 個の数字 \(e_0\cdots e_{L-1}\)（その \(F_t\) 軌道の数字）を対応させる写像は、全単射
\[
\mathbb Z/b^L\to\{0,\ldots,b-1\}^L.
\]
を誘導する。したがって、奇整数 \(Q_0,Q_0'\) とその \(T_t\) 軌道の繰り上がり語 \(c,c'\) について、
\[
c[0:L]=c'[0:L]\iff Q_0\equiv Q_0'\pmod{b^L},
\]
が成り立つ。特に、一つの軌道に沿って
\[
c[k:k+L]=c[0:L]\iff Q_k\equiv Q_0\pmod{b^L}.
\]
が成り立つ。
\label{thm-end:53}
\end{proposition}
\end{samepage}

% Markdown AST block 436: Proof.
\begin{proof}\leavevmode\label{proof:53}
\(L\) に関する帰納法で示す。\(L=1\) の場合、
\[
e_0\equiv\delta-az_0\pmod b
\]
であり、\(\gcd(a,b)=1\) が成り立つ（\(a=b+2\) で、\(b\) は奇数）。したがって、\(z_0\bmod b\mapsto e_0\) は全単射である。

% Markdown AST block 437: Para
帰納段階では、\(z_0'=z_0+b^{L+1}v\) ならば \(e_0=e_0'\) であり、補題52(d)により
\[
z_1'=F_t(z_0')=z_1+ab^Lv\equiv z_1\pmod{b^L},
\]
となる。したがって、残りの \(L\) 個の数字は帰納法の仮定により一致する（写像は矛盾なく定義される）。

% Markdown AST block 438: Para
逆に、数字語が一致するならば、\(e_0\) から \(z_0\bmod b\) が定まる。
\[
\begin{aligned}
&z_0=s+bq,\qquad z_0'=s+bq',\qquad\text{so}\\
&z_1=F_t(s)+aq,\qquad z_1'=F_t(s)+aq';
\end{aligned}
\]
と書く。帰納法の仮定により
\[
z_1\equiv z_1'\pmod{b^L},
\]
であり、\(a\) は法 \(b^L\) で可逆なので、\(q\equiv q'\pmod{b^L}\)、すなわち
\[
z_0\equiv z_0'\pmod{b^{L+1}}
\]
を得る（単射性）。

% Markdown AST block 439: Para
両方の集合の元の個数は \(b^{L+1}\) であるから、この写像は全単射である。

% Markdown AST block 440: Para
繰り上がりについては、補題52(c)により \(c_n\leftrightarrow e_n\) は全単射である。また、\(b\) は奇数なので、\(Q=2z+H\) から
\[
Q\equiv Q'\pmod{b^L}\iff z\equiv z'\pmod{b^L}.\qedhere
\]
を得る。
\end{proof}

% Markdown AST block 441: Proposition 54 (\(P\)-stage congruence certificate).
\begin{samepage}
\begin{proposition}[\(P\) 段階の合同証明書]\leavevmode\label{thm:54}
\(P>a\) を素数、\(Q_n=T_t^n(P)\) とし、正整数
\(k_1,\ldots,k_P\) が \(s_0=0\)、\(s_i=s_{i-1}+k_i\) および
\[
Q_{k_i}\equiv P\pmod{b^{s_{i-1}}}\qquad(i=1,\ldots,P).
\tag{33}\label{eq:33}
\]
を満たすとする。このとき、\(P\mid Q_n\) がある \(0<n\le s_P\) に対して成り立ち、\(Q_n>P\) であるから、\(Q_n\) は合成数である。
\label{thm-end:54}
\end{proposition}
\end{samepage}

% Markdown AST block 442: Proof.
\begin{proof}\leavevmode\label{proof:54}
命題53を \(L=s_{i-1}\) および \(k=k_i\) として適用すると、(33)から
\[
c[k_i:k_i+s_{i-1}]=c[0:s_{i-1}],
\]
を得る。これは \(t_i=k_i\) および \(L=0\) とした(31)である。

% Markdown AST block 443: Para
\(X=\mathbb Z/P\) 上で、
\[
f_c(x)=b^{-1}(ax+c)
\]
を用いて補題49を適用する（\(P>a>b\) より \(\gcd(P,ab)=1\) なので、これらは置換である）。対象を \(x=0\) とすると、ある
\(n=k_{i+1}+\cdots+k_j\le s_P\)
が
\[
Q_n\equiv G_n(0)=0.
\]
を満たす。

% Markdown AST block 444: Para
\(Q\ge3t\) ならば、
\[
\begin{aligned}
&\lfloor(Q+3t+1)/b\rfloor\ge1\qquad\text{and}\\
&T_t(Q)\ge Q+2;
\end{aligned}
\]
である。\(T_t\) は広義単調増加であり、\(P>a>3t\) なので、軌道は狭義単調増加となり、\(Q_n>P\) が成り立つ。\qedhere
\end{proof}

% Markdown AST block 445: Para
\emph{真の軌道への適用。} \(\lfloor\xi\rfloor=P\) とする。すべての項
\(\lfloor\xi r^n\rfloor\)、\(0\le n\le s_P\) が \(6\) と互いに素ならば、命題55（付録C）により、それらは \(T_t\) 軌道、すなわち初期値 \(P\) の軌道と一致する。命題54により、ある \(n\le s_P\) に合成数の項が存在する。そうでなければ、ある項 \(\ge P>3\) が \(2\) または \(3\) で割り切れ、合成数である。いずれの場合も、時刻 \(\le s_P\) に合成数の項が現れ、区間整合性の検査は必要ない。すべての \(t\) および \(P\) に対してこのような証明書が存在することは、証明されていない。

% Markdown AST block 446: Para
\emph{例。} \(9/7\)（\(t=1\)）および \(\xi=11\) の場合、真の軌道 \(\lfloor11(9/7)^n\rfloor\) の繰り上がりは \(c_0=c_{43}=-1\) を満たす。したがって、コピー条件(31)は \(t_1=1\)、\(t_2=43\) および \(L=0\) として成り立つ。補題49の証明の(1)により、\(G_{44}=G_1\circ G_{43}\) が \(\mathbb Z/11\) 上で成り立ち、剰余が帰還することを直接確認できる。
\[
11\mid Q_{44}=697829=11\cdot63439.
\]
（補題49そのものを適用するには十一個のコピー時刻が必要になるが、ここでは帰還を直接検証している。）

% Markdown AST block 447: Para
この証明書は \(T_1\) 軌道上ではなく真の軌道上にある（\(\lfloor99/7\rfloor=14\ne T_1(11)=15\) であるため）。また、\(P=11>a=9\) は、\(\gcd(11,63)=1\) であるから補題49そのものに対して許容される。

% Markdown AST block 448: Para
\emph{注意（定量的な形）。} \(R(L)\) を接頭語 \(w[0:L]\) の最初の正の帰還時刻とする。\(R(L)\le KL\) がすべての \(L\ge1\) に対して成り立つならば、\(L=1\) のコピーを
\[
t_k=R(S_{k-1}+1)\le
K(S_{k-1}+1),
\]
となるように選べる。したがって、
\[
S_k+1\le(K+1)^k,
\]
となり、法 \(M\) での初期状態への帰還は時刻 \((K+1)^M-1\) までに起こる。この族についてそのような \(K\) は証明されておらず、この評価も用いない。

% Markdown AST block 449: Appendix C. Non-branching families and the exact integer halting problem
\Needspace{6\baselineskip}\section{非分岐族と厳密な整数停止問題}\label{app:C}

% Markdown AST block 450: Para
本付録では族 \(r_t=(6t+3)/(6t+1)\)、\(t\ge1\)、すなわち \(9/7,15/13,21/19,\ldots\) を扱う（\(j=3t-2\) に当たる、系40の族 \((2j+7)/(2j+5)\) の元）。素数 \(2\) および \(3\) を避ける尾部が(32)の整数写像 \(T_t\) に従うこと、また固定した \(t\) に対する全係数の主張が厳密な整数プログラムの停止と同値であることを証明する。さらに、素数の後続項が一意となる別の族を構成する。これらの結果は、この族に対する停止を証明するものではなく、いずれの \(r_t\) についても全係数の定理を主張しない。

% Markdown AST block 451: C.1 The non-branching lemma
\Needspace{6\baselineskip}
\subsection{非分岐補題}\label{sec:C.1}

% Markdown AST block 452: Proposition 55 (non-branching).
\begin{samepage}
\begin{proposition}[非分岐]\leavevmode\label{thm:55}
ここで
\[
\begin{aligned}
&t\ge1,\qquad a=6t+3,\qquad b=6t+1,\\
&x>0,\qquad Q=\lfloor x\rfloor,\qquad Q'=\lfloor rx\rfloor.
\end{aligned}
\]
とする。\(Q\) が奇数で \(\gcd(Q',6)=1\) ならば、\(Q'=T_t(Q)\) である。
\label{thm-end:55}
\end{proposition}
\end{samepage}

% Markdown AST block 453: Proof.
\begin{proof}\leavevmode\label{proof:55}
繰り上がり \(c=bQ'-aQ\) は、(6)により
\[
-6t\le c\le6t+2
\]
を満たす。\(a,b,Q,Q'\) は奇数なので、\(c=2s\) は偶数であり、
\[
-3t\le s\le3t+1.
\]
を満たす。

% Markdown AST block 454: Para
\(3\mid a\)、\(3\nmid b\)、\(3\nmid Q'\) なので、
\[
c\equiv bQ'\not\equiv0\pmod3,
\]
であり、したがって \(3\nmid s\) である。これにより \(s=-3t\) は除かれ、
\[
-3t+1\le s\le3t+1,
\]
が残る。この区間にはちょうど \(b=6t+1\) 個の整数が含まれる。

% Markdown AST block 455: Para
法 \(b\) で \(a\equiv2\) なので、
\(2s=c\equiv-2Q\) および
\[
s\equiv-Q\pmod b;
\]
が成り立つ。この区間には各剰余類の整数がちょうど一つずつ含まれるため、\(s\) は一意である。

% Markdown AST block 456: Para
\(m=\lfloor(Q+3t+1)/b\rfloor\) とおくと、
\[
0\le Q+3t+1-bm\le6t
\]
より
\[
Q-3t+1\le bm\le Q+3t+1,
\]
を得る。したがって、\(s=bm-Q\) はこの区間に属し、\(\equiv-Q\) に等しい。よって
\[
Q'=(aQ+2s)/b=(a-2)Q/b+2m=Q+2m=T_t(Q).\qedhere
\]
である。
\end{proof}

% Markdown AST block 457: Para
ここでは \(Q\) が奇数であることと \(\gcd(Q',6)=1\) だけを用いており、\(3\nmid Q\) は必要ない。結論は候補の一意性である。\(T_t(Q)\) が \(3\) で割り切れる場合、許容される後続項は存在しない。

% Markdown AST block 458: C.2 Conjugation to a ceiling map and digit restriction
\Needspace{6\baselineskip}
\subsection{天井写像への共役と数字制約}\label{sec:C.2}

% Markdown AST block 459: Para
補題52により、写像 \(T_t\)（奇整数 \(Q=2z+H\) 上の写像）は
\[
F_t(z)=\lceil(az-\delta)/b\rceil,
\]
と共役であり、その数字は
\[
e=bF_t(z)-az+\delta\in\{0,\ldots,b-1\}
\]
で、繰り上がりは
\[
c=2(e-(3t-1)).
\]
である。\(t=2u\) が偶数の場合、\(\delta=0\)、\(H=6u-1\) であり、\(F_t(z)=\lceil az/b\rceil\) は通常の天井写像である。

% Markdown AST block 460: Lemma 56 (digit restriction).
\begin{samepage}
\begin{lemma}[数字制約]\leavevmode\label{thm:56}
(a)
\[
Q'=T_t(Q)\equiv2e-1\pmod3;
\]
が成り立つ。したがって、\(3\mid Q'\) であることと \(e\equiv2\pmod3\) であることは同値である。
(b) \(Q\) および \(Q'\) が \(a\) より大きい素数ならば、
\[
e\in D_t=\{0\le e\le6t:\ \gcd(e-(3t-1),(6t+3)(6t+1))=1\}.
\]
が成り立つ。
(c)
\[
d_t=|D_t|\le4t+1<b.
\]
\label{thm-end:56}
\end{lemma}
\end{samepage}

% Markdown AST block 461: Proof.
\begin{proof}\leavevmode\label{proof:56}
(a) \(e\equiv F_t(z)+\delta\pmod3\) である。これは \(b\equiv1\) および \(a\equiv0\) による。また、\(H\equiv-1-\delta\) なので、
\[
Q'=2F_t(z)+H\equiv2(e-\delta)-1-\delta\equiv2e-1\pmod3.
\]
を得る。

% Markdown AST block 462: Para
\textbf{(b)} 補題29により \(\gcd(c,ab)=1\) であり、\(ab\) は奇数なので、
\[
\gcd(e-(3t-1),ab)=1.
\]
が成り立つ。

% Markdown AST block 463: Para
\textbf{(c)} \(3t-1\equiv2\pmod3\) なので、すべての \(e\equiv2\pmod3\) は除外される。該当するものの個数は \(2t\) である。ここでは \(e\) を \([0,6t]\) の中で数えているため、
\[
d_t\le6t+1-2t.\qedhere
\]
を得る。
\end{proof}

% Markdown AST block 464: Table 9.
% Caption is rendered at the start of the immediately following table.

% Markdown AST block 465: Table 9
\begingroup
\small
\setlength{\tabcolsep}{4pt}
\renewcommand{\arraystretch}{1.15}
\setlength{\LTcapwidth}{\textwidth}
\begin{longtable}{@{}>{\raggedright\arraybackslash}p{\dimexpr0.08\dimexpr\linewidth-16pt\relax} >{\raggedright\arraybackslash}p{\dimexpr0.12\dimexpr\linewidth-16pt\relax} >{\raggedright\arraybackslash}p{\dimexpr0.8\dimexpr\linewidth-16pt\relax}@{}}
\caption{\(t\le3\) に対する許容数字集合。}\label{tab:9}\\
\toprule
\(t\) & \(a/b\) & \(D_t\) \\
\midrule
\endfirsthead
\toprule
\(t\) & \(a/b\) & \(D_t\) \\
\midrule
\endhead
\bottomrule
\endlastfoot
1 & \(9/7\) & \(\{0,\allowbreak{}1,\allowbreak{}3,\allowbreak{}4,\allowbreak{}6\}\) \\*
2 & \(15/13\) & \(\{1,\allowbreak{}3,\allowbreak{}4,\allowbreak{}6,\allowbreak{}7,\allowbreak{}9,\allowbreak{}12\}\) \\*
3 & \(21/19\) & \(\{0,\allowbreak{}3,\allowbreak{}4,\allowbreak{}6,\allowbreak{}7,\allowbreak{}9,\allowbreak{}10,\allowbreak{}12,\allowbreak{}13,\allowbreak{}16,\allowbreak{}18\}\) \\
\end{longtable}
\endgroup

% Markdown AST block 466: Para
\emph{注意（実数としての整合性）。} 真の軌道では \(x_n=2z_n+H+\theta_n\) なので、
\[
x_n/2-z_n\in[H/2,(H+1)/2).
\]
である。\(bz_{n+1}=az_n+e_n-\delta\) から、
\[
r^{-(n+1)}z_{n+1}-r^{-n}z_n=r^{-(n+1)}(e_n-\delta)/b
\]
が \(|e_n-\delta|\le b\) を満たす形で得られる。したがって、
\[
\Lambda=\lim r^{-n}z_n
\]
が存在し、
\[
z_n=\Lambda r^n-\sum_{j\ge0}(e_{n+j}-\delta)b^j/a^{j+1};
\]
が成り立つ。さらに、\(z_0>b/2\) ならば \(\Lambda>0\) である。真の軌道に対しては \(\xi=2\Lambda\) が定まる。すべての数字が \(D_t\) に属することからは、幅 \(1/2\) の区間条件も、各項が素数であることも導かれない。

% Markdown AST block 467: Remark.
\begin{remark}
整数上の天井写像 \(x\mapsto\lceil(p/q)x\rceil\) の反復は、\hyperref[ref-aev26]{AEV26, Conjecture 1.6} における法 \(q^k\) の剰余の一様分布に関する予想の対象である。同論文は整数の初期値を扱い、定理ではなく予想を述べている。本論文の内容はこれに依存しない。
\end{remark}

% Markdown AST block 468: C.3 The exact integer halting program
\Needspace{6\baselineskip}
\subsection{厳密な整数停止プログラム}\label{sec:C.3}

% Markdown AST block 469: Para
\(t\ge1\) と初期素数 \(P>a\) を固定する。\((Q,L,U,B)=(P,0,1,1)\) から始め、次を繰り返す。
\[
\begin{aligned}
Q^{+}&=T_t(Q),\quad c=bQ^{+}-aQ,\quad B^{+}=bB,\\
L^{+}&=\max(0,aL-cB),\\
U^{+}&=\min(B^{+},aU-cB),
\end{aligned}
\]
\(Q^{+}\) が合成数、または \(L^{+}\ge U^{+}\) ならば\emph{停止}する。それ以外の場合は \((Q,L,U,B)\) を \((Q^{+},L^{+},U^{+},B^{+})\) に置き換え、続行する。素数判定は厳密に行う。時間、メモリ、または反復回数の上限に達した場合は未判定と記録し、停止とはみなさない。

% Markdown AST block 470: Proposition 57 (exactness).
\begin{samepage}
\begin{proposition}[厳密性]\leavevmode\label{thm:57}
\(n\) 回のステップが受理された後、\(B=b^n\) であり、\([L/B,U/B)\) は、小数部分 \(\theta_n\) の集合に正確に一致する。ここで対象は、\(\xi\in[P,P+1)\) のうち、その軌道の整数部分が \(P=Q_0,\ldots,Q_n\) となるものである。特に、\(L^{+}\ge U^{+}\) は、有限列 \(Q_0,\ldots,Q_{n+1}\) がいかなる実数 \(\xi\) によっても実現されないことを意味する。
\label{thm-end:57}
\end{proposition}
\end{samepage}

% Markdown AST block 471: Proof.
\begin{proof}\leavevmode\label{proof:57}
写像
\[
\theta\mapsto\theta'=(a\theta-c)/b
\]
は単調増加なアフィン全単射であるから、\([L/B,U/B)\) の像は
\[
[(aL-cB)/(bB),(aU-cB)/(bB)),
\]
であり、それと \([0,1)\) との共通部分は
\[
[L^{+}/B^{+},U^{+}/B^{+}).
\]
である。

% Markdown AST block 472: Para
初期状態では、\(\theta_0\) は \([0,1)\) 全体を動く。帰納的に、\(\theta_{n+1}\) が像と \([0,1)\) の共通部分に属することと、それを生じさせる \(\theta_n\) が前の集合に属することは同値である。\qedhere
\end{proof}

% Markdown AST block 473: C.4 Equivalence with the all-coefficient statement
\Needspace{6\baselineskip}
\subsection{全係数の主張との同値性}\label{sec:C.4}

% Markdown AST block 474: Theorem 58 (halting equivalence).
\begin{samepage}
\begin{theorem}[停止との同値性]\leavevmode\label{thm:58}
\(t\ge1\) を固定すると、次の条件は同値である。

% Markdown AST block 475: OrderedList
\begin{enumerate}
\def\labelenumi{(\Alph{enumi})}
\item
  任意の \(\xi>0\) に対し、数列 \(\lfloor\xi r_t^n\rfloor\) は合成数の項を無限に多く含む。
\item
  C.3節のプログラムは、任意の初期素数 \(P>a\) に対して停止する。
\end{enumerate}
\label{thm-end:58}
\end{theorem}
\end{samepage}

% Markdown AST block 476: Proof.
\begin{proof}\leavevmode\label{proof:58}
(B)\(\Rightarrow\)(A)。(A)が成り立たないならば、ある \(\xi\) および \(n_0\) が存在して、\(Q_n\) はすべての \(n\ge n_0\) に対して素数である。(7)により \(n_1\ge n_0\) を \(Q_{n_1}>a\) を満たすように選び、\(\xi\) を \(\xi r^{n_1}\) に置き換える。このとき、\(P=\lfloor\xi\rfloor>a\) は素数であり、すべての項は \(>a>3\) なる素数で、\(6\) と互いに素である。命題55により、
\[
Q_{n+1}=T_t(Q_n)
\]
がすべての \(n\) に対して成り立つ。したがって、プログラムの \(Q\) は真の項と一致し、合成数になることはない。また、命題57により真の \(\theta_n\) はすべての区間に属するので、空の区間はない。よってプログラムは停止せず、(B)に矛盾する。

% Markdown AST block 477: Para
(A)\(\Rightarrow\)(B)。ある \(P>a\) に対してプログラムが停止しないと仮定する。このとき、すべての \(Q_n=T_t^n(P)\) は素数である（合成数でない整数 \(>a\ge9\) は素数である）。また、すべての区間は空でない。したがって、命題57により、各 \(n\) に対して、ある \(\xi_n\in[P,P+1)\) が存在して
\[
\lfloor\xi_nr^j\rfloor=Q_j
\]
が \(0\le j\le n\) について成り立つ。

% Markdown AST block 478: Para
ここで
\[
\begin{aligned}
&I_j=[Q_jr^{-j},(Q_j+1)r^{-j}]\qquad \text{and}\\
&J_n=I_0\cap\cdots\cap I_n.
\end{aligned}
\]
とおく。このとき、
\[
\xi_n\in J_n\ne\varnothing,
\]
および \(J_{n+1}\subseteq J_n\) が成り立ち、各 \(J_n\) はコンパクト区間 \(I_0=[P,P+1]\) の閉部分集合である。Cantorの共通部分定理により \(\bigcap_nJ_n\ne\varnothing\) であり、各幅は高々 \(r^{-n}\to0\) なので、共通部分は一点 \(\xi_*>0\) である。（半開区間では、有限交叉性から共通点は得られない。このため閉区間を用い、端点を別に扱う。）

% Markdown AST block 479: Para
したがって、
\[
Q_n\le\xi_*r^n\le Q_n+1
\]
がすべての \(n\) に対して成り立つ。

% Markdown AST block 480: Para
\(\xi_*r^n=Q_n+1\) がある \(n\) に対して成り立つならば、
\[
\xi_*=(Q_n+1)b^n/a^n
\]
は有理数である。その既約表示を \(\xi_*=u/v\) とすると、\(\xi_*r^m\in\mathbb Z\) には \(vb^m\mid ua^m\) が必要なので、\(v\mid a^m\) および \(b^m\mid u\) となる。これは有限個の \(m\) に対してしか成り立たない。なぜなら \(u\ne0\) だからである。

% Markdown AST block 481: Para
したがって、すべての十分大きな \(n\) に対して、
\[
Q_n\le\xi_*r^n<Q_n+1,
\]
である。すなわち、
\[
Q_n=\lfloor\xi_*r^n\rfloor
\]
は素数であり、(A)に矛盾する。\qedhere
\end{proof}

% Markdown AST block 482: Para
すべての \(t\) に対して(B)が成り立つかは未解決である。\(\xi_*\) の軌道は \(T_t\) 軌道と有限個の添字で異なる可能性があるが、尾部が一致すれば十分である。

% Markdown AST block 483: C.5 Removing prime forks for general bases
\Needspace{6\baselineskip}
\subsection{一般の底に対する素数分岐の除去}\label{sec:C.5}

% Markdown AST block 484: Para
\(a=b+d\)、\(\gcd(a,b)=1\)、\(1<a/b<2\) とする（したがって \(d<b\)）。\(x\in[Q,Q+1)\) に対し、\(\lfloor rx\rfloor\) が取りうる値は、整数 \(m\) のうち、\([m,m+1)\) が \([aQ/b,a(Q+1)/b)\) と交わるものである。後者の区間の長さは \(a/b<2\) なので、候補は高々三つの連続整数からなり、奇整数は高々二つしか含まれない。二つとも現れる場合には、必ず \(L,L+2\) の形である。

% Markdown AST block 485: Lemma 59 (two odd candidates).
\begin{samepage}
\begin{lemma}[二つの奇数候補]\leavevmode\label{thm:59}
\(Q\) が奇数で、\(L\) と \(L+2\) の両方が候補ならば、\(C=L+1\) とおくと
\[
\begin{aligned}
aQ&=bC-k,\quad 1\le k<d,\quad k\equiv a\pmod2,\\
bL&=a(Q-1)+(d+k),\\
b(L+2)&=a(Q+1)-(d-k).
\end{aligned}
\]
が成り立つ。
\label{thm-end:59}
\end{lemma}
\end{samepage}

% Markdown AST block 486: Proof.
\begin{proof}\leavevmode\label{proof:59}
\(L\) が候補であることと、\(aQ/b<L+1\) および \(L<a(Q+1)/b\) は同値である。\(L+2\) が候補であることと、\(aQ/b<L+3\) および \(L+2<a(Q+1)/b\) は同値である。

% Markdown AST block 487: Para
合わせると、\(aQ<bC\) および \(bC+b<aQ+a\)、すなわち
\[
bC-d<aQ<bC,
\]
である。したがって、
\[
k=bC-aQ\in\{1,\ldots,d-1\}.
\]
となる。

% Markdown AST block 488: Para
\(L\) は奇数なので、\(C\) および \(bC\) は偶数である。また、\(aQ\equiv a\pmod2\) であるから、\(k\equiv a\pmod2\) を得る。

% Markdown AST block 489: Para
各等式は
\[
\begin{aligned}
&bL=bC-b=aQ+k-b\qquad \text{and}\\
&b(L+2)=bC+b=aQ+k+b.
\end{aligned}\qedhere
\]
から従う。
\end{proof}

% Markdown AST block 490: Corollary 60 (no prime fork).
\begin{samepage}
\begin{corollary}[素数分岐が存在しないための条件]\leavevmode\label{thm:60}
任意の \(k\) で \(1\le k<d\) および \(k\equiv a\pmod2\) を満たすものに対して
\[
\gcd(k,b)>1,\qquad\text{or}\qquad \gcd(d+k,a)>1,\qquad\text{or}\qquad \gcd(d-k,a)>1,
\]
が成り立つならば、奇素数 \(Q>a\) で、二つの素数候補 \(L,L+2>a\) をもつものは存在しない。さらに、\(L,L+2>3\) が素数ならば \(3\mid C\) なので、素数 \(Q>3\) において分岐が起こるには
\[
3\mid a\iff3\mid k;
\]
が必要である。この同値性を満たさない \(k\) の値も除外される。
\label{thm-end:60}
\end{corollary}
\end{samepage}

% Markdown AST block 491: Proof.
\begin{proof}\leavevmode\label{proof:60}
\(p\mid\gcd(k,b)\) ならば、\(p\mid aQ\)、\(p\nmid a\)、\(p\mid Q\) であり、\(Q\) が素数であることから \(Q=p\le b<a\) が従う。

% Markdown AST block 492: Para
\(p\mid\gcd(d+k,a)\) ならば、\(p\mid bL\)、\(p\nmid b\)、\(p\mid L\)、\(L=p\le a\) である。

% Markdown AST block 493: Para
\(p\mid\gcd(d-k,a)\) ならば、同様に \(p\mid L+2\) である。

% Markdown AST block 494: Para
最後の主張について、\(L,L+1,L+2\) のうち \(3\) の倍数は \(C=L+1\) であり、
\[
aQ=bC-k\equiv-k\pmod3
\]
が \(3\nmid Q\) のもとで成り立つ。\qedhere
\end{proof}

% Markdown AST block 495: Para
\emph{例。} \(a=462u+77\)、\(b=462u+71\)（\(u\ge0\)）の場合を考える。\(d=6\)、\(a\) は奇数、\(k\in\{1,3,5\}\) である。\(7\mid a\) および \(11\mid a\) から、\(\gcd(d+k,a)>1\) が \(k=1,5\) に対して成り立つ。\(k=3\) の場合、\(3\mid k\) だが \(a\equiv2\pmod3\) なので、同値性は成り立たない。これらの組は既約であり（\(\gcd(a,b)\mid6\)、\(a\) は奇数、\(3\nmid a\)）、\(a<2b\) を満たす。したがって、この族では素数 \(Q>a\) の素数後続項は高々一つである。ただし、残る候補が大きい方とは限らない。

% Markdown AST block 496: C.6 Families with a prescribed prime-power difference
\Needspace{6\baselineskip}
\subsection{指定した素数冪の差をもつ族}\label{sec:C.6}

% Markdown AST block 497: Proposition 61 (families for every prime-power difference).
\begin{samepage}
\begin{proposition}[任意の素数冪の差に対する族]\leavevmode\label{thm:61}
\(d=p^h\ge2\) とし、
\[
K_d=\{k:1\le k<d,\ k\equiv d+1\pmod2\}.
\]
とおく。各 \(k\in K_d\) に対して、素数 \(\ell_k\ne p\) で \(d+k\) を割り切るものが存在する。\(P_d\) を、相異なる \(\ell_k\) の積とする。このとき、\(\gcd(P_d,2d)=1\) である。また、\(a>2d\) が
\[
a\equiv0\pmod{P_d}\qquad \text{and}\qquad a\equiv d+1\pmod{2d}
\]
（法 \(2dP_d\) の等差数列）および \(b=a-d\) を満たすならば、\(\gcd(a,b)=1\)、\(1<a/b<2\)、および系60の仮定を、
\[
\gcd(d+k,a)>1
\]
がすべての該当する \(k\) について成り立つという形で満たす。
\label{thm-end:61}
\end{proposition}
\end{samepage}

% Markdown AST block 498: Proof.
\begin{proof}\leavevmode\label{proof:61}
\(k\in K_d\) に対して、\(d+k\equiv2d+1\) は奇数であり、\(d<d+k<2d\) である。\(p\) の冪は \((d,2d)\) 内に存在しないため、\(d+k\ge3\) は素因数 \(\ell_k\ne p\) をもち、その素因数は必ず奇数である。

% Markdown AST block 499: Para
したがって、\(P_d\) は奇数で、\(p\) と互いに素である。よって \(\gcd(P_d,2d)=1\) となり、二つの合同式は両立する。

% Markdown AST block 500: Para
\(a\equiv1\pmod d\) より、
\[
\gcd(a,b)=\gcd(a,d)=1;
\]
および
\[
a>2d\iff a<2b;
\]
を得る。また、\(a\equiv d+1\pmod2\) によって \(k\equiv a\pmod2\) と \(k\in K_d\) は同値となり、
\[
\ell_k\mid\gcd(d+k,a).\qedhere
\]
が成り立つ。
\end{proof}

% Markdown AST block 501: Table 10.
% Caption is rendered at the start of the immediately following table.

% Markdown AST block 502: Table 10
\begingroup
\small
\setlength{\tabcolsep}{4pt}
\renewcommand{\arraystretch}{1.15}
\setlength{\LTcapwidth}{\textwidth}
\begin{longtable}{@{}>{\raggedright\arraybackslash}p{\dimexpr0.08\dimexpr\linewidth-24pt\relax} >{\raggedright\arraybackslash}p{\dimexpr0.3\dimexpr\linewidth-24pt\relax} >{\raggedright\arraybackslash}p{\dimexpr0.3\dimexpr\linewidth-24pt\relax} >{\raggedright\arraybackslash}p{\dimexpr0.32\dimexpr\linewidth-24pt\relax}@{}}
\caption{小さい \(d\) に対して命題61が与える族。}\label{tab:10}\\
\toprule
\(d\) & \(a\) & \(b\) & 範囲 \\
\midrule
\endfirsthead
\toprule
\(d\) & \(a\) & \(b\) & 範囲 \\
\midrule
\endhead
\bottomrule
\endlastfoot
2 & \(12u+3\) & \(12u+1\) & \(u\ge1\) \\*
3 & \(30u+10\) & \(30u+7\) & \(u\ge0\) \\*
4 & \(280u+245\) & \(280u+241\) & \(u\ge0\) \\*
5 & \(210u+126\) & \(210u+121\) & \(u\ge0\) \\
\end{longtable}
\endgroup

% Markdown AST block 503: Para
\(d\) が奇数の場合、分子 \(a\) は偶数である。補題59は偶奇条件 \(k\equiv a\pmod2\) によってこれを扱う。これらは素数後続項について一様な規則をもつ族であり、合成数の出現が証明された族ではない。

% Markdown AST block 504: Proposition 62 (ceiling form).
\begin{samepage}
\begin{proposition}[天井関数による表示]\leavevmode\label{thm:62}
命題61の設定のもとで、\(b=2dh+1\)、\(H=2h-1\)、\(Q=2z+H\) と書く。\(r(Q+1)\) 未満の最大の奇整数は \(2\lceil az/b\rceil+H\) である。したがって、\(Q\) の後続項に奇数の候補が存在するならば、その最大値は
\[
Q'=2\lceil az/b\rceil+H,
\]
であり、素数後続項 \(>a\) が素数 \(Q>a\) に対して存在する場合には、この値に限られる。\(z'=\lceil az/b\rceil\) および \(e=bz'-az\in[0,b)\) とおくと、
\[
\begin{aligned}
\gcd(Q,b)&=\gcd(2e+d+1,b),\\
\gcd(Q',a)&=\gcd(2e+2d+1,a),
\end{aligned}
\]
となる。したがって、素数 \(>a\) の尾部では、すべての数字が
\[
D_{a,b}=\{0\le e<b:\gcd(2e+d+1,b)=\gcd(2e+2d+1,a)=1\}.
\]
に属する。
\label{thm-end:62}
\end{proposition}
\end{samepage}

% Markdown AST block 505: Proof.
\begin{proof}\leavevmode\label{proof:62}
\(y=(Q+1)/2=z+h\) とおく。\(2ry\) 未満の最大の奇整数は \(2y'-1\) であり、\(y'=\lceil ry-1/2\rceil\) を満たす。

% Markdown AST block 506: Para
\(b\) は奇数なので、\(ry\) は半整数ではない。したがって、
\[
y'=\lfloor ry+1/2\rfloor
=\lfloor(2ay+b)/(2b)\rfloor=\lfloor(ay+(b-1)/2)/b\rfloor.
\]
となる。

% Markdown AST block 507: Para
\((b-1)/2=dh\) および \(a=b+d\) とおくと、
\[
ay+dh=az+bh+(b-1),
\]
であるから、
\[
\begin{aligned}
&y'=h+\lfloor(az+b-1)/b\rfloor=h+\lceil az/b\rceil\qquad\text{and}\\
&Q'=2y'-1=2\lceil az/b\rceil+H.
\end{aligned}
\]
を得る。

% Markdown AST block 508: Para
最大公約数の等式については、
\[
dQ=2dz+b-1-d\equiv-2e-1-d\pmod b
\]
（\(2e\equiv-2az\equiv-2dz\) による）であり、さらに \(\gcd(d,b)=1\) である。同様に、\(dQ'=2dz'+b-1-d\) に加えて \(bz'\equiv e\pmod a\) および \(b\equiv-d\pmod a\) を用いると、
\[
dQ'\equiv-(2e+2d+1)\pmod a.\qedhere
\]
を得る。
\end{proof}

% Markdown AST block 509: Para
これらの族に対し、十分大きいすべての初期値が最終的に数字制約または区間条件を満たさなくなることを示す、一様な停止定理は証明されていない。

% Markdown AST block 510: Appendix D. Prime-prefix counting, the word sieve, and finite height
\Needspace{6\baselineskip}\section{素数接頭列の計数、語の篩、および有限高さ}\label{app:D}

% Markdown AST block 511: Para
本付録では定理47を証明する。D.1節では、軌道が \(N\) ステップにわたって素数を保つ初期素数の個数を評価する。D.2節では、一般の底に対する繰り上がり語の篩と、それにより得られる冪型の上界を述べる。D.3--D.4節では、この篩を用いて、\(6/5\) および \(9/7\) に対する明示した高さ以下のすべての初期整数を完全に扱う。いずれの結果も全係数に対する定理ではない。

% Markdown AST block 512: D.1 A counting bound along the family
\Needspace{6\baselineskip}
\subsection{族に沿った個数の上界}\label{sec:D.1}

% Markdown AST block 513: Para
\(t\ge1\)、\(N\ge1\) および \(X\ge a=6t+3\) に対して、\(S_t(X,N)\) を、素数 \(a<P\le X\) のうち、C.3節のプログラムが \(N\) ステップを受理するものの集合とする。

% Markdown AST block 514: Proposition 63.
\begin{samepage}
\begin{proposition}\leavevmode\label{thm:63}
\[
|S_t(X,N)|\le d_t^N(\lfloor X/(2b^N)\rfloor+1),
\]
が成り立つ。ここで \(d_t=|D_t|\le4t+1\) である。\(N=N(X)\) が
\(b^N\ge X\) を満たす最小の整数ならば、
\[
|S_t(X,N(X))|\le d_t\,X^{\alpha_t}
\]
であり、ここで
\[
\alpha_t=\log d_t/\log b<1.
\]
である。同じ上界は、プログラムが決して停止しない初期素数の集合に対しても成り立つ。
\label{thm-end:63}
\end{proposition}
\end{samepage}

% Markdown AST block 515: Proof.
\begin{proof}\leavevmode\label{proof:63}
各 \(P\in S_t(X,N)\) は、長さ \(N\) で、すべての数字が \(D_t\) に属する数字語を定める（補題56）。したがって、それは高々 \(d_t^N\) 個の語のいずれかである。命題53により、その語は \(z_0\bmod b^N\)、したがって \(P=2z_0+H\) の法 \(2b^N\) での剰余を定める。法 \(2b^N\) の一つの剰余類は、高々 \(\lfloor X/(2b^N)\rfloor+1\) 個の元を \((a,X]\) の中にもつ。

% Markdown AST block 516: Para
\(N=N(X)\) の場合、\(2b^N>X\) より各語に対応する元は高々一つである。また、\(b^{N-1}<X\) より
\[
d_t^N\le d_t\cdot d_t^{\log_bX}
=d_tX^{\alpha_t};
\]
を得る。\(\alpha_t<1\) である。これは \(d_t<b\) による。停止しない初期素数は、\(S_t(X,N)\) に任意の \(N\) に対して属する。\qedhere
\end{proof}

% Markdown AST block 517: Para
\(CX^\alpha\) 型で \(\alpha<1\) を満たす上界から、集合が空であることは導かれない。単一の停止しない軌道の尾部でも、開始点は \(O(\log X)\) 個しか \(X\) 以下に存在しない。

% Markdown AST block 518: D.2 The word sieve for a general base
\Needspace{6\baselineskip}
\subsection{一般の底に対する語の篩}\label{sec:D.2}

% Markdown AST block 519: Para
\begin{samepage}\leavevmode\label{layout:carry-intro-start}
\(a>b\ge2\) は互いに素とし、\(c_0\cdots c_{\ell-1}\) を繰り上がり語とする。\(C_0=0\)、\(C_{j+1}=aC_j+b^jc_j\) とおくと、(8)により
\[
b^jQ_j=a^jQ_0+C_j
\]
が成り立つ。
\label{layout:carry-intro-end}\end{samepage}

% Markdown AST block 520: Lemma 64.
\begin{samepage}
\begin{lemma}\leavevmode\label{thm:64}
(i) \emph{整数性。} 合同式
\[
a^\ell Q_0\equiv-C_\ell\pmod{b^\ell}
\]
は \(Q_0\bmod b^\ell\) を定め、
\[
b^j\mid a^jQ_0+C_j
\]
がすべての \(j\le\ell\) について成り立つことを意味する。したがって、\(X\le b^\ell\) の場合、与えられた語をもつ整数 \(Q_0\in[1,X]\) は高々一つである。
(ii) \emph{厳密な区間。}
\[
\theta_j=(a^j\theta_0-C_j)/b^j,
\]
とおくと、この語がある実数軌道片によって実現されることと、
\[
[0,1)\cap\bigcap_{j=1}^{\ell}\Bigl[\frac{C_j}{a^j},\frac{C_j+b^j}{a^j}\Bigr)\ne\varnothing .
\]
は同値である。
(iii) \emph{固定約数。} 素数 \(p\nmid ab\) に対して、\(p\mid Q_j\) であることと \(Q_0\equiv R_j\pmod p\) は同値である。ここで、\(R_0=0\) とし、
\[
R_{j+1}=R_j-c_jb^ja^{-j-1}\pmod p.
\]
とおく。
\[
\{R_0,\ldots,R_\ell\}=\mathbb F_p,
\]
ならば、すべての初期整数に対して \(p\mid Q_j\) がある \(j\le\ell\) について成り立つ。根は高々 \(\ell+1\) 個であるから、素数 \(p\le\ell+1\) だけが \(\mathbb F_p\) を覆いうる。
\label{thm-end:64}
\end{lemma}
\end{samepage}

% Markdown AST block 521: Proof.
\begin{proof}\leavevmode\label{proof:64}
(i)
\[
C_\ell=a^{\ell-j}C_j+\sum_{i\ge j}a^{\ell-1-i}b^ic_i
\equiv a^{\ell-j}C_j\pmod{b^j},
\]
なので、\(b^\ell\mid a^\ell Q_0+C_\ell\) より
\[
b^j\mid a^{\ell-j}(a^jQ_0+C_j)
\]
を得る。また、\(\gcd(a,b)=1\) であるから、
\[
b^j\mid a^jQ_0+C_j.
\]
となる。

% Markdown AST block 522: Para
\textbf{(ii)} \(\theta_j\in[0,1)\) であることと
\[
\theta_0\in[C_j/a^j,(C_j+b^j)/a^j);
\]
は同値である。写像 \(\theta_0\mapsto\theta_j\) は、命題57と同様に厳密である。

% Markdown AST block 523: Para
\textbf{(iii)}
\[
p\mid Q_j\iff a^jQ_0+C_j\equiv0\iff Q_0\equiv-C_ja^{-j},
\]
\[
\text{and}\quad
-(aC_j+b^jc_j)a^{-j-1}=R_j-c_jb^ja^{-j-1}.\qedhere
\]
\end{proof}

% Markdown AST block 524: Para
\(A_\ell\) を、長さ \(\ell\) の語のうち、(1)すべての文字が \(\mathcal C(a,b)\) に属し、(2)(ii)の区間が空でなく、(3)素数 \(p\le\ell+1\) で \(p\nmid2ab\) を満たすもので \(\{R_0,\ldots,R_\ell\}=\mathbb F_p\) となるものが存在しない、という条件を満たすものの個数とする。\(Q_0,\ldots,Q_\ell\) が \(\max(a,\ell+1)\) より大きい素数ならば、その軌道片の繰り上がり語は、補題29により(1)を、真の \(\theta_0\) を用いて(2)を満たす。また、覆う素数があれば、ある \(Q_j\) を割り切って \(Q_j=p\le\ell+1\) を強制するため、(3)も満たす。固定約数が存在しないことは必要条件にすぎず、素数組に関する予想を用いてこれを十分条件とすることはしない。

% Markdown AST block 525: Proposition 65 (power-type bound).
\begin{samepage}
\begin{proposition}[冪型の上界]\leavevmode\label{thm:65}
\(\ell\) を固定し、\(\lambda=A_\ell^{1/\ell}\) および \(\alpha=\log\lambda/\log b\) とおく。
\[
A_s\le C\lambda^{s-1}
\]
が \(0\le s<\ell\) に対して成り立つと仮定する。このとき、素数 \(P\) で \(\max(a,\ell+1)<P\le X\) を満たし、ある \(\xi\in[P,P+1)\) が \(N(X)+1\) 個の素数項を生じさせるものの個数は、高々 \(CX^\alpha\) である。ここで、\(N(X)\) は \(b^{N(X)}\ge X\) を満たす最小の整数である。
\label{thm-end:65}
\end{proposition}
\end{samepage}

% Markdown AST block 526: Proof.
\begin{proof}\leavevmode\label{proof:65}
\(N=k\ell+s\) と書き、\(0\le s<\ell\) とする。すべての \(N+1\) 項が \(\max(a,\ell+1)\) より大きい素数である軌道片の繰り上がり語は、\(k\) 個の長さ \(\ell\) のブロックと、長さ \(s\) の一つのブロックに分かれる。各ブロックは、その開始点での小数部分を初期値として(1)--(3)を満たす。ブロック間の整合性を無視すると過大計数になるため、語の個数は高々
\[
A_\ell^kA_s\le
C\lambda^{N-1}
\]
である。

% Markdown AST block 527: Para
\(N=N(X)\) の場合、\(b^{N-1}<X\) より
\[
\lambda^{N-1}=b^{\alpha(N-1)}<X^\alpha,
\]
を得る。また、補題64(i)により、各語は \([1,X]\) 内に高々一つの初期整数をもつ。異なる \(P\) は異なる語を与える。\qedhere
\end{proof}

% Markdown AST block 528: Para
仮定 \(A_s\le C\lambda^{s-1}\) は、整数不等式 \(A_\ell\le C^\ell\)（\(s=0\)）および \(A_s^{\ell}\le C^{\ell}A_\ell^{s-1}\)（\(1\le s<\ell\)）と同値である。表11は、D.5節の方法で計算した二つの底についての完全な個数を示す。指数は
\[
18035747<7^9=40353607\qquad \text{and}\qquad 3955763<5^{10}=9765625,
\]
によって認証される。すなわち、\(\alpha<1/2\) が \(9/7\) に対して、\(\alpha<1/4\) が \(6/5\) に対して成り立つということである。定数 \(C=13\) および \(C=10\) は、不等式を満たす最小の整数である。\(6/5\) の場合、アルファベットは \(\mathcal C(6,5)=\{-1,1\}\) であり、素数後続項が存在するには \(Q\equiv1\pmod5\)（このとき \(c=-1\)）、または \(Q\equiv4\pmod5\)（このとき \(c=1\)）が必要である。これは \(c\equiv-Q\pmod5\) による。\(9/7\) の場合、アルファベットは \(\mathcal C(9,7)=\{-4,-2,2,4,8\}\) である。

% Markdown AST block 529: Table 11.
% Caption is rendered at the start of the immediately following table.

% Markdown AST block 530: Table 11
\begingroup
\small
\setlength{\tabcolsep}{4pt}
\renewcommand{\arraystretch}{1.15}
\setlength{\LTcapwidth}{\textwidth}
\begin{longtable}{@{}>{\raggedright\arraybackslash}p{\dimexpr0.08\dimexpr\linewidth-32pt\relax} >{\raggedright\arraybackslash}p{\dimexpr0.06\dimexpr\linewidth-32pt\relax} >{\raggedright\arraybackslash}p{\dimexpr0.65\dimexpr\linewidth-32pt\relax} >{\raggedright\arraybackslash}p{\dimexpr0.06\dimexpr\linewidth-32pt\relax} >{\raggedright\arraybackslash}p{\dimexpr0.15\dimexpr\linewidth-32pt\relax}@{}}
\caption{篩の個数 \(A_s\)（\(0\le s\le\ell\)）。}\label{tab:11}\\
\toprule
底 & \(\ell\) & \(A_0,\ldots,A_\ell\) & \(C\) & 範囲 \\
\midrule
\endfirsthead
\toprule
底 & \(\ell\) & \(A_0,\ldots,A_\ell\) & \(C\) & 範囲 \\
\midrule
\endhead
\bottomrule
\endlastfoot
\(9/7\) & 18 & 1,\allowbreak{} 5,\allowbreak{} 19,\allowbreak{} 67,\allowbreak{} 199,\allowbreak{} 530,\allowbreak{} 1315,\allowbreak{} 3130,\allowbreak{} 7250,\allowbreak{} 16511,\allowbreak{} 37155,\allowbreak{} 82942,\allowbreak{} 183838,\allowbreak{} 404586,\allowbreak{} 883163,\allowbreak{} 1909362,\allowbreak{} 4083475,\allowbreak{} 8631875,\allowbreak{} 18035747 & 13 & \(P>19\) \\*
\(6/5\) & 40 & 1,\allowbreak{} 2,\allowbreak{} 4,\allowbreak{} 8,\allowbreak{} 15,\allowbreak{} 28,\allowbreak{} 50,\allowbreak{} 79,\allowbreak{} 124,\allowbreak{} 191,\allowbreak{} 285,\allowbreak{} 417,\allowbreak{} 601,\allowbreak{} 864,\allowbreak{} 1233,\allowbreak{} 1757,\allowbreak{} 2493,\allowbreak{} 3522,\allowbreak{} 4950,\allowbreak{} 6965,\allowbreak{} 9717,\allowbreak{} 13582,\allowbreak{} 18786,\allowbreak{} 26085,\allowbreak{} 35937,\allowbreak{} 49661,\allowbreak{} 67971,\allowbreak{} 93166,\allowbreak{} 126734,\allowbreak{} 172576,\allowbreak{} 233228,\allowbreak{} 315696,\allowbreak{} 423604,\allowbreak{} 568869,\allowbreak{} 757642,\allowbreak{} 1010562,\allowbreak{} 1335866,\allowbreak{} 1766366,\allowbreak{} 2315828,\allowbreak{} 3038864,\allowbreak{} 3955763 & 10 & \(P>41\) \\
\end{longtable}
\endgroup

% Markdown AST block 531: D.3 Recovering the initial integer
\Needspace{6\baselineskip}
\subsection{初期整数の復元}\label{sec:D.3}

% Markdown AST block 532: Lemma 66.
\begin{samepage}
\begin{lemma}\leavevmode\label{thm:66}
\(N\ge2\) とし、\(c\) を長さ \(N\) の語とする。
(a) \emph{（初期値。）}
\[
s(c)=(-C_N(a^N)^{-1})\bmod b^N\in[0,b^N).
\]
とおく。素数 \(P\le b^N\) で、その軌道の繰り上がり語が \(c\) となるものは、すべて \(s(c)\) に等しい。特に、\(s(c)\) が偶数、または \(s(c)\le\max(a,N+1)\) ならば、そのような素数は \(\max(a,N+1)\) を超えない。
(b) \emph{（区間。）} この語が \(\lfloor\xi\rfloor=s(c)\) を満たす実数軌道片によって実現されることと、補題64(ii)の区間が空でないことは同値である。
(c) \emph{（根。）} 素数 \(p\le N+1\) が \(p\nmid2ab\) を満たし、\(\mathbb F_p\) を覆う場合、または
\[
s(c)\bmod p\in\{R_0,\ldots,R_N\},
\]
の場合、初期整数 \(P=s(c)>N+1\) と繰り上がり語 \(c\) をもつすべての軌道には、項 \(Q_j\) で \(p\mid Q_j\)、\(Q_j\ge P>p\) を満たすものが存在し、したがって合成数となる。
\label{thm-end:66}
\end{lemma}
\end{samepage}

% Markdown AST block 533: Proof.
\begin{proof}\leavevmode\label{proof:66}
(a) 補題64(i)により \(P\equiv s(c)\pmod{b^N}\) である。\(0<P\le b^N\) であり、\(b^N\) は合成数なので、\(P=s(c)\) となる（また、\(s(c)=0\) ならば \(P=b^N\) となり、素数ではない）。

% Markdown AST block 534: Para
\textbf{(b)} は補題64(ii)である。

% Markdown AST block 535: Para
\textbf{(c)} は補題64(iii)と \(Q_j\ge Q_0=P\) から従う。\qedhere
\end{proof}

% Markdown AST block 536: D.4 Proof of Theorem 47
\Needspace{6\baselineskip}
\subsection{定理47の証明}\label{sec:D.4}

% Markdown AST block 537: Proof of Theorem 47.
\begin{proof}[定理47の証明]\leavevmode\label{proof:47}
(a) \(6<P=\lfloor\xi\rfloor\le5^{40}\) とし、\(Q_0,\ldots,Q_{40}\) のどの項も合成数ではないと仮定する。すべて \(\ge P\ge7\) なので、すべて素数である。\(N=40\) とおく。\(P\le41=\max(6,N+1)\) ならば \(P\in\{7,11,\ldots,41\}\) であり、これら十個の素数は個別に除外される（表13）。

% Markdown AST block 538: Para
\(P\ge43\) ならば、すべての項は \(\max(a,N+1)\) より大きい素数である。したがって、繰り上がり語 \(c\) は長さ \(40\) で、\(a/b=6/5\) としたD.2節の(1)--(3)を満たす。また、補題66(a)により \(P=s(c)\) であり、これは奇数かつ \(41\) より大きい。補題66(c)により、
\[
s(c)\bmod p\notin\{R_j\}
\]
がすべての素数 \(p\in\{7,11,\ldots,41\}\) に対して成り立つ。

% Markdown AST block 539: Para
(1)--(3)を満たすすべての \(A_{40}=3955763\) 個の語を完全に列挙すると（表12）、初期値が偶数であるものが \(1978668\) 個、根判定で棄却されるものが \(1977095\) 個である。初期値が \(\le41\) となるものはなく、生き残るものもない。この矛盾により(a)が証明された。

% Markdown AST block 540: OrderedList
\begin{enumerate}
\def\labelenumi{(\alph{enumi})}
\setcounter{enumi}{1}
\item
  \(9<P\le7^{18}\) とし、\(Q_0,\ldots,Q_{18}\) のどの項も合成数ではないと仮定する。\(N=18\) とおくと、素数 \(P\le19=\max(9,N+1)\)、すなわち \(11,13,17,19\) は個別に除外される（表13）。

  \(P\ge23\) の場合、\(a/b=9/7\)、素数 \(p\in\{5,11,13,17,19\}\) を用いて \(A_{18}=18035747\) 個の語に同じ篩を適用すると、偶数の初期値が \(9018851\) 個、根による棄却が \(9012304\) 個となる。初期値が \(\le19\) となるものはなく、\(4592\) 個の候補 \(P\) が生き残る。

  生き残った各候補では、すべての項が素数で、したがって \(6\) と互いに素である。命題55により、軌道は
  \[
  T(Q)=Q+2\lfloor(Q+4)/7\rfloor
  \]
  （写像 \(T_1\)）に従う。この \(T\) 軌道の繰り上がりは篩の語と一致する。実際、繰り上がりは
  \[
  c\equiv-2Q\pmod7,
  \]
  によって定まり、文字 \(-4,-2,2,4,8\) は相異なる剰余 \(3,5,2,4,1\) を法 \(7\) でもつ。

  \(4592\) 個の各候補は、ある時刻 \(j\le7\) に真の約数をもつ。すなわち、整数 \(1<d<T^j(P)\) で \(d\mid T^j(P)\) を満たすものが存在する。記録された証拠で用いる約数は高々 \(9923\) であり、それぞれ厳密な等式
  \[
  T^j(P)=d\cdot e.
  \]
  として検査されている。したがって、\(Q_j\) はある \(j\le7\le18\) に対して合成数となり、矛盾する。\qedhere
\end{enumerate}
\end{proof}

% Markdown AST block 541: Table 12.
% Caption is rendered at the start of the immediately following table.

% Markdown AST block 542: Table 12
\begingroup
\small
\setlength{\tabcolsep}{2.5pt}
\renewcommand{\arraystretch}{1.15}
\setlength{\LTcapwidth}{\textwidth}
\begin{longtable}{@{} r r r r r r r r@{}}
\caption{有限高さの篩。「語数」は \(A_N\) を表す。「偶数」および「根棄却」は、初期値が偶数である語、および補題66(c)で棄却される語の数を表す。「小初期値」は初期値 \(\le\max(a,N+1)\) の数を表す。}\label{tab:12}\\
\toprule
\shortstack[r]{底} & \shortstack[r]{\(N\)} & \shortstack[r]{高さ\\\(b^N\)} & \shortstack[r]{語数} & \shortstack[r]{偶数\\初期値} & \shortstack[r]{根による\\棄却} & \shortstack[r]{小さい\\初期値} & \shortstack[r]{残存数} \\
\midrule
\endfirsthead
\toprule
\shortstack[r]{底} & \shortstack[r]{\(N\)} & \shortstack[r]{高さ\\\(b^N\)} & \shortstack[r]{語数} & \shortstack[r]{偶数\\初期値} & \shortstack[r]{根による\\棄却} & \shortstack[r]{小さい\\初期値} & \shortstack[r]{残存数} \\
\midrule
\endhead
\bottomrule
\endlastfoot
\(6/5\) & 40 & \(5^{40}\) & 3955763 & 1978668 & 1977095 & 0 & 0 \\*
\(9/7\) & 18 & \(7^{18}\) & 18035747 & 9018851 & 9012304 & 0 & 4592 \\
\end{longtable}
\endgroup

% Markdown AST block 543: Para
\(4592\) 個の残存候補について、約数を用いる時刻の最大値は \(7\) である。候補ごとのこれらの時刻の分布は試し割りの順序と範囲に依存するため、報告しない。証明で用いる主張は、各候補に対してある時刻 \(j\le7\) に真の約数が存在することである。

% Markdown AST block 544: Table 13.
% Caption is rendered at the start of the immediately following table.

% Markdown AST block 545: Table 13
\begingroup
\small
\setlength{\tabcolsep}{4pt}
\renewcommand{\arraystretch}{1.15}
\setlength{\LTcapwidth}{\textwidth}
\begin{longtable}{@{}>{\raggedright\arraybackslash}p{\dimexpr0.1\dimexpr\linewidth-24pt\relax} >{\raggedright\arraybackslash}p{\dimexpr0.08\dimexpr\linewidth-24pt\relax} >{\raggedright\arraybackslash}p{\dimexpr0.67\dimexpr\linewidth-24pt\relax} >{\raggedright\arraybackslash}p{\dimexpr0.15\dimexpr\linewidth-24pt\relax}@{}}
\caption{小さい初期素数。\(x\in[P,P+1)\) に対する後続項
\(\lfloor rx\rfloor\) の候補は、\(6/5\) では高々二つ、\(9/7\) では高々三つである（C.5節）。各行には、必要な場合にすべての候補を列挙した。\(6/5\) の場合、「\(Q\equiv s\pmod5\)」で \(s\notin\{1,4\}\) とは、\(Q\) の後続候補に \(30\) と互いに素なものがないことを意味する（D.2節）。したがって、どの候補も少なくとも \(8\) であるため合成数である。\(30\) と互いに素な候補が一つだけ存在する場合は、その値を示す。「空区間」とは、全項が素数である \(T\) 軌道について、C.3節の厳密な区間が指定したステップで空になることを意味する。したがって、この軌道はいかなる \(\xi\) によっても実現されず、一方、この軌道から外れる場合は項が \(2\) または \(3\) で割り切れるため、合成数となる。}\label{tab:13}\\
\toprule
底 & \(P\) & 証拠 & 合成数出現時刻の上限 \\
\midrule
\endfirsthead
\toprule
底 & \(P\) & 証拠 & 合成数出現時刻の上限 \\
\midrule
\endhead
\bottomrule
\endlastfoot
\(6/5\) & 7 & \(7\equiv2\pmod5\) & 1 \\
\(6/5\) & 11 & \(30\) と互いに素な候補は \(13\) のみ。\(13\equiv3\pmod5\) & 2 \\
\(6/5\) & 13 & \(13\equiv3\pmod5\) & 1 \\
\(6/5\) & 17 & \(17\equiv2\pmod5\) & 1 \\
\(6/5\) & 19 & \(30\) と互いに素な候補は \(23\) のみ。\(23\equiv3\pmod5\) & 2 \\
\(6/5\) & 23 & \(23\equiv3\pmod5\) & 1 \\
\(6/5\) & 29 & 両候補 \(34\) と \(35=5\cdot7\) は合成数 & 1 \\
\(6/5\) & 31 & \(30\) と互いに素な候補は \(37\) のみ。\(37\equiv2\pmod5\) & 2 \\
\(6/5\) & 37 & \(37\equiv2\pmod5\) & 1 \\
\(6/5\) & 41 & 両候補 \(49=7\cdot7\) と \(50\) は合成数 & 1 \\
\(9/7\) & 11 & 両候補 \(14\) と \(15=3\cdot5\) は合成数 & 1 \\
\(9/7\) & 13 & ステップ4の後に空区間 (\(L=2450>U=2401\),\allowbreak{} \(B=7^4\)) & 4 \\
\(9/7\) & 17 & ステップ3の後に空区間 (\(L=350>U=343\),\allowbreak{} \(B=7^3\)) & 3 \\
\(9/7\) & 19 & 両候補 \(24\) と \(25=5\cdot5\) は合成数 & 1 \\
\end{longtable}
\endgroup

% Markdown AST block 546: Para
定理47は \(\lfloor\xi\rfloor>5^{40}\) または \(\lfloor\xi\rfloor>7^{18}\) について何も述べない。初期整数 \(s+kb^N\) が最初の \(N\) 個について \(s\) と同じ繰り上がりをもつ場合、各項は
\[
Q_j(s)+ka^jb^{N-j},
\]
であり、\(k=0\) に対して見つかった約数が、\(k\ne0\) に対する項も割り切るとは限らない。

% Markdown AST block 547: D.5 Verification of the finite-height computation
\Needspace{6\baselineskip}
\subsection{有限高さ計算の検証}\label{sec:D.5}

% Markdown AST block 548: Para
表11--13の個数と \(4592\) 個の残存候補の一覧は、9.3節の実装Aによって生成した。これは研究ノートに同梱された同系統の二つの構成であり、厳密な区間を異なる座標で追跡する。一方は現在の小数区間を用い、初期値を一桁ずつ持ち上げる。他方は共通分母 \(a^N\) をもつ初期区間を用い、葉で初期値を解く。両者は全深さの個数、全分類の個数、および残存候補一覧で一致する。また、六つの小規模事例
\[
(a,b,N)\in\{(3,2,5),(4,3,5),(5,3,5),
(6,5,6),(7,5,5),(9,7,5)\}
\]
では、\(9,11,37,50,54,530\) 個の語について、直積と厳密な有理数演算により全語を列挙するPythonオラクルとも一致する。表11の篩の個数は全範囲で再計算した。有限高さの計算とその再現手順は \href{supplement/COMPUTATION.md}{計算ガイド} に記載している。

% Markdown AST block 549: Para
これらとは独立に、9.3節の実装B（別のプログラマが、提供されたプログラムではなく補題64--66の仕様から記述したコード）は、両方の底について篩を再列挙した。初期値、根による棄却、残存候補一覧を再計算し、各残存候補の約数等式と小さい素数に関する各証拠を検査し、表11の定数を自身の個数から最小の整数として再計算した。要素ごとに比較したすべての量（個数配列、分類別個数、残存候補集合、\(4592\) 個の証拠等式、小さい素数に関する十四個の証拠、六つの小規模事例、および定数）は一致した。比較の説明は \href{supplement/COMPUTATION.md}{計算ガイド}、参照データは \href{checks/reference.json}{checks/reference.json} に記載している。この一致は、両実装が同じ有限対象を計算したことの証拠である。それらの対象から定理47へ進む論拠はD.2--D.4節に示した。

% Markdown AST block 550: Appendix E. Long prime prefixes and return-word obstructions
\Needspace{6\baselineskip}\section{長い素数接頭列と帰還語の障害}\label{app:E}

% Markdown AST block 551: E.1 Arbitrarily long prime prefixes along the family
\subsection{族に沿った任意に長い素数接頭列}\label{sec:E.1}

% Markdown AST block 552: Lemma 67 (transplanting a prime progression).
\begin{samepage}
\begin{lemma}[素数等差数列の移植]\leavevmode\label{thm:67}
\(L\ge3\) とし、\(P,P+D,\ldots,P+(L-1)D\) は \(3\) より大きい素数とする。また、
\[
P>48L^2D^2.
\]
とする。\(b\) を \(b\equiv1\pmod6\) および \(|b-2P/D|\le3\) を満たす整数とし、\(a=b+2\)、\(t=(b-1)/6\)、\(\xi=P+1/2\) とおく。このとき、\(t\ge1\)、\(P>a\) であり、
\[
\lfloor\xi(a/b)^j\rfloor=P+jD
\]
が \(0\le j<L\) に対して成り立つ。
\label{thm-end:67}
\end{lemma}
\end{samepage}

% Markdown AST block 553: Proof.
\begin{proof}\leavevmode\label{proof:67}
\emph{\(6\mid D\)。} \(D\) が奇数ならば、\(P,P+D\) のいずれかは偶数になる。\(3\nmid D\) ならば、三つの項 \(P,P+D,P+2D\) は法 \(3\) のすべての剰余を覆う。

% Markdown AST block 554: Para
\emph{\(b\) の存在。} どの実数からも距離 \(3\) 以内に整数 \(\equiv1\pmod6\) が存在する。
\[
2P/D>96L^2D\ge96\cdot9\cdot6,
\]
より、
\[
b\ge2P/D-3\ge P/D\ge7,
\]
を得るので、\(t\ge1\) となる。

% Markdown AST block 555: Para
\emph{最初の評価。} \(b=2P/D+\varepsilon\) と書き、\(|\varepsilon|\le3\) とする。このとき、
\[
(2P+1)/b-D=(1-D\varepsilon)/b
\]
であり、\(|1-D\varepsilon|\le4D\) なので、
\[
|(2P+1)/b-D|\le4D^2/P.
\]
を得る。

% Markdown AST block 556: Para
\emph{二項展開。} \(u=2/b\) とおくと、
\[
(1+u)^j=1+ju+R_j
\]
であり、ここで
\[
0\le R_j\le\sum_{k\ge2}(ju)^k/k!\le\tfrac12(ju)^2e^{ju};
\]
である。ここでは
\[
ju\le2L/b\le2LD/P<1,
\]
なので、\(e^{ju}<3\) および
\[
R_j\le\tfrac32(ju)^2=6j^2/b^2\le6j^2D^2/P^2.
\]
が成り立つ。したがって、
\[
\begin{aligned}
&\Bigl|\bigl(P+\tfrac12\bigr)(1+u)^j-\bigl(P+jD+\tfrac12\bigr)\Bigr|\\
&\le j\Bigl|\frac{2P+1}{b}-D\Bigr|+\bigl(P+\tfrac12\bigr)R_j\\
&\le\frac{4jD^2}{P}+\frac{8j^2D^2}{P}\\
&\le\frac{12L^2D^2}{P}<\frac14,
\end{aligned}
\]
を得る。ここで \((P+1/2)/P\le4/3\) と、\(4j+8j^2\le12L^2\)（\(j<L\) に対して成り立つ）を用いた。よって
\[
\xi(a/b)^j\in(P+jD+1/4,P+jD+3/4)
\]
であり、その整数部分は \(P+jD\) である。

% Markdown AST block 557: Para
\emph{\(P>a\)。}
\[
a=b+2\le2P/D+5\le P/3+5<P
\]
が \(D\ge6\) より成り立つ。\qedhere
\end{proof}

% Markdown AST block 558: Theorem 68 (arbitrarily long prime prefixes).
\begin{samepage}
\begin{theorem}[任意に長い素数接頭列]\leavevmode\label{thm:68}
任意の \(L\ge3\) と任意の \(T\) に対し、ある \(t\ge T\) と有理数 \(\xi>0\) が存在して、\(\lfloor\xi((6t+3)/(6t+1))^j\rfloor\) は \(0\le j<L\) に対して素数であり、初項は \(6t+3\) を超える。したがって、C.3節のプログラムが停止するまでのステップ数には、\(t\) と初期素数の両方について一様な上界は存在しない。
\label{thm-end:68}
\end{theorem}
\end{samepage}

% Markdown AST block 559: Proof.
\begin{proof}\leavevmode\label{proof:68}
外部入力として、素数の狭い等差数列に関する定理を用いる。\(k\ge2\) および \(\delta>0\) に対して、任意の部分集合 \(A\)（\([1,Y]\) 内の素数の部分集合）が
\[
|A|\ge\delta Y/\log Y
\]
を満たすならば、すべての十分大きな \(Y\) に対して、非自明な \(k\) 項の等差数列で、その公差 \(D\) が
\[
|D|\le C_{k,\delta}(\log Y)^{(k-1)2^{k-2}}
\]
を満たすものが含まれる（\hyperref[ref-shao15]{SHAO15, Theorem 1.3}）。指数が \(k\) および \(\delta\) だけに依存する定性的な形は、\hyperref[ref-tz14]{TZ14, Theorem 5} を等差数列に特殊化したものである。（{[}TZ14{]} および {[}SHAO15{]} では文字 \(L\) は指数を表すが、ここでの \(L\) は長さ \(k\) である。）

% Markdown AST block 560: Para
\(k=L\) および
\[
A=\{p\text{ prime}:Y/2<p\le Y\};
\]
とおく。素数定理により、
\[
|A|\ge\tfrac13Y/\log Y
\]
が十分大きな \(Y\) に対して成り立つので、\(\delta=1/3\) は許容される。

% Markdown AST block 561: Para
得られる等差数列は \((Y/2,Y]\) 内にあるから、その最小項 \(P\) は \(P>Y/2\) を満たす。全項は \(3\) を十分大きな \(Y\) に対して超える。必要なら順序を逆にすれば、\(D>0\) であり、
\[
D\le C(\log Y)^{(L-1)2^{L-2}}.
\]
を満たす。

% Markdown AST block 562: Para
したがって、
\[
P/(48L^2D^2)\to\infty
\]
が \(Y\to\infty\) のとき成り立つ。十分大きな \(Y\) に対して補題67を適用でき、
\[
t=(b-1)/6\ge(P/D-1)/6\to\infty.\qedhere
\]
を得る。
\end{proof}

% Markdown AST block 563: Para
この主張の量化は
\[
\forall L\,\forall T\,\exists t\ge T\,\exists\xi;
\]
である。どの底についても無限の素数尾部の存在を主張せず、固定した \(t\) に対する停止の主張にも矛盾しない。外部定理の暗黙の定数と \(Y\) の閾値は明示されていないため、この方法では明示的な \(t\) は得られない。狭い等差数列に関する定理が不可欠である。\(D\) が \(P\) と同程度である等差数列では、補題67の仮定は満たされない。

% Markdown AST block 564: Para
\emph{例（九つの素数項）。}
\[
\begin{aligned}
&t=1654617,\qquad a=9927705,\qquad b=9927703,\\
&\xi=2084817839/2,\qquad P=1042408919,\qquad D=210,
\end{aligned}
\]
とおくと、
\[
\lfloor\xi(a/b)^j\rfloor=P+210j
\]
が \(0\le j\le8\) に対して成り立ち、これら九個の整数は素数である。一方、
\[
\lfloor\xi(a/b)^9\rfloor=1042410809=11\cdot94764619.
\]
である。仮定
\[
P>48\cdot81\cdot210^2=171460800
\]
は成り立つ。また、\(b\) は、整数 \(\equiv1\pmod6\) であって、\(2P/D\) に最も近いものである。整数部分は厳密な有理数演算で計算し、素数であることは試し割りで確認した。この例は何らかの主張への反例ではない。

% Markdown AST block 565: E.2 Return words as an obstruction
\Needspace{6\baselineskip}
\subsection{障害としての帰還語}\label{sec:E.2}

% Markdown AST block 566: Para
次の補題は定理42と同種のものである。最終的周期的でない繰り上がり列をもつ右向きの擬軌道を構成し、補題38を適用する。

% Markdown AST block 567: Lemma 69 (return-word obstruction).
\begin{samepage}
\begin{lemma}[帰還語の障害]\leavevmode\label{thm:69}
\(\gcd(m,ab)=1\)、\(q\in(\mathbb Z/m)^\times\) とし、\(U,V\) を次の条件を満たす繰り上がり語とする。\\
\textbf{(1)} 各語に沿って
\[
bz\equiv aq'+c\pmod m
\]
を \(q\) から辿ると、すべての剰余は単元であり、最後の剰余は \(q\) となる。\\
\textbf{(2)} 閉区間 \(I\subset(0,1)\) が存在し、逆向き写像
\[
\theta\mapsto(b\theta+c)/a
\]
を \(I\) の点から各語に沿って適用すると、すべての値は \((0,1)\) 内に保たれ、\(I\) は \(I\) 内に写る。\\
\textbf{(3)} \(UV\ne VU\)。\\
\textbf{(4)} すべての文字は
\[
\gcd(c,ab)=1\qquad\text{and}\qquad c\equiv b-a\pmod2.
\]
を満たす。
このとき、任意の \(M\) で、そのすべての素因数が \(2abm\) を割り切るものに対し、最終的周期的でない繰り上がり列をもつ右向きの \(M\) 擬軌道が存在する。また、\(\mathfrak M_0\)（または擬軌道に対して健全な操作による拡張）で、すべての \(M_t\mid M\) を満たす実行は成功しない。
\label{thm-end:69}
\end{lemma}
\end{samepage}

% Markdown AST block 568: Proof.
\begin{proof}\leavevmode\label{proof:69}
\emph{(i) \(m\) の素数冪。} \(\tilde m\) を、\(m\) と、\(M\) のうち \(m\) の素因数からなる部分との最小公倍数とする。語 \(w\) の長さを \(\ell\) とすると、剰余写像
\[
\pi_w:x\mapsto(a/b)^\ell x+C(w)/b^\ell
\]
を法 \(\tilde m\) で考えたもの
はアフィン置換である（\(p\nmid ab\)）。\(\pi_U(q)\equiv q\pmod m\) なので、\(\pi_U\) は有限なファイバー
\[
F=\{x\in\mathbb Z/\tilde m:x\equiv q\pmod m\},
\]
を置換する。\(\pi_V\) についても同様である。したがって、任意の \(\hat x\in F\) に対して、ある \(s,t\ge1\) が存在して
\[
\pi_U^s(\hat x)=\hat x\qquad \text{and}\qquad \pi_V^t(\hat x)=\hat x.
\]
を満たす。途中の剰余を法 \(m\) に還元すると、(1)の単元になるため、法 \(\tilde m\) でも単元である。

% Markdown AST block 569: Para
\emph{(ii) 非可換性。} \(U^sV^t=V^tU^s\) ならば、\(U^s\) と \(V^t\) は共通の原始語 \(z\) の冪である。\(w^k\) の原始根は \(w\) の原始根に等しいから、\(U\) と \(V\) はともに \(z\) の冪であり、可換となって(3)に矛盾する。

% Markdown AST block 570: Para
\emph{(iii) 小数部分。} \(U,V\) は \(I\) を \(I\) 内へ写すので、\(U^s,V^t\) も同様であり、すべての途中の値は \((0,1)\) に属する。同じ長さで相異なるブロック \(X_0=U^sV^t\) および \(X_1=V^tU^s\) を、最終的周期的でない二進列に従って連結する。定理42の証明の入れ子状の閉区間（各逆向きブロック写像は傾き \((b/a)^{|X_i|}\) の縮小写像である）により、\(\theta_n\) がすべての \(n\ge0\) に対して定まり、
\[
a\theta_n-b\theta_{n+1}=c_n
\]
が厳密に成り立つ。また、ブロックの復号により、繰り上がり列は最終的周期的でないことが分かる。

% Markdown AST block 571: Para
\emph{(iv) \(2ab\) を割り切る素数。} \(p\mid a\)、\(p\mid b\)、および \(p=2\)（\(a,b\) は奇数）について、(4)を用い、定理42の証明と同様に法 \(p^e\) の単元剰余を構成する。中国剰余定理ですべての剰余を組み合わせて法 \(\tilde m\cdot(\text{the
rest of }M)\) の擬軌道を得て、\(M\) に還元する。補題38により結論を得る。\qedhere
\end{proof}

% Markdown AST block 572: Corollary 70 (explicit witnesses).
\begin{samepage}
\begin{corollary}[明示的な証拠]\leavevmode\label{thm:70}
(a) \(9/7\) の場合、
\[
m=715=5\cdot11\cdot13,\qquad q=713,\qquad I=[1/5,3/10],
\]
および表14の語 \(U,V\) を用いると、\(\mathfrak M_0\) の実行で法に現れる素数が \(S_{13}=\{2,3,5,7,11,13\}\) だけであるものは成功しない。特に、法が \(4290\) の冪を割り切る実行は成功しない。
(b) \(6/5\) の場合、
\[
m=17017=7\cdot11\cdot13\cdot17,\qquad q=12854,\qquad I=[9/10,49/50],
\]
および表15の語 \(U,V\) を用いると、法に現れる素数が
\(S_{17}=\{2,3,5,7,11,13,17\}\) だけである実行は成功しない。
\label{thm-end:70}
\end{corollary}
\end{samepage}

% Markdown AST block 573: Proof.
\begin{proof}\leavevmode\label{proof:70}
補題69の仮定は厳密演算で検証される。表14--15の剰余列は、単元を経由して \(q\) に帰還する。\(I\) の端点の逆向き像は \((0,1)\) 内にとどまり、\(I\) に到達する（写像は単調増加なアフィン写像なので、端点の検査で十分である）。語の長さは異なり、どちらも他方の周期構造の接頭語ではないので、\(UV\ne VU\) である（\(9/7\) の場合、\(U\) と \(V\) は位置 \(2\) で異なる。\(6/5\) の場合は位置 \(19\) で異なる）。また、すべての文字はそれぞれ \(\mathcal C(9,7)=\{-4,-2,2,4,8\}\) および \(\mathcal C(6,5)=\{-1,1\}\) に属する。\qedhere
\end{proof}

% Markdown AST block 574: Table 14.
% Caption is rendered at the start of the immediately following table.

% Markdown AST block 575: Table 14
\begingroup
\small
\setlength{\tabcolsep}{4pt}
\renewcommand{\arraystretch}{1.15}
\setlength{\LTcapwidth}{\textwidth}
\begin{longtable}{@{}>{\raggedright\arraybackslash}p{\dimexpr0.16\dimexpr\linewidth-16pt\relax} >{\raggedright\arraybackslash}p{\dimexpr0.38\dimexpr\linewidth-16pt\relax} >{\raggedright\arraybackslash}p{\dimexpr0.46\dimexpr\linewidth-16pt\relax}@{}}
\caption{\(9/7\) の証拠（\(m=715\)、\(q=713\)）。剰余は \(q\) から帰還までを列挙する。}\label{tab:14}\\
\toprule
語 & 文字列 & 剰余列 \\
\midrule
\endfirsthead
\toprule
語 & 文字列 & 剰余列 \\
\midrule
\endhead
\bottomrule
\endlastfoot
\(U\) （長さ18） & \(2,\allowbreak{}-2,\allowbreak{}2,\allowbreak{}-2,\allowbreak{}4,\allowbreak{}-4,\allowbreak{}2,\allowbreak{}2,\allowbreak{}2,\allowbreak{}2,\allowbreak{}-4,\allowbreak{}4,\allowbreak{}-2,\allowbreak{}4,\allowbreak{}4,\allowbreak{}-4,\allowbreak{}2,\allowbreak{}4\) & 713,\allowbreak{} 202,\allowbreak{} 668,\allowbreak{} 42,\allowbreak{} 258,\allowbreak{} 128,\allowbreak{} 164,\allowbreak{} 109,\allowbreak{} 549,\allowbreak{} 604,\allowbreak{} 164,\allowbreak{} 6,\allowbreak{} 519,\allowbreak{} 667,\allowbreak{} 41,\allowbreak{} 564,\allowbreak{} 316,\allowbreak{} 713,\allowbreak{} 713 \\*
\(V\) （長さ17） & \(2,\allowbreak{}-2,\allowbreak{}-4,\allowbreak{}8,\allowbreak{}-4,\allowbreak{}4,\allowbreak{}-2,\allowbreak{}4,\allowbreak{}-4,\allowbreak{}4,\allowbreak{}2,\allowbreak{}2,\allowbreak{}2,\allowbreak{}-4,\allowbreak{}2,\allowbreak{}4,\allowbreak{}4\) & 713,\allowbreak{} 202,\allowbreak{} 668,\allowbreak{} 654,\allowbreak{} 127,\allowbreak{} 367,\allowbreak{} 166,\allowbreak{} 111,\allowbreak{} 654,\allowbreak{} 636,\allowbreak{} 614,\allowbreak{} 279,\allowbreak{} 359,\allowbreak{} 564,\allowbreak{} 316,\allowbreak{} 713,\allowbreak{} 713,\allowbreak{} 713 \\
\end{longtable}
\endgroup

% Markdown AST block 576: Table 15.
% Caption is rendered at the start of the immediately following table.

% Markdown AST block 577: Table 15
\begingroup
\small
\setlength{\tabcolsep}{4pt}
\renewcommand{\arraystretch}{1.15}
\setlength{\LTcapwidth}{\textwidth}
\begin{longtable}{@{}>{\raggedright\arraybackslash}p{\dimexpr0.16\dimexpr\linewidth-16pt\relax} >{\raggedright\arraybackslash}p{\dimexpr0.27\dimexpr\linewidth-16pt\relax} >{\raggedright\arraybackslash}p{\dimexpr0.57\dimexpr\linewidth-16pt\relax}@{}}
\caption{\(6/5\) の証拠（\(m=17017\)、\(q=12854\)）。\(+\) と \(-\) はそれぞれ繰り上がり \(+1\) と \(-1\) を表す。}\label{tab:15}\\
\toprule
語 & 文字列 & 剰余列 \\
\midrule
\endfirsthead
\toprule
語 & 文字列 & 剰余列 \\
\midrule
\endhead
\bottomrule
\endlastfoot
\(U\) （長さ33） & \(+^{15}-+^{3}-^{3}+^{9}-^{2}\) & 12854,\allowbreak{} 15425,\allowbreak{} 8300,\allowbreak{} 16767,\allowbreak{} 6507,\allowbreak{} 11212,\allowbreak{} 16858,\allowbreak{} 13423,\allowbreak{} 9301,\allowbreak{} 7758,\allowbreak{} 2503,\allowbreak{} 13214,\allowbreak{} 15857,\allowbreak{} 5415,\allowbreak{} 13305,\allowbreak{} 5756,\allowbreak{} 6907,\allowbreak{} 11692,\allowbreak{} 417,\allowbreak{} 3904,\allowbreak{} 8088,\allowbreak{} 6302,\allowbreak{} 14369,\allowbreak{} 226,\allowbreak{} 13885,\allowbreak{} 6452,\allowbreak{} 11146,\allowbreak{} 9972,\allowbreak{} 15370,\allowbreak{} 8234,\allowbreak{} 9881,\allowbreak{} 8454,\allowbreak{} 13548,\allowbreak{} 12854 \\*
\(V\) （長さ39） & \(+^{15}-+^{12}-^{2}+^{3}-+^{3}-^{2}\) & 12854,\allowbreak{} 15425,\allowbreak{} 8300,\allowbreak{} 16767,\allowbreak{} 6507,\allowbreak{} 11212,\allowbreak{} 16858,\allowbreak{} 13423,\allowbreak{} 9301,\allowbreak{} 7758,\allowbreak{} 2503,\allowbreak{} 13214,\allowbreak{} 15857,\allowbreak{} 5415,\allowbreak{} 13305,\allowbreak{} 5756,\allowbreak{} 6907,\allowbreak{} 11692,\allowbreak{} 417,\allowbreak{} 3904,\allowbreak{} 4685,\allowbreak{} 12429,\allowbreak{} 14915,\allowbreak{} 7688,\allowbreak{} 2419,\allowbreak{} 2903,\allowbreak{} 13694,\allowbreak{} 16433,\allowbreak{} 12913,\allowbreak{} 12092,\allowbreak{} 4300,\allowbreak{} 11967,\allowbreak{} 747,\allowbreak{} 4300,\allowbreak{} 15370,\allowbreak{} 8234,\allowbreak{} 9881,\allowbreak{} 8454,\allowbreak{} 13548,\allowbreak{} 12854 \\
\end{longtable}
\endgroup

% Markdown AST block 578: Para
\(6/5\) の場合、法 \(1001\) でのより短い証拠（根 \(380\)、\(I=[4/5,19/20]\)、長さ \(27\) の二つの語）も存在するが、(b)には法 \(17017\) の証拠で十分である。(i)の持ち上げにおける位数の例を挙げる。\(9/7\) については、法 \(715\cdot55\) で \((s,t)=(10,20)\)、法 \(715\cdot169\) で \((26,52)\) とできる。\(6/5\) については、法 \(17017\cdot77\) で \((14,14)\)、法 \(17017\cdot289\) で \((136,136)\) とできる。

% Markdown AST block 579: Para
8節と同様、これらは手続きが終了しないことに関する主張である。剰余 \(q_n\) は剰余類の列をなし、単一の軌道の整数部分ではない。補題41により、この証拠の任意の有限窓は、任意に大きい係数をもつ真の軌道により \(\forall N\,\exists\xi_N\) の形で実現される。しかし、証拠全体を実現する単一の \(\xi\) は得られない。また、実現する整数が素数であるとは主張せず、指定した素数と互いに素であることだけを主張している。定理58とは異なり、ここでは初期整数を固定しないため、コンパクト性の議論は適用できない。

% Markdown AST block 580: E.3 Eliminating one return-word language with a further prime
\Needspace{6\baselineskip}
\subsection{追加の素数による一つの帰還語言語の排除}\label{sec:E.3}

% Markdown AST block 581: Proposition 71.
\begin{samepage}
\begin{proposition}\leavevmode\label{thm:71}
\(a=9\)、\(b=7\) とし、\(U,V\) を表14の語とする。\(U\) および \(V\) をそれぞれ一本の辺とみなし、剰余 \(q\in\{1,\ldots,16\}\)（法 \(17\)）上にグラフを作る。辺 \(q\xrightarrow{w}z\) を置くのは、\(w\) を \(q\) から法 \(17\) で辿る際、最後を含むすべての剰余が0でない場合とする。このグラフの辺はちょうど次の十二本である。
\[
\begin{aligned}
&4\xrightarrow{U}16,\ 4\xrightarrow{V}16,\ 5\xrightarrow{U}1,\
\\
&6\xrightarrow{V}4,\ 8\xrightarrow{V}9,\ 9\xrightarrow{U}9,\
\\
&10\xrightarrow{V}14,\ 11\xrightarrow{U}13,\ 11\xrightarrow{V}8,\
\\
&14\xrightarrow{U}2,\ 15\xrightarrow{U}4,\ 16\xrightarrow{V}12,
\end{aligned}
\]
ランクを
\[
\operatorname{rank}(1,\ldots,16)=(0,0,0,2,1,3,0,1,0,2,2,0,0,1,3,1)
\]
とすると、自己ループ \(9\xrightarrow{U}9\) 以外のすべての辺でランクは狭義に減少する。したがって、比 \(9/7\) の真の軌道で、すべての項が素数であり、繰り上がり列が \(U\) と \(V\) の無限連結となるものは存在しない。
\label{thm-end:71}
\end{proposition}
\end{samepage}

% Markdown AST block 582: Proof.
\begin{proof}\leavevmode\label{proof:71}
辺一覧とランク不等式は有限の検証である。すべての項が素数である軌道では、\(17\nmid Q_n\) が、\(Q_n>17\) となった時点から成り立つ。したがって、ブロック境界における剰余 \(Q_n\bmod17\) はグラフ内の無限ウォークをなす。自己ループ以外のすべての辺でランクが狭義に減少するため、有限個のブロックの後は \(9\xrightarrow{U}9\) だけを用いる。繰り上がり列は最終的に \(U^\infty\) となり、周期 \(18\) で周期的となって補題4に矛盾する。\qedhere
\end{proof}

% Markdown AST block 583: Para
辺の表は \(a\bmod17\) および \(b\bmod17\) だけに依存するため、すべての組 \(a\equiv9\)、\(b\equiv7\pmod{17}\) に対して同じ十二本の辺が生じる。しかし、\(U\) および \(V\) が帰還語となるのは \(9/7\) の場合だけである（\(a=111\)、\(b=109\) の場合、語 \(U\) を \(713\) から法 \(715\) で辿っても帰還せず、非単元に出会う）。したがって、この族の他の元については何も主張しない。命題71は、\(U\) と \(V\) が生成する言語を排除するが、\(9/7\) のすべての繰り上がり列を扱うわけではない。また、素数 \(17\) を加えることで \(9/7\) の証明書が得られることも意味しない。

% Markdown AST block 584: E.4 Explicit two-carry certificates, prime selection, and the whole-word interval test
\Needspace{6\baselineskip}
\subsection{二つの繰り上がりによる明示的な証明書、素数の選択、全語区間判定}\label{sec:E.4}

% Markdown AST block 585: Para
\emph{定理42の具体例。} \(a/b=101/2\) について \(R=15\)（補助素数は \(3,5\)）とすると、
\[
\begin{aligned}
&\Delta=99,\qquad m=1,\qquad h=30,\\
&c=69,\qquad \gcd(69,202)=1;
\end{aligned}
\]
である。許容される最小の \(L\) は \(0\) であり、
\[
\begin{aligned}
&x_0=23/33,\qquad x_1=3379/3399,\\
&I=[23/33,3379/3399],\quad\text{and}\quad U=(69),\quad V=(99,69).
\end{aligned}
\]
となる。したがって、\(\mathfrak M_0\) の実行で、法に現れる素数が \(\{2,3,5,101\}\) だけであるものは成功しない。

% Markdown AST block 586: Para
\(a/b=101/100\) について \(R=21\)（素数は \(3,7\)）とすると、
\[
\begin{aligned}
&\Delta=1,\qquad m=1,\qquad h=42,\\
&c=-41,\qquad \gcd(41,10100)=1;
\end{aligned}
\]
である。閾値は \(-c/b=41/100\) であり、最小の \(L\) で \(x_L>41/100\) を満たすものを考える。ここで
\[
x_L=1-42\cdot100^L/(101^{L+1}-100^{L+1}),
\]
とすると、その値は \(L=54\) である。したがって、\(I=[x_{54},x_{55}]\) であり、対応する語は \(U=1^{54}(-41)\) および \(V=1^{55}(-41)\) である。この例は、負の繰り上がりも厳密に扱われることを示す。

% Markdown AST block 587: Para
\(5/2\) について \(R=1\) とすると、
\[
\begin{aligned}
&\Delta=3,\qquad m=1,\qquad h=2,\qquad c=1,\qquad L=0,\\
&U=(1),\qquad V=(3,1),
\end{aligned}
\]
である。これは、この底に対する系43の証拠である。

% Markdown AST block 588: Lemma 72 (primes producing a common nonzero translation).
\begin{samepage}
\begin{lemma}[共通の非零平行移動を生じさせる素数]\leavevmode\label{thm:72}
\(W=\{w_1,\ldots,w_s\}\) を各要素が空でない語である集合とし、それぞれの長さを \(\ell_i\)、(8)で定めた量を \(C_i=C(w_i)\) とする。\(g=\gcd(\ell_1,\ldots,\ell_s)\) とし、
\[
D=\gcd\bigl(a^g-b^g,\ C_ib^{\ell_1}-C_1b^{\ell_i}\ (2\le i\le s)\bigr).
\]
とおく。素数 \(p\nmid ab\) に対して、剰余写像
\[
F_i(q)=(a^{\ell_i}q+C_i)/b^{\ell_i}
\]
がすべて法 \(p\) で同一の平行移動 \(q\mapsto q+\tau\)（\(\tau\ne0\)）に合同であるための必要十分条件は、\(p\mid D\) および \(p\nmid C_1\) である。この場合、\(W\) の語の無限連結は、いずれも \(p\) による持ち上げを生き残れない。ブロック境界での剰余は \(q,q+\tau,q+2\tau,\ldots\) であり、\(p\) ブロック以内に0になる。
\label{thm-end:72}
\end{lemma}
\end{samepage}

% Markdown AST block 589: Proof.
\begin{proof}\leavevmode\label{proof:72}
\(F_i\) の傾きは \((a/b)^{\ell_i}\) である。すべての傾きが \(1\) であることと、\(a/b\) の法 \(p\) での位数がすべての \(\ell_i\) を割り切ることは同値である。これは、その位数が \(g\) を割り切ること、すなわち \(p\mid a^g-b^g\) と同値である。

% Markdown AST block 590: Para
平行移動 \(C_i/b^{\ell_i}\) が一致することと
\[
C_ib^{\ell_1}\equiv C_1b^{\ell_i}\pmod p,
\]
は同値である。また、共通の平行移動が0でないことと \(p\nmid C_1\) は同値である（\(b\) は単元）。\qedhere
\end{proof}

% Markdown AST block 591: Para
これは \(W\) の連結だけを排除する。成分を除去できるのは、その成分のすべての帰還語が \(W\) に属する場合だけである。定理42の二つの繰り上がりからなる語 \(U=\Delta^Lc\)、\(V=\Delta^{L+1}c\) については、\(g=1\) および \(D\mid a-b\) である。\(p\mid\Delta\) および \(p\nmid c\) ならば、両方の語は \(q\mapsto q+c/b\) として作用する。\(101/2\) の場合、\(D=99\)、\(C_1=69\) である。素数 \(3\) は \(C_1\) を割り切るため使えないが、\(11\) は平行移動
\[
\tau=69\cdot2^{-1}\equiv7\pmod{11},
\]
を与える。したがって、\(3\) および \(5\) を生き残る二語による分岐は、\(11\) を加えることで排除される。\(5/2\) の場合には \(D=3\) であり、\(3\) が同じ役割を果たす。\(\Delta\in
\{1,2\}\) の場合、\(D\mid2\) であり、この方法では奇素数は得られない。これは、この仕組み自体の限界である。

% Markdown AST block 592: Proposition 73 (whole-word interval test).
\begin{samepage}
\begin{proposition}[全語区間判定]\leavevmode\label{thm:73}
始点セル \([j/K,(j+1)/K)\) と語 \(c_0\cdots c_{\ell-1}\) をもつマクロ辺に対して、
\[
\begin{aligned}
(L_0,U_0,B_0)&=(j,j+1,K)\qquad\text{and}\\
B_{i+1}&=bB_i,\\
L_{i+1}&=\max(0,aL_i-c_iB_i),\\
U_{i+1}&=\min(B_{i+1},aU_i-c_iB_i).
\end{aligned}
\]
と定義する。このとき、\([L_i/B_i,U_i/B_i)\) は、このセルから始まり、これらの繰り上がりをもつ実数軌道片の、\(i\) 文字後の小数部分の集合に正確に一致する。したがって、\(L_i\ge U_i\) がある \(i\) に対して成り立つ場合（または最終区間が終点セルと交わらない場合）にマクロ辺を棄却する操作は、擬軌道に対して健全である。この辺を通る任意の擬軌道、特に任意の真の軌道には、非空性の証拠となる厳密な小数部分が存在する。この操作はウォークに対しては健全ではない。(E3)が表すのは隣接する二つのセルの両立性だけであり、長さ \(\ge2\) の道で \(G(M,K)\) 内にあるものが実数の小数部分をもつとは限らないためである。
\label{thm-end:73}
\end{proposition}
\end{samepage}

% Markdown AST block 593: Proof.
\begin{proof}\leavevmode\label{proof:73}
写像
\[
\theta\mapsto(a\theta-c)/b
\]
は単調増加なアフィン全単射なので、\([L/B,U/B)\) の像は
\[
[(aL-cB)/(bB),(aU-cB)/(bB)),
\]
であり、\([0,1)\) との共通部分を取ると更新式を得る。厳密性は命題57と同様、帰納法から従う。\qedhere
\end{proof}

% Markdown AST block 594: Para
この判定を追加すると定理37のクラスが変わるかどうかは、本論文では決定していない。この判定を用いた探索計算は採用しておらず、5節の証明書でも用いていない。定理42と補題69の証拠は、すべての途中の時刻でこの判定を通過する。この点で、見かけ上の局所的な道とは異なる。

% BEGIN PUBLIC ADDITION acknowledgments
\section*{謝辞およびAIの関与に関する開示}

本研究の発端は、著者がミルズの定数に関心を持ったことである。ここでミルズの定数とは、すべての正整数 \(n\) に対して \(\lfloor A^{3^n}\rfloor\) が素数となる実数 \(A>1\) の最小値を指す。この最小の定数の無理数性が証明されたこと \cite{SA25} を知り、著者は ChatGPT 上の \href{https://help.openai.com/en/articles/20001354-gpt-5-6}{GPT-6 Pro}（GPT-6 Astra を基盤とするモデル） に、より強い「すべてのミルズ数は無理数である」という主張を証明するよう求めた。その後、著者自身の記憶では約500回、「頑張れ」「諦めるな」といった継続を促す指示を繰り返した。しかし、この探索では求めた証明には到達しなかった。そこで、それまでに得られた内容から論文化に値する数学的な成果を見いだして論文にまとめるようAIに依頼した。本稿はその依頼から生まれたものであり、すべてのミルズ数に関する当初の問いを解決するものではない。

著者は、本研究で主に用いたAIである GPT-6 Astra（ChatGPT では GPT-6 Pro として利用し、Codex でも使用）に謝意を表する。また、論文化の途中で Codex の週次利用上限（Weekly Limit）に達した際、一部の作業を Claude Code で補助した Claude Fable 5.1 にも、その限定的な寄与について謝意を表する。最初の予想を提示した後の数学的な展開、論文化に値する成果の見いだし、証明の試み、計算プログラム、文章の執筆、および Lean による形式化は、すべてAIが行った。これらの作業に対し、著者は数学的に意義のある指示や助言を一切していない。形式化や、Codex 内で執筆担当と批評担当のエージェントが改稿を繰り返してから組版する自動的な作業手順を求めたのも、数学的内容を提供せずに、作業と手順を指定したものである。これらはAIによる自動的な批評であり、人間の数学者による独立した査読ではない。

この過程で著者が行った数学的なことは、「最小のミルズ数だけでなく、すべてのミルズ数が無理数なのではないか」という最初の予想を持ったことだけである。それ以外に、著者は数学的な貢献を何もしていない。証明の着想、数学的議論、計算、結果の数学的評価を提供したこともない。その後の関与は、作業を続けること、論文化できる内容を見いだすこと、論文を書くこと、Lean による証明を作ることを求める非数学的な指示と、公開するという判断に限られた。数学的証明を自ら検証したわけではない。結果はおそらく正しく数学的にも意義があるだろうという著者の見方は、検証を行ったという報告ではない。読者には、議論、再現可能な計算、および形式証明の成果物それ自体を評価していただきたい。Lean による証明の範囲と信頼上の仮定は第9.4節に明示している。

著者は、人間が最初の問いを示し、その後の数学的・技術的な作業をすべてAIが行う研究が、数学を含む多くの分野で今後ますます一般的になると考えている。研究を始めた人間自身が内容に踏み込んだ助言をせず、成果を十分理解していない場合も含まれる。これは著者個人の将来予想であり、本稿の数学的結果から導かれる結論ではない。この経緯を記すのは、成果がどのように生まれたか、AIが何を担ったか、そして著者自身の貢献がどこまでに限られるかを明らかにするためである。
% END PUBLIC ADDITION acknowledgments

% Markdown AST block 595: References
\begin{thebibliography}{SHAO15}

% Markdown AST block 596: reference AFS08
\bibitem[AFS08]{AFS08}\phantomsection\label{ref-afs08}
Shigeki Akiyama, Christiane Frougny, and Jacques Sakarovitch,
\emph{Powers of rationals modulo 1 and rational base number systems},
Israel Journal of Mathematics \textbf{168} (2008), 53--91.
\href{https://doi.org/10.1007/s11856-008-1056-4}{Published article}.

% Markdown AST block 597: reference AD04
\bibitem[AD04]{AD04}\phantomsection\label{ref-ad04}
Giedrius Alkauskas and Artūras Dubickas, \emph{Prime and composite
numbers as integer parts of powers}, Acta Mathematica Hungarica
\textbf{105} (2004), no. 3, 249--256.
\href{https://doi.org/10.1023/B:AMHU.0000049291.38284.2c}{Published article}.

% Markdown AST block 598: reference AEV26
\bibitem[AEV26]{AEV26}\phantomsection\label{ref-aev26}
Mélodie Andrieu, Shalom Eliahou, and Léo Vivion, \emph{A
Normality Conjecture on Rational Base Number Systems}, arXiv:2510.11723
{[}math.NT{]}, version 2, April 7, 2026 (version 1: October 6, 2025).
\href{https://arxiv.org/abs/2510.11723v2}{Preprint}.

% Markdown AST block 599: reference CW13
\bibitem[CW13]{CW13}\phantomsection\label{ref-cw13}
Fintan Costello and Paul Watts, \emph{A short note on
Jacobsthal's function}, arXiv:1306.1064 {[}math.NT{]}, version 1, June 5,
2013. \href{https://arxiv.org/abs/1306.1064v1}{Preprint}.

% Markdown AST block 600: reference Dub06
\bibitem[Dub06]{Dub06}\phantomsection\label{ref-dub06}
Artūras Dubickas, \emph{Truncatable primes and unavoidable sets of
divisors}, Acta Mathematica Universitatis Ostraviensis \textbf{14}
(2006), 21--25.
\href{https://klevas.mif.vu.lt/~dubickas/files/dvifai/cekko1.pdf}{Author's copy}.

% Markdown AST block 600: reference Dub08
\bibitem[Dub08]{Dub08}\phantomsection\label{ref-dub08}
Artūras Dubickas, \emph{On the powers of 3/2 and other rational
numbers}, Mathematische Nachrichten \textbf{281} (2008), no. 7, 951--958.
\href{https://doi.org/10.1002/mana.200510651}{Published article}.

% Markdown AST block 600: reference Dub09
\bibitem[Dub09]{Dub09}\phantomsection\label{ref-dub09}
Artūras Dubickas, \emph{Prime and composite integers close to powers of
a number}, Monatshefte für Mathematik \textbf{158} (2009), no. 3,
271--284. \href{https://doi.org/10.1007/s00605-008-0042-6}{Published article}.

% Markdown AST block 601: reference DJ22
\bibitem[DJ22]{DJ22}\phantomsection\label{ref-dj22}
Artūras Dubickas and Lukas Jonuška, \emph{Divisibility of integers
obtained from truncated periodic sequences}, International Journal of
Number Theory \textbf{18} (2022), no. 1, 165--174.
\href{https://doi.org/10.1142/S1793042122500129}{Published article}.

% Markdown AST block 602: reference DN
\bibitem[DN]{DN}\phantomsection\label{ref-dn}
Artūras Dubickas and Aivaras Novikas, \emph{Integer parts of powers of
rational numbers}, Mathematische Zeitschrift \textbf{251} (2005), no. 3,
635--648. \href{https://doi.org/10.1007/s00209-005-0827-4}{Published article}.

% Markdown AST block 603: reference FLP
\bibitem[FLP]{FLP}\phantomsection\label{ref-flp}
Leopold Flatto, Jeffrey C. Lagarias, and Andrew D. Pollington,
\emph{On the range of fractional parts \(\{\xi(p/q)^n\}\)},
Acta Arithmetica \textbf{70} (1995), no. 2, 125--147.
\href{https://doi.org/10.4064/aa-70-2-125-147}{Published article}.

% Markdown AST block 604: reference FS
\bibitem[FS]{FS}\phantomsection\label{ref-fs}
W. Forman and H. N. Shapiro, \emph{An arithmetic property of certain
rational powers}, Communications on Pure and Applied Mathematics
\textbf{20} (1967), no. 3, 561--573.
\href{https://doi.org/10.1002/cpa.3160200305}{Published article}.

% Markdown AST block 605: reference IW78
\bibitem[IW78]{IW78}\phantomsection\label{ref-iw78}
Henryk Iwaniec, \emph{On the problem of Jacobsthal}, Demonstratio
Mathematica \textbf{11} (1978), no. 1, 225--231.
\href{https://doi.org/10.1515/dema-1978-0121}{Published article}.

% Markdown AST block 606: reference M68
\bibitem[M68]{M68}\phantomsection\label{ref-m68}
K. Mahler, \emph{An unsolved problem on the powers of 3/2},
Journal of the Australian Mathematical Society \textbf{8} (1968), no. 2,
313--321. \href{https://doi.org/10.1017/S1446788700005371}{Published article}.

% Markdown AST block 607: reference N
\bibitem[N]{N}\phantomsection\label{ref-n}
Aivaras Novikas, \emph{Composite numbers in the sequences of integers},
doctoral dissertation, Vilnius University, 2012.
\href{https://epublications.vu.lt/object/elaba:2089125/index.html}{University repository}.

% Markdown AST block 608: reference SHAO15
\bibitem[SHAO15]{SHAO15}\phantomsection\label{ref-shao15}
Xuancheng Shao, \emph{Narrow arithmetic progressions in the
primes}, arXiv:1509.04955 {[}math.NT{]}, version 1, September 16, 2015;
published in International Mathematics Research Notices \textbf{2017}, no. 2,
391--428. \href{https://arxiv.org/abs/1509.04955v1}{Preprint},
\href{https://doi.org/10.1093/imrn/rnv393}{Published article}.

% Markdown AST block 609: reference S26
\bibitem[S26]{S26}\phantomsection\label{ref-s26}
Ralf Stephan, \emph{On the composites among \([\xi 7^n]\)}, author-uploaded
preprint dated August 15, 2026.
\href{https://www.researchgate.net/publication/412294027_On_the_composites_among_x7}{Preprint};
\href{https://github.com/rwst/On-Composites}{Lean files}.

% Markdown AST block 610: reference TZ14
\bibitem[TZ14]{TZ14}\phantomsection\label{ref-tz14}
Terence Tao and Tamar Ziegler, \emph{Narrow progressions in the
primes}, arXiv:1409.1327 {[}math.NT{]}, version 2, October 10, 2014;
published in \emph{Analytic Number Theory} (Springer, 2015), 357--379.
\href{https://arxiv.org/abs/1409.1327v2}{Preprint},
\href{https://doi.org/10.1007/978-3-319-22240-0_22}{Published chapter}.

% Markdown AST block 611: reference T
\bibitem[T]{T}\phantomsection\label{ref-t}
Robert Tarjan, \emph{Depth-First Search and Linear Graph Algorithms},
SIAM Journal on Computing \textbf{1} (1972), no. 2, 146--160.
\href{https://doi.org/10.1137/0201010}{Published article}.

% BEGIN PUBLIC ADDITION reference-sa25
\bibitem[SA25]{SA25}\phantomsection\label{ref-sa25}
Kota Saito, \emph{Mills' constant is irrational},
Mathematika \textbf{71} (2025), no. 3, e70027.
\href{https://doi.org/10.1112/mtk.70027}{Published article},
\href{https://arxiv.org/abs/2404.19461}{Preprint}.
% END PUBLIC ADDITION reference-sa25

\end{thebibliography}
\end{document}
